\begin{proof}[Proof of \cref{obj:thm:unrestricted-near-complete-arrival-envelope}]
\begin{enumerate}
\item Choose the universal constant \(\overline C>0\) supplied by
\cref{obj:thm:unrestricted-mixed-count-upper}. For every admissible law \(P\) and the specified sampling law, that result gives
\[
E_P\!\left[
\left(\widehat\tau^{\mathrm{mix}}_{n,d,q}(\mathbf O_n)-\tau(P)\right)^2
\right]
\le \mathfrak U_{n,d,q}
\le \overline C\,r_{n,d,q},
\]
and
\[
E_P\!\left[
\left(\widehat\tau^{\mathrm{mix}}_{n,d,q}(\mathbf O_n)-\tau(P)\right)^2
\right]
\le
E_P\!\left[
E_\omega\!\left[
\left(\widetilde\tau(\mathbf O_n;\omega)-\tau(P)\right)^2
\,\middle|\,\mathbf O_n
\right]\right].
\]
Its envelope is exactly the branchwise envelope in the statement, and its auxiliary output includes the zero output on the Poisson tail. We retain this same constant throughout.
% lean: CausalSmith.Stat.MarRareqLogfrontier.unrestricted_mixed_count_upper

\item We first record how a deterministic risk bound yields a minimax bound. Write
\[
\mathcal O_d:=\{0,\ldots,d-1\}\times\{0,1\}^{4}
\]
for the observed-record space. For a bounded estimator kernel \(\mathsf T\), define its risk by
\[
\mathcal R_P(\mathsf T)
:=
\int_{\mathcal O_d^n}\int_{[-1,1]}
(h-\tau(P))^2\,\mathsf T(\mathbf o,dh)\,
\mathcal L_P(\mathbf O_n)(d\mathbf o).
\]
These risks are nonnegative.

Let \(f_{\mathrm{det}}:\mathcal O_d^n\to[-1,1]\) be measurable, and let \(b_{\mathrm{up}}\ge0\). Suppose
\[
E_P\!\left[
\left(f_{\mathrm{det}}(\mathbf O_n)-\tau(P)\right)^2
\right]\le b_{\mathrm{up}}
\qquad
\text{for every }P\in\mathcal M^{+}_{n,d,q}.
\]
The parameter-independent kernel
\[
\mathsf T_{\mathrm{det}}(\mathbf o,\cdot)
:=\operatorname{Dirac}\!\left(f_{\mathrm{det}}(\mathbf o)\right)
\]
belongs to the decision class in \cref{obj:def:unrestricted-minimax-risk}, and satisfies
\[
\mathcal R_P(\mathsf T_{\mathrm{det}})
=
E_P\!\left[
\left(f_{\mathrm{det}}(\mathbf O_n)-\tau(P)\right)^2
\right].
\]
Consequently,
\[
\mathfrak R^{+}_{n,d,q}
\le
\sup_{P\in\mathcal M^{+}_{n,d,q}}
E_P\!\left[
\left(f_{\mathrm{det}}(\mathbf O_n)-\tau(P)\right)^2
\right]
\le b_{\mathrm{up}}.
\]
For a nonempty model class, the last inequality follows by taking the supremum of the pointwise bounds. For an empty class, the real supremum is defined to be zero, and the inequality follows from \(b_{\mathrm{up}}\ge0\).
% lean: CausalSmith.Stat.MarRareqLogfrontier.unrestricted_deterministic_upper

\item Fix \(n,d\ge1\) and \(0<q\le1\). Binary potential outcomes and
\cref{obj:def:ate-functional} give
\[
\tau(P)=P(Y(1)=1)-P(Y(0)=1)\in[-1,1].
\]
Indeed, each probability belongs to \([0,1]\). The zero map is measurable and takes values in \([-1,1]\), and its risk under the canonical product probability law is
\[
E_P\!\left[(0-\tau(P))^2\right]=\tau(P)^2\le1.
\]
Applying the deterministic comparison just established yields
\[
\mathfrak R^{+}_{n,d,q}\le1.
\]
% lean: CausalSmith.Stat.MarRareqLogfrontier.ate_mem_unit_interval

\item The rate in \cref{obj:def:frontier-rate} satisfies
\[
r_{n,d,q}
=
\min\left\{
1,\,
\frac1{nq}
+
\left[\frac{d}{nq\log(e+nq)}\right]^2
\right\}\ge0,
\qquad
\overline C\,r_{n,d,q}\ge0.
\]
The estimator in \cref{obj:def:mixed-count-estimator} is measurable. Its auxiliary output lies in \([-1,1]\) for every sample and every randomization, so averaging gives
\[
-1
\le
E_\omega\!\left[
\widetilde\tau(\mathbf o;\omega)\mid\mathbf o
\right]
=
\widehat\tau^{\mathrm{mix}}_{n,d,q}(\mathbf o)
\le1.
\]
Under the canonical product law, the coordinate observations are measurable and independent, with
\[
\mathcal L_P(O_i)=\mathcal L_P(O_1),
\qquad i=1,\ldots,n,
\]
so they satisfy \cref{obj:ass:iid-sampling}. The first step therefore applies to every law in the model class. The deterministic comparison, with \(b_{\mathrm{up}}=\overline C\,r_{n,d,q}\), gives
\[
\mathfrak R^{+}_{n,d,q}\le\overline C\,r_{n,d,q}.
\]
% lean: CausalSmith.Stat.MarRareqLogfrontier.unrestricted_near_complete_arrival_envelope

\item We next establish the observed-contrast bound. Fix
\(P\in\mathcal M^{+}_{n,d,q}\). For an observed record, write its coordinates as
\[
o=(X(o),A(o),S(o),R(o),(RY)(o)),
\]
and define
\[
\xi_{\mathrm{obs}}(o)
:=
2(2A(o)-1)\mathbf1\{R(o)=1,\ (RY)(o)=1\}.
\]
For a sample \(\mathbf o=(o_1,\ldots,o_n)\in\mathcal O_d^n\), define
\[
\overline\xi_n(\mathbf o)
:=
\frac1n\sum_{i=1}^n\xi_{\mathrm{obs}}(o_i).
\]
The canonical product law gives
\[
E_P[\overline\xi_n(\mathbf O_n)]
=E_P[\xi_{\mathrm{obs}}(O_1)],
\qquad
\operatorname{Var}_P(\overline\xi_n(\mathbf O_n))
=\frac1n\operatorname{Var}_P(\xi_{\mathrm{obs}}(O_1)).
\]
Since every observed coordinate is binary,
\[
-2\le\xi_{\mathrm{obs}}(o)\le2,
\qquad
\operatorname{Var}_P(\xi_{\mathrm{obs}}(O_1))
\le E_P[\xi_{\mathrm{obs}}(O_1)^2]\le4.
\]
Expanding the squared error around the mean now gives
\[
E_P\!\left[
\left(\overline\xi_n(\mathbf O_n)-\tau(P)\right)^2
\right]
=
\frac{\operatorname{Var}_P(\xi_{\mathrm{obs}}(O_1))}{n}
+
\left(E_P[\xi_{\mathrm{obs}}(O_1)]-\tau(P)\right)^2.
\]
By \cref{obj:def:complete-arrival-ht-estimator},
\[
\widehat\tau^{\mathrm{HT}}_n(\mathbf o)
=
\operatorname{clip}_{[-1,1]}(\overline\xi_n(\mathbf o)).
\]
Because \(\tau(P)\in[-1,1]\), considering whether
\(\overline\xi_n(\mathbf o)<-1\), belongs to \([-1,1]\), or exceeds \(1\), gives
\[
\left(
\widehat\tau^{\mathrm{HT}}_n(\mathbf o)-\tau(P)
\right)^2
\le
\left(\overline\xi_n(\mathbf o)-\tau(P)\right)^2.
\]
% lean: CausalSmith.Stat.MarRareqLogfrontier.observed_contrast_risk_le

\item To bound the bias, define the missing-outcome masses for both arms by
\[
\Delta_a(P)
:=
2P(A=a,R=0,Y=1),
\qquad a\in\{0,1\}.
\]
For every cell indexed by \(a,x,s\), the occupied-cell arrival condition implies
\[
q\,p_{axs}(P)
\le P(A=a,X=x,S=s,R=1).
\]
For an occupied cell this is the arrival-floor inequality in
\cref{obj:ass:occupied-cell-arrival}; for a null cell its left side is zero and its right side is nonnegative. Partitioning a cell according to \(R\) therefore gives
\[
\begin{aligned}
P(A=a,X=x,S=s,R=0)
&=
p_{axs}(P)-P(A=a,X=x,S=s,R=1)\\
&\le(1-q)p_{axs}(P).
\end{aligned}
\]
Balanced assignment gives \(P(A=1)=1/2\) and
\(P(A=0)=1-P(A=1)=1/2\). Thus, for each \(a\in\{0,1\}\),
\[
\sum_{x=0}^{d-1}\sum_{s=0}^{1}p_{axs}(P)=\frac12,
\]
and the finite cell partition yields
\[
\begin{aligned}
P(A=a,R=0)
&=
\sum_{x=0}^{d-1}\sum_{s=0}^{1}
P(A=a,X=x,S=s,R=0)\\
&\le
(1-q)\sum_{x=0}^{d-1}\sum_{s=0}^{1}p_{axs}(P)
=\frac{1-q}{2}.
\end{aligned}
\]
The event \(\{A=a,R=0,Y=1\}\) is contained in \(\{A=a,R=0\}\), so
\[
0\le\Delta_a(P)\le2P(A=a,R=0)\le1-q.
\]

Outcome consistency, randomized independence, and balanced assignment give
\[
\begin{aligned}
P(A=a,Y=1)
&=P(A=a,Y(a)=1)\\
&=P(A=a)P(Y(a)=1)
=\frac12P(Y(a)=1).
\end{aligned}
\]
Using the observed-record map and partitioning the outcome-one event according to arrival, we obtain
\[
\begin{aligned}
E_P[\xi_{\mathrm{obs}}(O_1)]
&=
2P(A=1,R=1,Y=1)-2P(A=0,R=1,Y=1)\\
&=
2P(A=1,Y=1)-2P(A=0,Y=1)
-\Delta_1(P)+\Delta_0(P)\\
&=\tau(P)-\Delta_1(P)+\Delta_0(P).
\end{aligned}
\]
The bounds on the two deficits imply
\[
-(1-q)
\le
E_P[\xi_{\mathrm{obs}}(O_1)]-\tau(P)
\le1-q,
\]
and hence, since \(1-q\ge0\),
\[
\left(E_P[\xi_{\mathrm{obs}}(O_1)]-\tau(P)\right)^2
\le(1-q)^2.
\]
Together with the preceding variance and clipping bounds, this gives
\[
\begin{aligned}
E_P\!\left[
\left(\widehat\tau^{\mathrm{HT}}_n(\mathbf O_n)-\tau(P)\right)^2
\right]
&\le
E_P\!\left[
\left(\overline\xi_n(\mathbf O_n)-\tau(P)\right)^2
\right]\\
&\le\frac4n+(1-q)^2.
\end{aligned}
\]
The clipped contrast is measurable on the finite sample space and takes values in \([-1,1]\). Since \(4/n+(1-q)^2\ge0\), the deterministic comparison gives both
\[
\mathfrak R^{+}_{n,d,q}\le\frac4n+(1-q)^2
\]
and
\[
\sup_{P\in\mathcal M^{+}_{n,d,q}}
E_P\!\left[
\left(\widehat\tau^{\mathrm{HT}}_n(\mathbf O_n)-\tau(P)\right)^2
\right]
\le\frac4n+(1-q)^2.
\]
% lean: CausalSmith.Stat.MarRareqLogfrontier.observed_contrast_risk_le

\item For the lower bound, first construct a complete-arrival pair. Define
\[
\delta_{\mathrm{gap}}:=\frac1{4\sqrt n},
\qquad
0<\delta_{\mathrm{gap}}\le\frac14,
\qquad
\delta_{\mathrm{gap}}^2=\frac1{16n}.
\]
For \(\iota\in\{0,1\}\), let \(P_{\mathrm{test},\iota}\) be the full-data law specified by
\[
\begin{gathered}
X=0,\qquad S(0)=S(1)=0,\qquad Y(0)=0,\qquad R=1,\\
A\sim\operatorname{Bernoulli}(1/2),\qquad
Y(1)\sim\operatorname{Bernoulli}(1/2+\iota\delta_{\mathrm{gap}}),
\qquad A\perp\!\!\!\perp Y(1),\\
S=S(A)=0,\qquad Y=Y(A)=A\,Y(1).
\end{gathered}
\]
The baseline value \(0\) is available for every \(d\ge1\), and the outcome parameters lie in \([1/2,3/4]\). Independence and the fair treatment coin verify randomized independence and balanced assignment. The displayed definitions verify consistency. Constant arrival verifies arrival MAR and gives arrival probability one in every occupied cell. Therefore
\[
P_{\mathrm{test},0},P_{\mathrm{test},1}
\in\mathcal M^{+}_{n,d,1},
\qquad
\tau(P_{\mathrm{test},\iota})
=\frac12+\iota\delta_{\mathrm{gap}}.
\]

Define the observed laws and their sample laws by
\[
Q_{\mathrm{test},\iota}
:=\mathcal L_{P_{\mathrm{test},\iota}}(O_1),
\qquad
Q_{\mathrm{test},\iota}^{(n)}
:=Q_{\mathrm{test},\iota}^{\otimes n}.
\]
Their common support consists of the three records
\[
\begin{aligned}
o_{\mathrm{test},\mathrm{control}}&:=(0,0,0,1,0),\\
o_{\mathrm{test},\mathrm{zero}}&:=(0,1,0,1,0),\\
o_{\mathrm{test},\mathrm{one}}&:=(0,1,0,1,1),
\end{aligned}
\qquad
\mathcal Z_{\mathrm{test}}
:=
\{o_{\mathrm{test},\mathrm{control}},
o_{\mathrm{test},\mathrm{zero}},
o_{\mathrm{test},\mathrm{one}}\}.
\]
The atom probabilities are
\[
\begin{aligned}
Q_{\mathrm{test},\iota}(\{o_{\mathrm{test},\mathrm{control}}\})&=\frac12,\\
Q_{\mathrm{test},\iota}(\{o_{\mathrm{test},\mathrm{zero}}\})&=\frac14-\frac{\iota\delta_{\mathrm{gap}}}{2},\\
Q_{\mathrm{test},\iota}(\{o_{\mathrm{test},\mathrm{one}}\})&=\frac14+\frac{\iota\delta_{\mathrm{gap}}}{2},
\end{aligned}
\]
with probability zero elsewhere. In particular,
\(Q_{\mathrm{test},1}\ll Q_{\mathrm{test},0}\).

For finite laws with a positive reference law on their common support, define
\[
\chi^2(Q_{\mathrm{test},1}\Vert Q_{\mathrm{test},0})
:=
\sum_{o\in\mathcal Z_{\mathrm{test}}}
\frac{
\left(Q_{\mathrm{test},1}(\{o\})-Q_{\mathrm{test},0}(\{o\})\right)^2
}{
Q_{\mathrm{test},0}(\{o\})
}.
\]
Summing the squared probability ratios gives
\[
\begin{aligned}
1+\chi^2(Q_{\mathrm{test},1}\Vert Q_{\mathrm{test},0})
&=
\sum_{o\in\mathcal Z_{\mathrm{test}}}
\frac{Q_{\mathrm{test},1}(\{o\})^2}{Q_{\mathrm{test},0}(\{o\})}\\
&=
\frac{(1/2)^2}{1/2}
+
\frac{(1/4-\delta_{\mathrm{gap}}/2)^2}{1/4}
+
\frac{(1/4+\delta_{\mathrm{gap}}/2)^2}{1/4}\\
&=1+2\delta_{\mathrm{gap}}^2.
\end{aligned}
\]
Thus
\[
\chi^2(Q_{\mathrm{test},1}\Vert Q_{\mathrm{test},0})
=2\delta_{\mathrm{gap}}^2=\frac1{8n}.
\]
% lean: CausalSmith.Stat.MarRareqLogfrontier.completeFamilyLaw_obs_chi

\item Define the sample likelihood ratio on \(\mathcal O_d^n\) by
\[
L_{\mathrm{test},n}(\mathbf o)
:=
\begin{cases}
\displaystyle\prod_{i=1}^{n}
\frac{Q_{\mathrm{test},1}(\{o_i\})}
{Q_{\mathrm{test},0}(\{o_i\})},
&\mathbf o\in\mathcal Z_{\mathrm{test}}^n,\\
0,&\mathbf o\notin\mathcal Z_{\mathrm{test}}^n.
\end{cases}
\]
The product laws satisfy
\[
E_{Q_{\mathrm{test},0}^{(n)}}[L_{\mathrm{test},n}]=1,
\]
and independence gives
\[
\begin{aligned}
1+\chi^2(Q_{\mathrm{test},1}^{(n)}\Vert Q_{\mathrm{test},0}^{(n)})
&=
E_{Q_{\mathrm{test},0}^{(n)}}[L_{\mathrm{test},n}^2]\\
&=
\left(
1+\chi^2(Q_{\mathrm{test},1}\Vert Q_{\mathrm{test},0})
\right)^n.
\end{aligned}
\]
Here the product divergence is defined by the same finite-sum formula on
\(\mathcal Z_{\mathrm{test}}^n\). Since \(1+t\le e^t\) for \(t\ge0\),
\[
\chi^2(Q_{\mathrm{test},1}^{(n)}\Vert Q_{\mathrm{test},0}^{(n)})
\le e^{1/8}-1<1.
\]
For the numerical inequality, monotonicity gives \(e^{1/8}\le e^{1/2}\), while
\[
(e^{1/2})^2=e<3<4,
\]
so \(e^{1/8}<2\).

Define total variation for these sample laws by
\[
\operatorname{TV}(Q_{\mathrm{test},1}^{(n)},Q_{\mathrm{test},0}^{(n)})
:=
\frac12
\sum_{\mathbf o\in\mathcal Z_{\mathrm{test}}^n}
\left|
Q_{\mathrm{test},1}^{(n)}(\{\mathbf o\})
-
Q_{\mathrm{test},0}^{(n)}(\{\mathbf o\})
\right|.
\]
Cauchy--Schwarz applied to the likelihood ratio yields
\[
\begin{aligned}
\operatorname{TV}(Q_{\mathrm{test},1}^{(n)},Q_{\mathrm{test},0}^{(n)})
&=
\frac12E_{Q_{\mathrm{test},0}^{(n)}}|L_{\mathrm{test},n}-1|\\
&\le
\frac12\sqrt{
E_{Q_{\mathrm{test},0}^{(n)}}[(L_{\mathrm{test},n}-1)^2]
}\\
&=
\frac12\sqrt{
\chi^2(Q_{\mathrm{test},1}^{(n)}\Vert Q_{\mathrm{test},0}^{(n)})
}
\le\frac12.
\end{aligned}
\]
Consequently, for every measurable \(\mathcal E\subseteq\mathcal O_d^n\),
\[
\begin{aligned}
Q_{\mathrm{test},1}^{(n)}(\mathcal E^c)
+Q_{\mathrm{test},0}^{(n)}(\mathcal E)
&=
1-\left(
Q_{\mathrm{test},1}^{(n)}(\mathcal E)
-Q_{\mathrm{test},0}^{(n)}(\mathcal E)
\right)\\
&\ge
1-\operatorname{TV}(Q_{\mathrm{test},1}^{(n)},Q_{\mathrm{test},0}^{(n)})
\ge\frac12.
\end{aligned}
\]
% lean: CausalSmith.Stat.MarRareqLogfrontier.complete_arrival_product_testing_half

\item Take any parameter-independent bounded kernel \(\mathsf T\) in the decision class, and define its barycenter by
\[
b_{\mathsf T}(\mathbf o)
:=
\int_{[-1,1]}h\,\mathsf T(\mathbf o,dh).
\]
By \cref{obj:lem:randomized-kernel-barycenter}, this is a measurable map into \([-1,1]\), and
\[
E_{P_{\mathrm{test},\iota}}\!\left[
\left(b_{\mathsf T}(\mathbf O_n)-\tau(P_{\mathrm{test},\iota})\right)^2
\right]
\le
\mathcal R_{P_{\mathrm{test},\iota}}(\mathsf T),
\qquad \iota\in\{0,1\}.
\]
Define the two large-error events and the nearest-target test by
\[
\mathcal E_{\mathrm{test},\iota}
:=
\left\{\mathbf o:
\frac{\delta_{\mathrm{gap}}}{2}
\le
\left|b_{\mathsf T}(\mathbf o)-\tau(P_{\mathrm{test},\iota})\right|
\right\},
\qquad \iota\in\{0,1\},
\]
and
\[
\mathcal A_{\mathrm{test}}
:=
\left\{\mathbf o:
\left|b_{\mathsf T}(\mathbf o)-\tau(P_{\mathrm{test},1})\right|
\le
\left|b_{\mathsf T}(\mathbf o)-\tau(P_{\mathrm{test},0})\right|
\right\}.
\]
The target separation and the triangle inequality give, for every sample,
\[
\delta_{\mathrm{gap}}
\le
\left|b_{\mathsf T}(\mathbf o)-\tau(P_{\mathrm{test},0})\right|
+
\left|b_{\mathsf T}(\mathbf o)-\tau(P_{\mathrm{test},1})\right|.
\]
On \(\mathcal A_{\mathrm{test}}\), the first error is at least
\(\delta_{\mathrm{gap}}/2\); on its complement, the second error is at least
\(\delta_{\mathrm{gap}}/2\). Hence
\[
\mathcal A_{\mathrm{test}}\subseteq\mathcal E_{\mathrm{test},0},
\qquad
\mathcal A_{\mathrm{test}}^c\subseteq\mathcal E_{\mathrm{test},1}.
\]
Applying the testing bound and then event monotonicity gives
\[
\frac12
\le
Q_{\mathrm{test},1}^{(n)}(\mathcal A_{\mathrm{test}}^c)
+
Q_{\mathrm{test},0}^{(n)}(\mathcal A_{\mathrm{test}})
\le
Q_{\mathrm{test},1}^{(n)}(\mathcal E_{\mathrm{test},1})
+
Q_{\mathrm{test},0}^{(n)}(\mathcal E_{\mathrm{test},0}),
\]
so
\[
\max_{\iota\in\{0,1\}}
Q_{\mathrm{test},\iota}^{(n)}(\mathcal E_{\mathrm{test},\iota})
\ge\frac14.
\]
Integrating squared error over each large-error event yields
\[
\left(\frac{\delta_{\mathrm{gap}}}{2}\right)^2
Q_{\mathrm{test},\iota}^{(n)}(\mathcal E_{\mathrm{test},\iota})
\le
E_{P_{\mathrm{test},\iota}}\!\left[
\left(b_{\mathsf T}(\mathbf O_n)-\tau(P_{\mathrm{test},\iota})\right)^2
\right].
\]
Whether the probability for \(\iota=0\) is smaller than that for \(\iota=1\), or conversely, multiplication by the nonnegative factor
\((\delta_{\mathrm{gap}}/2)^2\) preserves their ordering. Therefore
\[
\begin{aligned}
\frac{\delta_{\mathrm{gap}}^2}{16}
&\le
\left(\frac{\delta_{\mathrm{gap}}}{2}\right)^2
\max_{\iota\in\{0,1\}}
Q_{\mathrm{test},\iota}^{(n)}(\mathcal E_{\mathrm{test},\iota})\\
&=
\max_{\iota\in\{0,1\}}
\left\{
\left(\frac{\delta_{\mathrm{gap}}}{2}\right)^2
Q_{\mathrm{test},\iota}^{(n)}(\mathcal E_{\mathrm{test},\iota})
\right\}\\
&\le
\max_{\iota\in\{0,1\}}
E_{P_{\mathrm{test},\iota}}\!\left[
\left(b_{\mathsf T}(\mathbf O_n)-\tau(P_{\mathrm{test},\iota})\right)^2
\right]\\
&\le
\max_{\iota\in\{0,1\}}
\mathcal R_{P_{\mathrm{test},\iota}}(\mathsf T).
\end{aligned}
\]
For every bounded kernel and every full-data law, bounded actions and
\(\tau(P)\in[-1,1]\) also give
\[
0\le\mathcal R_P(\mathsf T)\le4.
\]
Thus the worst-case risks are finite. The decision class is nonempty because it contains the zero estimator, and the complete-arrival class contains the constructed pair. Taking the supremum over that class and then the infimum over kernels gives
\[
\mathfrak R^{+}_{n,d,1}
\ge
\frac{\delta_{\mathrm{gap}}^2}{16}
=
\frac1{256n}.
\]
% lean: CausalSmith.Stat.MarRareqLogfrontier.complete_arrival_two_point_minimax_floor

\item Finally, every complete-arrival law belongs to the model with arrival floor \(q\). Indeed, all other conditions are shared, and for each occupied cell of such a law,
\[
q\,p_{axs}(P)
\le p_{axs}(P)
\le P(A=a,X=x,S=s,R=1).
\]
Together with \(0<q\le1\), this proves
\[
\mathcal M^{+}_{n,d,1}\subseteq\mathcal M^{+}_{n,d,q}.
\]
For every kernel \(\mathsf T\), class inclusion gives
\[
\sup_{P\in\mathcal M^{+}_{n,d,1}}\mathcal R_P(\mathsf T)
\le
\sup_{P\in\mathcal M^{+}_{n,d,q}}\mathcal R_P(\mathsf T).
\]
The risks lie in \([0,4]\), and both classes contain the complete-arrival pair. Taking infima over the same nonempty decision class therefore yields
\[
\frac1{256n}
\le
\mathfrak R^{+}_{n,d,1}
\le
\mathfrak R^{+}_{n,d,q}.
\]
Combining this with the three upper bounds proves
\[
\frac1{256n}
\le
\mathfrak R^{+}_{n,d,q}
\le
\min\left\{
1,\,
\overline C\,r_{n,d,q},\,
\frac4n+(1-q)^2
\right\}.
\]
The first step supplies the mixed-count envelope and its auxiliary-risk comparison, and the observed-contrast calculation supplies the asserted uniform witness bound.
% lean: CausalSmith.Stat.MarRareqLogfrontier.unrestricted_near_complete_arrival_envelope
\end{enumerate}
\end{proof}

\begin{proof}[Proof of \cref{obj:thm:unrestricted-large-alphabet-deficit-lower}]
\begin{enumerate}
\item Put
\[
\delta:=1-q\in[0,1).
\]
Since \(d\ge1024n^2\ge1\), the hypotheses of
\cref{obj:thm:unrestricted-near-complete-arrival-envelope} hold. Its lower bound and the last term of its upper envelope give
\[
\frac1{256n}\le\mathfrak R^{+}_{n,d,q},
\qquad
\mathfrak R^{+}_{n,d,q}\le\frac4n+\delta^2.
\]
If \(\delta^2\le n^{-1}\), then
\[
\frac{\delta^2}{1024}
\le\frac1{1024n}
\le\frac1{256n}
\le\mathfrak R^{+}_{n,d,q}.
\]
This case includes \(q=1\). For the remaining case, assume
\[
\delta^2>n^{-1},
\qquad\text{so in particular}\qquad n^{-1}\le\delta^2.
\]
We establish the deficit lower bound in this regime. % lean: CausalSmith.Stat.MarRareqLogfrontier.unrestricted_large_alphabet_deficit_lower

\item Write an ordered observed sample as
\[
\mathbf o=(o_1,\ldots,o_n)
\in\left(\{0,\ldots,d-1\}\times\{0,1\}^4\right)^n,
\]
with each record ordered as \((X,A,S,R,RY)\). The decision class contains the constant zero kernel
\[
\mathsf T_0(\mathbf o,\cdot):=\operatorname{Dirac}(0).
\]
Fix an arbitrary \(\mathsf T\in\mathcal T_n\), the bounded, parameter-independent kernel class of
\cref{obj:def:unrestricted-minimax-risk}, and define its risk under a full-data law \(P\) by
\[
\mathcal R_{\mathsf T}(P)
:=
\int\!\int_{[-1,1]}
(w-\tau(P))^2\,
\mathsf T(\mathbf o,dw)\,
\mathcal L_P(\mathbf O_n)(d\mathbf o).
\]
Define the separation between the two target centers by
\[
h_{\mathrm{sep}}
:=\frac12-\frac q{1+q}
=\frac{\delta}{2(1+q)}.
\]
Because \(1<1+q\le2\) and \(\delta\ge0\), this identity gives
\[
\frac{\delta}{4}\le h_{\mathrm{sep}}\le\frac{\delta}{2},
\qquad h_{\mathrm{sep}}\ge0.
\]

Construct a comparison law \(P_\star\) by taking a uniform baseline, zero surrogates and control outcome, and a Bernoulli treated potential outcome:
\[
X\sim\operatorname{Unif}\{0,\ldots,d-1\},
\qquad
S(0)=S(1)=Y(0)=0,
\qquad
Y(1)\mid X\sim\operatorname{Bernoulli}\!\left(\frac q{1+q}\right).
\]
Draw \(A\sim\operatorname{Bernoulli}(1/2)\) independently of the baseline and potential variables, set
\[
S=S(A),\qquad Y=Y(A),
\]
and draw the arrival indicator conditionally independently of \(Y(1)\), with
\[
P_\star(R=1\mid X,A,Y(1))
=
\begin{cases}
1,&A=0,\\
(1+q)/2,&A=1.
\end{cases}
\]
These distributions are well defined because
\[
0<\frac q{1+q}\le\frac12,
\qquad
q\le\frac{1+q}{2}\le1.
\]
Evaluating the average treatment effect of
\cref{obj:def:ate-functional} yields
\[
\tau(P_\star)
=E_{P_\star}[Y(1)-Y(0)]
=\frac q{1+q}
=\frac12-h_{\mathrm{sep}}.
\]
% lean: CausalSmith.Stat.MarRareqLogfrontier.largeAlphabet_center_separation
% lean: CausalSmith.Stat.MarRareqLogfrontier.largeAlphabetComparisonLaw_ate

\item For each
\[
\xi=(\xi_x)_{x=0}^{d-1}\in\{0,1\}^{\{0,\ldots,d-1\}},
\]
construct a law \(P_\xi\) with the same uniform baseline and independent fair treatment assignment, and with
\[
S(0)=S(1)=Y(0)=0,\qquad Y(1)=\xi_X,
\qquad S=S(A),\qquad Y=Y(A).
\]
Its arrival draw satisfies, for \(x\in\{0,\ldots,d-1\}\) and \(a\in\{0,1\}\),
\[
P_\xi(R=1\mid X=x,A=a)=1-a\delta\xi_x.
\]
Thus every arrival probability is either \(q\) or \(1\), and
\[
\tau(P_\xi)
=\frac1d\sum_{x=0}^{d-1}\xi_x.
\]
Put the product fair-coin prior on these vectors:
\[
\nu_{\mathrm{cell}}
:=\bigotimes_{x=0}^{d-1}\operatorname{Bernoulli}(1/2),
\qquad
\nu_{\mathrm{cell}}(\{\xi\})=2^{-d}.
\]
Define the sample laws and their mixture by
\[
Q_\xi:=\mathcal L_{P_\xi}(\mathbf O_n),
\qquad
Q_\star:=\mathcal L_{P_\star}(\mathbf O_n),
\qquad
Q_{\mathrm{mix}}
:=\int Q_\xi\,\nu_{\mathrm{cell}}(d\xi)
=\sum_\xi2^{-d}Q_\xi.
\]
For each fixed full-data law, the sample consists of \(n\) independent observed records as in \cref{obj:ass:iid-sampling}. The mixture draws \(\xi\) once for the entire sample.

Define the test event, the exceptional prior vectors, and the two risks by
\[
E_{\mathrm{test}}
:=
\left\{(\mathbf o,w):
w\le\frac12-\frac{h_{\mathrm{sep}}}{2}\right\},
\qquad
w\in\mathbb R,
\]
\[
F_{\mathrm{bad}}
:=
\left\{\xi:
\left|\tau(P_\xi)-\frac12\right|>
\frac{h_{\mathrm{sep}}}{4}\right\},
\]
\[
\mathcal R_{\mathrm{prior}}(\mathsf T)
:=\int\mathcal R_{\mathsf T}(P_\xi)\,\nu_{\mathrm{cell}}(d\xi),
\qquad
\mathcal R_\star(\mathsf T):=\mathcal R_{\mathsf T}(P_\star).
\]
The notation \(Q\ltimes\mathsf T\) denotes the joint sample--action law used in
\cref{obj:lem:randomized-kernel-tv-comparison}.

If \(\xi\notin F_{\mathrm{bad}}\) and \((\mathbf o,w)\in E_{\mathrm{test}}\), then
\[
\tau(P_\xi)-w
\ge
\left(\frac12-\frac{h_{\mathrm{sep}}}{4}\right)
-\left(\frac12-\frac{h_{\mathrm{sep}}}{2}\right)
=\frac{h_{\mathrm{sep}}}{4}.
\]
Consequently, for every \(\xi\),
\[
\frac{h_{\mathrm{sep}}^2}{16}
(Q_\xi\ltimes\mathsf T)(E_{\mathrm{test}})
\le
\mathcal R_{\mathsf T}(P_\xi)
+\frac{h_{\mathrm{sep}}^2}{16}
\mathbf1_{F_{\mathrm{bad}}}(\xi).
\]
For exceptional vectors this follows from the probability bound by one and nonnegativity of risk; for the other vectors it follows by integrating the displayed pointwise loss bound. Since the same kernel is used for every \(\xi\), finite averaging gives
\[
(Q_{\mathrm{mix}}\ltimes\mathsf T)(E_{\mathrm{test}})
=
\int
(Q_\xi\ltimes\mathsf T)(E_{\mathrm{test}})
\,\nu_{\mathrm{cell}}(d\xi).
\]
Averaging the member-wise risk inequality therefore yields
\[
\frac{h_{\mathrm{sep}}^2}{16}
(Q_{\mathrm{mix}}\ltimes\mathsf T)(E_{\mathrm{test}})
\le
\mathcal R_{\mathrm{prior}}(\mathsf T)
+\frac{h_{\mathrm{sep}}^2}{16}\nu_{\mathrm{cell}}(F_{\mathrm{bad}}).
\]

Under \(P_\star\), an action in \(E_{\mathrm{test}}^c\) satisfies
\[
w-\tau(P_\star)
>
\left(\frac12-\frac{h_{\mathrm{sep}}}{2}\right)
-\left(\frac12-h_{\mathrm{sep}}\right)
=\frac{h_{\mathrm{sep}}}{2}.
\]
In particular, its squared loss is at least \(h_{\mathrm{sep}}^2/16\), so
\[
\frac{h_{\mathrm{sep}}^2}{16}
(Q_\star\ltimes\mathsf T)(E_{\mathrm{test}}^c)
\le\mathcal R_\star(\mathsf T).
\]
% lean: CausalSmith.Stat.MarRareqLogfrontier.largeAlphabetLaw_ate
% lean: CausalSmith.Stat.MarRareqLogfrontier.largeAlphabet_prior_test_risk
% lean: CausalSmith.Stat.MarRareqLogfrontier.largeAlphabet_comparison_test_risk

\item We next compare the two sample laws. For a single treated visit to any baseline cell \(x\), averaging the fair coordinate \(\xi_x\) gives the following observed-mark probabilities:
\[
\begin{array}{c|c|c}
(R,RY)
&
\displaystyle\int P_\xi((R,RY)\mid X=x,A=1)\,
\nu_{\mathrm{cell}}(d\xi)
&
P_\star((R,RY)\mid X=x,A=1)
\\ \hline
(1,1)&q/2&
\displaystyle\frac{1+q}{2}\frac q{1+q}=q/2
\\
(1,0)&1/2&
\displaystyle\frac{1+q}{2}\left(1-\frac q{1+q}\right)=1/2
\\
(0,0)&\delta/2&
\displaystyle1-\frac{1+q}{2}=\delta/2
\\
(0,1)&0&0
\end{array}
\]
In the prior column, the first probability comes from \(\xi_x=1\) and arrival, the second from \(\xi_x=0\), and the third from \(\xi_x=1\) and nonarrival. A control visit has mark \((R,RY)=(1,0)\) with probability one under every law. All surrogates equal zero.

Define the event of distinct treated baseline cells by
\[
E_{\mathrm{distinct}}
:=
\bigcap_{1\le i<j\le n}
\{A_i\ne1\ \text{or}\ A_j\ne1\ \text{or}\ X_i\ne X_j\}.
\]
For any sample atom \(\mathbf o=(o_1,\ldots,o_n)\) in this event, each treated-record probability depends on a different prior coordinate, while each control-record probability is constant in \(\xi\). Independence of the prior coordinates and the one-visit calculation therefore give
\[
\begin{aligned}
Q_{\mathrm{mix}}(\{\mathbf o\})
&=
\int
\prod_{i=1}^n
\mathcal L_{P_\xi}(O_1)(\{o_i\})\,
\nu_{\mathrm{cell}}(d\xi)
\\
&=
\prod_{i=1}^n
\int
\mathcal L_{P_\xi}(O_1)(\{o_i\})\,
\nu_{\mathrm{cell}}(d\xi)
\\
&=
\prod_{i=1}^n
\mathcal L_{P_\star}(O_1)(\{o_i\})
=
Q_\star(\{\mathbf o\}).
\end{aligned}
\]
The equality also covers atoms with impossible surrogate or arrived-outcome marks, whose probabilities are zero under both constructions. % lean: CausalSmith.Stat.MarRareqLogfrontier.largeAlphabetSampleMixture_matching

\item Under each \(P_\xi\) and under \(P_\star\), the probability of any treated baseline cell is
\[
P_\xi(A=1,X=x)=P_\star(A=1,X=x)=\frac1{2d},
\qquad x\in\{0,\ldots,d-1\}.
\]
Thus independence of distinct observations gives, for \(1\le i<j\le n\),
\[
\begin{aligned}
Q_\xi\{A_i=A_j=1,\ X_i=X_j\}
&\le
\sum_{x=0}^{d-1}\left(\frac1{2d}\right)^2
=\frac1{4d},
\\
Q_\star\{A_i=A_j=1,\ X_i=X_j\}
&\le\frac1{4d}.
\end{aligned}
\]
Taking a union over the \(\binom n2\) pairs yields
\[
Q_\xi(E_{\mathrm{distinct}}^c),
\ Q_\star(E_{\mathrm{distinct}}^c)
\le
\binom n2\frac1{4d}
\le\frac{n^2}{8d}
\le\frac1{8192}.
\]
When \(n=1\), the pair set is empty and these collision probabilities are zero. Averaging the first bound over the prior gives
\[
Q_{\mathrm{mix}}(E_{\mathrm{distinct}}^c)
=
\int Q_\xi(E_{\mathrm{distinct}}^c)\,
\nu_{\mathrm{cell}}(d\xi)
\le\frac1{8192}.
\]

On the finite sample space, agreement of the atom masses on
\(E_{\mathrm{distinct}}\) implies
\[
\begin{aligned}
\operatorname{TV}(Q_{\mathrm{mix}},Q_\star)
&\le
\frac12\sum_{\mathbf o}
\left|Q_{\mathrm{mix}}(\{\mathbf o\})
-Q_\star(\{\mathbf o\})\right|
\\
&\le
\frac12\left\{
Q_{\mathrm{mix}}(E_{\mathrm{distinct}}^c)
+Q_\star(E_{\mathrm{distinct}}^c)
\right\}
\le\frac1{8192}.
\end{aligned}
\]
Indeed, the summands vanish on \(E_{\mathrm{distinct}}\), and elsewhere each absolute difference is bounded by the sum of the two atom masses. Applying
\cref{obj:lem:randomized-kernel-tv-comparison} to the common kernel and the measurable event \(E_{\mathrm{test}}\) now gives
\[
1-\frac1{8192}
\le
1-\operatorname{TV}(Q_{\mathrm{mix}},Q_\star)
\le
(Q_{\mathrm{mix}}\ltimes\mathsf T)(E_{\mathrm{test}})
+
(Q_\star\ltimes\mathsf T)(E_{\mathrm{test}}^c).
\]
% lean: CausalSmith.Stat.MarRareqLogfrontier.largeAlphabet_actual_collision_testing

\item Under \(\nu_{\mathrm{cell}}\), the independent fair coordinates have
\[
E_{\nu_{\mathrm{cell}}}[\xi_x]=\frac12,
\qquad
\operatorname{Var}_{\nu_{\mathrm{cell}}}(\xi_x)=\frac14.
\]
Consequently,
\[
E_{\nu_{\mathrm{cell}}}\!\left[\sum_{x=0}^{d-1}\xi_x\right]
=\frac d2,
\qquad
\operatorname{Var}_{\nu_{\mathrm{cell}}}
\!\left(\sum_{x=0}^{d-1}\xi_x\right)
=\frac d4.
\]
The regime \(n^{-1}\le\delta^2\), together with \(\delta\ge0\), implies
\[
\delta>0,\qquad h_{\mathrm{sep}}\ge\frac{\delta}{4}>0,
\qquad
1\le n\delta^2,\qquad
\delta^2\le16h_{\mathrm{sep}}^2.
\]
These inequalities and the alphabet-size condition give the concentration budget
\[
d h_{\mathrm{sep}}^2
\ge1024n^2h_{\mathrm{sep}}^2
\ge64n^2\delta^2
\ge64n
\ge64.
\]
Since \(\tau(P_\xi)=d^{-1}\sum_x\xi_x\), Chebyshev's inequality applied to the sum of the prior coordinates yields
\[
\begin{aligned}
\nu_{\mathrm{cell}}(F_{\mathrm{bad}})
&=
\nu_{\mathrm{cell}}
\left\{
\left|\sum_{x=0}^{d-1}\xi_x-\frac d2\right|
>\frac{d h_{\mathrm{sep}}}{4}
\right\}
\\
&\le
\frac{d/4}{(d h_{\mathrm{sep}}/4)^2}
=
\frac4{d h_{\mathrm{sep}}^2}
\le\frac1{16}.
\end{aligned}
\]
% lean: CausalSmith.Stat.MarRareqLogfrontier.largeAlphabet_center_prior_concentration

\item Multiplying the testing and concentration bounds by the nonnegative factor \(h_{\mathrm{sep}}^2/16\) gives
\[
\frac{h_{\mathrm{sep}}^2}{16}
\left(1-\frac1{8192}\right)
\le
\frac{h_{\mathrm{sep}}^2}{16}
\left\{
(Q_{\mathrm{mix}}\ltimes\mathsf T)(E_{\mathrm{test}})
+
(Q_\star\ltimes\mathsf T)(E_{\mathrm{test}}^c)
\right\},
\]
\[
\frac{h_{\mathrm{sep}}^2}{16}
\nu_{\mathrm{cell}}(F_{\mathrm{bad}})
\le\frac{h_{\mathrm{sep}}^2}{256}.
\]
Combining these with the two test-risk inequalities established above, we obtain
\[
\begin{aligned}
\mathcal R_{\mathrm{prior}}(\mathsf T)
+\mathcal R_\star(\mathsf T)
&\ge
\frac{h_{\mathrm{sep}}^2}{16}
\left\{
1-\frac1{8192}-\frac1{16}
\right\}
\\
&\ge\frac{h_{\mathrm{sep}}^2}{32},
\end{aligned}
\]
where the last inequality uses
\(1-1/8192-1/16\ge1/2\). % lean: CausalSmith.Stat.MarRareqLogfrontier.largeAlphabet_deficit_minimax_lower

\item To pass to worst-case risk, define
\[
\mathcal R_{\mathrm{sup}}(\mathsf T)
:=
\sup_{P\in\mathcal M^{+}_{n,d,q}}
\mathcal R_{\mathsf T}(P).
\]
Binary potential outcomes give \(\tau(P)\in[-1,1]\). Since the kernel takes values in \([-1,1]\), its loss and risk satisfy
\[
0\le(w-\tau(P))^2\le4
\quad(w\in[-1,1]),
\qquad
0\le\mathcal R_{\mathsf T}(P)\le4.
\]
Thus the member-law risks are bounded above.

Every constructed law belongs to the model class in
\cref{obj:def:unrestricted-arrival-model-class}. The independent fair treatment draw verifies
\cref{obj:ass:randomized-independence,obj:ass:balanced-randomization}, and the defining identities for the realized variables verify
\cref{obj:ass:surrogate-consistency,obj:ass:outcome-consistency}.
Within an occupied cell, the outcome under \(P_\xi\) is deterministic, while under \(P_\star\) its arrival draw is independent of the outcome; these properties verify
\cref{obj:ass:arrival-mar}. For every \(a\in\{0,1\}\) and \(x\in\{0,\ldots,d-1\}\), the cell probabilities of
\cref{obj:synth_16} are
\[
p_{ax0}(P_\xi)=p_{ax0}(P_\star)=\frac1{2d},
\qquad
p_{ax1}(P_\xi)=p_{ax1}(P_\star)=0.
\]
On the occupied cells, their arrival probabilities are
\[
\rho_{ax0}(P_\xi)=1-a\delta\xi_x\ge q,
\qquad
\rho_{ax0}(P_\star)
=
\begin{cases}
1,&a=0,\\
(1+q)/2,&a=1,
\end{cases}
\ge q.
\]
This verifies \cref{obj:ass:occupied-cell-arrival}, and hence
\[
P_\xi,P_\star\in\mathcal M^{+}_{n,d,q}.
\]

Each member-law risk is therefore at most
\(\mathcal R_{\mathrm{sup}}(\mathsf T)\). Since
\(\nu_{\mathrm{cell}}\) is a probability law,
\[
\mathcal R_{\mathrm{prior}}(\mathsf T)
\le
\int\mathcal R_{\mathrm{sup}}(\mathsf T)\,
\nu_{\mathrm{cell}}(d\xi)
=
\mathcal R_{\mathrm{sup}}(\mathsf T),
\qquad
\mathcal R_\star(\mathsf T)
\le\mathcal R_{\mathrm{sup}}(\mathsf T).
\]
It follows that
\[
2\mathcal R_{\mathrm{sup}}(\mathsf T)
\ge
\mathcal R_{\mathrm{prior}}(\mathsf T)+\mathcal R_\star(\mathsf T)
\ge\frac{h_{\mathrm{sep}}^2}{32}.
\]
Using \(h_{\mathrm{sep}}\ge\delta/4\ge0\), we conclude
\[
\mathcal R_{\mathrm{sup}}(\mathsf T)
\ge\frac{h_{\mathrm{sep}}^2}{64}
\ge\frac{\delta^2}{1024}.
\]
This holds for every kernel in the stated decision class. Taking its infimum, as in
\cref{obj:def:unrestricted-minimax-risk}, proves
\[
\frac{\delta^2}{1024}\le\mathfrak R^{+}_{n,d,q}
\]
in the remaining regime. % lean: CausalSmith.Stat.MarRareqLogfrontier.largeAlphabet_deficit_minimax_lower

\item The two cases establish the deficit lower bound throughout the stated parameter range. Combining it with the initial bounds gives
\[
\max\left\{\frac1{256n},\frac{\delta^2}{1024}\right\}
\le\mathfrak R^{+}_{n,d,q}
\le\frac4n+\delta^2.
\]
Adding the two lower inequalities and dividing by two yields
\[
\mathfrak R^{+}_{n,d,q}
\ge
\frac1{512n}+\frac{\delta^2}{2048}
\ge
\frac1{2048}\left\{\frac1n+\delta^2\right\}.
\]
Finally, \(\delta^2\ge0\) gives
\[
\mathfrak R^{+}_{n,d,q}
\le\frac4n+\delta^2
\le4\left\{\frac1n+\delta^2\right\}.
\]
Substituting \(\delta=1-q\) proves the asserted envelope. % lean: CausalSmith.Stat.MarRareqLogfrontier.unrestricted_large_alphabet_deficit_lower
\end{enumerate}
\end{proof}

\begin{proof}[Proof of \cref{obj:prop:unrestricted-uniform-near-complete-elbow}]
\begin{enumerate}
\item We first record two bounds. By \cref{obj:thm:unrestricted-near-complete-arrival-envelope}, for every \(n,d\ge1\),
\[
\mathfrak R^{+}_{n,d,q_n}\le \frac4n+(1-q_n)^2.
\]
For \(n\ge1\), the alphabet size \(1024n^2\) satisfies
\[
1\le1024n^2.
\]
Thus \cref{obj:thm:unrestricted-large-alphabet-deficit-lower}, applied at this alphabet size, supplies
\[
\frac1{2048}\left\{\frac1n+(1-q_n)^2\right\}
\le \mathfrak R^{+}_{n,1024n^2,q_n}.
\]
Multiplication by \(n>0\), using
\[
n\cdot\frac1{2048}\left\{\frac1n+(1-q_n)^2\right\}
=\frac1{2048}\left\{1+n(1-q_n)^2\right\},
\]
gives
\[
\frac1{2048}\left\{1+n(1-q_n)^2\right\}
\le n\mathfrak R^{+}_{n,1024n^2,q_n}.
\]
% lean: CausalSmith.Stat.MarRareqLogfrontier.unrestricted_uniform_near_complete_elbow

\item Multiplying the upper bound from the preceding step by \(n>0\) yields, for every \(d\ge1\),
\[
n\mathfrak R^{+}_{n,d,q_n}
\le n\left\{\frac4n+(1-q_n)^2\right\}
=4+n(1-q_n)^2.
\]
For any real \(b\) satisfying the stated strict inequality, the same step gives
\[
b<
\frac1{2048}\left\{1+n(1-q_n)^2\right\}
\le n\mathfrak R^{+}_{n,1024n^2,q_n}.
\]
Hence \(d=1024n^2\) is the required integer witness.
% lean: CausalSmith.Stat.MarRareqLogfrontier.unrestricted_uniform_near_complete_elbow

\item Suppose first that the arrival-rate condition holds with \(M\ge0\). For all sufficiently large \(n\ge1\),
\[
1-q_n\le\frac{M}{\sqrt n}.
\]
Since \(q_n\le1\) and \(\sqrt n>0\), this implies
\[
0\le(1-q_n)\sqrt n\le M.
\]
Squaring and using \((\sqrt n)^2=n\), we obtain
\[
n(1-q_n)^2
=\bigl((1-q_n)\sqrt n\bigr)^2
\le M^2,
\qquad
(1-q_n)^2\le\frac{M^2}{n}.
\]
Consequently, for every \(d\ge1\) and all such \(n\),
\[
\mathfrak R^{+}_{n,d,q_n}
\le\frac4n+(1-q_n)^2
\le\frac4n+\frac{M^2}{n}
=\frac{4+M^2}{n}.
\]
This proves the uniform risk bound with \(D=4+M^2\).

Conversely, suppose the uniform risk bound holds with a real constant \(D\). Define the nonnegative real constant
\[
M_{\mathrm{risk}}:=\sqrt{\max\{2048D,0\}}.
\]
For all sufficiently large \(n\ge1\), applying that bound at \(d=1024n^2\) and multiplying by \(n\) gives
\[
n\mathfrak R^{+}_{n,1024n^2,q_n}\le D.
\]
Combining this with the lower bound from the first step yields
\[
\frac1{2048}\left\{1+n(1-q_n)^2\right\}\le D.
\]
In particular,
\[
\bigl(\sqrt n(1-q_n)\bigr)^2
=n(1-q_n)^2
\le2048D
\le\max\{2048D,0\}
=M_{\mathrm{risk}}^2.
\]
Both \(\sqrt n(1-q_n)\) and \(M_{\mathrm{risk}}\) are nonnegative, so
\[
\sqrt n(1-q_n)\le M_{\mathrm{risk}},
\qquad
1-q_n\le\frac{M_{\mathrm{risk}}}{\sqrt n}.
\]
This is the arrival-rate condition.
% lean: CausalSmith.Stat.MarRareqLogfrontier.unrestricted_uniform_near_complete_elbow

\item The lower bound in \cref{obj:thm:unrestricted-near-complete-arrival-envelope} directly supplies, for every \(n,d\ge1\),
\[
\frac1{256n}\le\mathfrak R^{+}_{n,d,q_n}.
\]
% lean: CausalSmith.Stat.MarRareqLogfrontier.unrestricted_near_complete_arrival_envelope

\item Finally, assume that \(\bigl(\sqrt n(1-q_n)\bigr)_{n\ge0}\) is unbounded above, and fix \(C\in\mathbb R\). Suppose, for contradiction, that the claimed strict risk inequality holds for only finitely many integers \(n\). Choose an integer \(n_{\mathrm{tail}}\ge1\) such that
\[
n\mathfrak R^{+}_{n,1024n^2,q_n}\le C
\qquad\text{for every }n\ge n_{\mathrm{tail}}.
\]
Define the nonnegative real constant
\[
M_{\mathrm{tail}}:=\sqrt{\max\{2048C,0\}}.
\]
For every \(n\ge n_{\mathrm{tail}}\), the first step implies
\[
\frac1{2048}\left\{1+n(1-q_n)^2\right\}\le C,
\]
and hence
\[
\bigl(\sqrt n(1-q_n)\bigr)^2
=n(1-q_n)^2
\le2048C
\le\max\{2048C,0\}
=M_{\mathrm{tail}}^2.
\]
Nonnegativity therefore gives
\[
\sqrt n(1-q_n)\le M_{\mathrm{tail}}
\qquad(n\ge n_{\mathrm{tail}}).
\]
The remaining terms form a finite set. Explicitly, the finite maximum below exists and bounds every term:
\[
\sqrt n(1-q_n)
\le
\max\left(
\{M_{\mathrm{tail}}\}
\cup
\left\{\sqrt{n'}(1-q_{n'}):
n'\in\mathbb N,\ n'<n_{\mathrm{tail}}\right\}
\right)
\qquad(n\in\mathbb N).
\]
This contradicts unboundedness above. Thus there are infinitely many integers \(n\) for which
\[
C<n\mathfrak R^{+}_{n,1024n^2,q_n}.
\]
% lean: CausalSmith.Stat.MarRareqLogfrontier.unrestricted_uniform_near_complete_elbow
\end{enumerate}
\end{proof}

\begin{proof}[Proof of \cref{obj:prop:unrestricted-near-complete-phase-boundary}]
\begin{enumerate}
\item For all sufficiently large \(n\), we have
\[
n\ge1,
\qquad
1-q_n\le\frac{M}{\sqrt n}.
\]
Fix such an \(n\). Since \(d_n\ge1\) and \(0<q_n\le1\), \cref{obj:thm:unrestricted-near-complete-arrival-envelope} applies. Retaining the last term in its upper-bound minimum gives
\[
\frac{1}{256n}
\le \mathfrak R^{+}_{n,d_n,q_n}
\le \frac4n+(1-q_n)^2.
\]
% lean: CausalSmith.Stat.MarRareqLogfrontier.unrestricted_near_complete_arrival_envelope

\item Since \(n>0\), we have \(\sqrt n>0\), so multiplying the deficit bound by \(\sqrt n\) yields
\[
(1-q_n)\sqrt n\le M.
\]
Moreover, \(q_n\le1\) implies
\[
0\le(1-q_n)\sqrt n.
\]
Thus, using \(M\ge0\), we may square the inequality to obtain
\[
\bigl((1-q_n)\sqrt n\bigr)^2\le M^2.
\]
The identity \((\sqrt n)^2=n\) then gives
\[
n(1-q_n)^2
=
\bigl((1-q_n)\sqrt n\bigr)^2
\le M^2.
\]
Dividing by \(n>0\), we conclude that
\[
(1-q_n)^2\le\frac{M^2}{n}.
\]
% lean: CausalSmith.Stat.MarRareqLogfrontier.unrestricted_near_complete_phase_boundary

\item Substituting this bound into the upper estimate gives
\[
\mathfrak R^{+}_{n,d_n,q_n}
\le \frac4n+(1-q_n)^2
\le \frac4n+\frac{M^2}{n}
=\frac{4+M^2}{n}.
\]
Together with the lower estimate, this proves both asserted inequalities for all sufficiently large \(n\).
% lean: CausalSmith.Stat.MarRareqLogfrontier.unrestricted_near_complete_phase_boundary
\end{enumerate}
\end{proof}

\begin{proof}[Proof of \cref{obj:thm:unrestricted-complete-arrival-frontier}]
Fix integers \(n,d\ge1\), with sampling as in \cref{obj:ass:iid-sampling}.

\begin{enumerate}
\item Since \(0<1\le1\), specialize \cref{obj:thm:unrestricted-near-complete-arrival-envelope} to \(q=1\). Its lower bound and observed-outcome contrast bound give
\[
\frac{1}{256n}\le \mathfrak R^{+}_{n,d,1},
\qquad
\sup_{P\in\mathcal M^{+}_{n,d,1}}
E_P\!\left[
\left\{\widehat\tau^{\mathrm{HT}}_n(\mathbf O_n)-\tau(P)\right\}^{2}
\right]
\le \frac4n+(1-1)^2=\frac4n.
\]
% lean: CausalSmith.Stat.MarRareqLogfrontier.unrestricted_near_complete_arrival_envelope

\item The estimator in \cref{obj:def:complete-arrival-ht-estimator} is measurable, takes values in \([-1,1]\), and depends on the observed sample and the fixed parameters \(n,d\). Define its deterministic probability kernel by
\[
\mathsf T_{\mathrm{HT}}(\mathbf o,E_{\mathrm{out}})
:=
\mathbf 1\!\left\{
\widehat\tau^{\mathrm{HT}}_n(\mathbf o)\in E_{\mathrm{out}}
\right\},
\quad
\mathbf o\in
\bigl(\{0,\ldots,d-1\}\times\{0,1\}^{4}\bigr)^n,
\quad
E_{\mathrm{out}}\subseteq[-1,1]\ \text{Borel}.
\]
This parameter-independent kernel belongs to the decision class \(\mathcal T_n\) of \cref{obj:def:unrestricted-minimax-risk}. For every \(\mathsf T\in\mathcal T_n\) and \(P\in\mathcal M^{+}_{n,d,1}\), squared loss gives
\[
0\le
\int\!\int_{[-1,1]}
(h-\tau(P))^2\,
\mathsf T(\mathbf o,dh)\,
\mathcal L_P(\mathbf O_n)(d\mathbf o).
\]
Evaluating the defining minimax infimum at \(\mathsf T_{\mathrm{HT}}\) therefore yields
\[
\mathfrak R^{+}_{n,d,1}
\le
\sup_{P\in\mathcal M^{+}_{n,d,1}}
\int\!\int_{[-1,1]}
(h-\tau(P))^2\,
\mathsf T_{\mathrm{HT}}(\mathbf o,dh)\,
\mathcal L_P(\mathbf O_n)(d\mathbf o).
\]
By the definition of this kernel, for every such \(P\),
\[
\int\!\int_{[-1,1]}
(h-\tau(P))^2\,
\mathsf T_{\mathrm{HT}}(\mathbf o,dh)\,
\mathcal L_P(\mathbf O_n)(d\mathbf o)
=
E_P\!\left[
\left\{\widehat\tau^{\mathrm{HT}}_n(\mathbf O_n)-\tau(P)\right\}^{2}
\right].
\]
Consequently,
\[
\mathfrak R^{+}_{n,d,1}
\le
\sup_{P\in\mathcal M^{+}_{n,d,1}}
E_P\!\left[
\left\{\widehat\tau^{\mathrm{HT}}_n(\mathbf O_n)-\tau(P)\right\}^{2}
\right].
\]
% lean: CausalSmith.Stat.MarRareqLogfrontier.complete_arrival_frontier_direct

\item Combining these inequalities gives
\[
\frac{1}{256n}
\le \mathfrak R^{+}_{n,d,1}
\le
\sup_{P\in\mathcal M^{+}_{n,d,1}}
E_P\!\left[
\left\{\widehat\tau^{\mathrm{HT}}_n(\mathbf O_n)-\tau(P)\right\}^{2}
\right]
\le \frac4n.
\]
The constants \(1/256\) and \(4\) are independent of \(n\) and \(d\), establishing order \(n^{-1}\) uniformly over \(d\ge1\).
% lean: CausalSmith.Stat.MarRareqLogfrontier.unrestricted_complete_arrival_frontier
\end{enumerate}
\end{proof}

\begin{proof}[Proof of \cref{obj:prop:unrestricted-complete-arrival-phase-boundary}]
\begin{enumerate}
\item Choose the real comparison constants
\[
c:=\frac{1}{256},
\qquad C:=4.
\]
Then \(0<c\le C\). Fix any sequence of natural alphabet sizes \((d_n)_{n\ge0}\) satisfying \(d_n\ge1\) for every \(n\). Define the arrival-floor sequence and deficit constant by
\[
q_n:=1\quad(n\ge0),
\qquad M:=0.
\]
Thus \(M\ge0\) and \(q_n\in(0,1]\) for every \(n\). For every \(n\ge1\),
\[
1-q_n=0=\frac{M}{\sqrt n},
\]
so the arrival-deficit condition holds for all sufficiently large \(n\).
% lean: CausalSmith.Stat.MarRareqLogfrontier.unrestricted_complete_arrival_phase_boundary

\item Applying \cref{obj:prop:unrestricted-near-complete-phase-boundary} with these choices gives, for all sufficiently large \(n\),
\[
\frac{1}{256n}
\le \mathfrak R^{+}_{n,d_n,1}
\le \frac{4}{n},
\]
since \(q_n=1\) and \(4+M^2=4\). These are precisely the required bounds with the chosen \(c\) and \(C\). Both constants were fixed before choosing the alphabet-size sequence, so they are independent of that sequence.
% lean: CausalSmith.Stat.MarRareqLogfrontier.unrestricted_near_complete_phase_boundary
% lean: CausalSmith.Stat.MarRareqLogfrontier.unrestricted_complete_arrival_phase_boundary
\end{enumerate}
\end{proof}

\begin{proof}[Proof of \cref{obj:thm:matched-minimax-frontier}]
\begin{enumerate}
\item Choose universal constants \(\underline c\) and \(C_0\) supplied by
\cref{obj:thm:full-experiment-lower,obj:thm:mixed-count-upper}, respectively, and set
\[
0<\underline c,\qquad 0<C_0,\qquad
\overline C:=\max\{\underline c,C_0\}.
\]
Thus
\[
0<\underline c\le\overline C<\infty,
\qquad C_0\le\overline C.
\]
Fix \(n,d,q\) satisfying the hypotheses of the theorem.
By \cref{obj:thm:full-experiment-lower},
\[
\underline c\,r_{n,d,q}\le\mathfrak R_{n,d,q}.
\]
% lean: CausalSmith.Stat.MarRareqLogfrontier.matched_minimax_frontier

\item We next construct an admissible decision kernel from the deterministic estimator in
\cref{obj:def:mixed-count-estimator}. For each observed sample \(o\), its defining finite average is
\[
\widehat\tau^{\mathrm{mix}}_{n,d,q}(o)
=
\sum_{\nu=0}^{n}
\Pr_\omega(K_0=\nu)\,
E_\omega\!\left[
\widetilde\tau(o;\omega)\mid K_0=\nu
\right].
\]
The omitted Poisson tail contributes zero by the estimator's definition.
Each conditional expectation is the uniform average over the \(3^\nu\) stream assignments.
Clipping, together with the zero-output branches, gives
\[
-1\le \widetilde\tau(o;\omega)\le1,
\qquad
-1\le
E_\omega\!\left[
\widetilde\tau(o;\omega)\mid K_0=\nu
\right]
\le1
\quad(0\le\nu\le n).
\]
The prefix weights satisfy
\[
\Pr_\omega(K_0=\nu)\ge0,\qquad
\sum_{\nu=0}^{n}\Pr_\omega(K_0=\nu)
=\Pr_\omega(K_0\le n)\le1.
\]
Consequently, at every observed sample,
\[
-1
\le-\Pr_\omega(K_0\le n)
\le\widehat\tau^{\mathrm{mix}}_{n,d,q}(o)
\le\Pr_\omega(K_0\le n)
\le1.
\]
The observed-sample space is finite with measurable singletons, so this deterministic map is measurable. Its Dirac kernel,
\[
\mathsf T_{\mathrm{mix}}(o,dh)
:=
\delta_{\widehat\tau^{\mathrm{mix}}_{n,d,q}(o)}(dh),
\]
therefore belongs to the decision class \(\mathcal T_n\) of
\cref{obj:def:minimax-risk}.
% lean: CausalSmith.Stat.MarRareqLogfrontier.mixed_count_estimator_regular

\item For every \(\mathsf T\in\mathcal T_n\) and \(P\in\mathcal M_{n,d,q}\), the joint squared-error risk is nonnegative:
\[
E_P\!\left[(\mathsf T-\tau(P))^2\right]
=
\int\!\int_{[-1,1]}
(h-\tau(P))^2\,\mathsf T(o,dh)\,
\mathcal L_P(\mathbf O_n)(do)
\ge0,
\]
because both integrals have nonnegative integrands. For the Dirac kernel just constructed,
\[
E_P\!\left[(\mathsf T_{\mathrm{mix}}-\tau(P))^2\right]
=
E_P\!\left[
\left(
\widehat\tau^{\mathrm{mix}}_{n,d,q}(\mathbf O_n)-\tau(P)
\right)^2
\right].
\]
Moreover,
\[
nq>0,\qquad
r_{n,d,q}
=
\min\left\{
1,\,
\frac1{nq}
+
\left[\frac{d}{nq\log(e+nq)}\right]^2
\right\}
\ge0.
\]
For each \(P\in\mathcal M_{n,d,q}\),
\cref{obj:thm:mixed-count-upper} now yields the uniform bound
\[
E_P\!\left[(\mathsf T_{\mathrm{mix}}-\tau(P))^2\right]
\le \mathfrak U_{n,d,q}
\le C_0\,r_{n,d,q}
\le\overline C\,r_{n,d,q}.
\]
The last inequality uses \(C_0\le\overline C\) and \(r_{n,d,q}\ge0\).
% lean: CausalSmith.Stat.MarRareqLogfrontier.mixed_count_upper

\item Nonnegativity of all the risks permits comparison of the minimax value with the worst-case risk of any admissible estimator. In particular,
\[
\mathfrak R_{n,d,q}
\le
\sup_{P\in\mathcal M_{n,d,q}}
E_P\!\left[(\mathsf T_{\mathrm{mix}}-\tau(P))^2\right].
\]
If \(\mathcal M_{n,d,q}\) is nonempty, the uniform bound from the preceding step gives
\[
\sup_{P\in\mathcal M_{n,d,q}}
E_P\!\left[(\mathsf T_{\mathrm{mix}}-\tau(P))^2\right]
\le\overline C\,r_{n,d,q}.
\]
If \(\mathcal M_{n,d,q}\) is empty, the real-supremum convention assigns its worst-case risk the value zero. Since \(\overline C\ge C_0>0\) and \(r_{n,d,q}\ge0\), in this case
\[
\sup_{P\in\mathcal M_{n,d,q}}
E_P\!\left[(\mathsf T_{\mathrm{mix}}-\tau(P))^2\right]
=0\le\overline C\,r_{n,d,q}.
\]
Thus, in either case,
\[
\underline c\,r_{n,d,q}
\le\mathfrak R_{n,d,q}
\le\overline C\,r_{n,d,q}.
\]
The constants were chosen independently of \(n,d,q\), as required.
% lean: CausalSmith.Stat.MarRareqLogfrontier.matched_minimax_frontier
\end{enumerate}
\end{proof}

\begin{proof}[Proof of \cref{obj:prop:fixed-d-effective-sample-reduction}]
Write
\[
N(t)=n(t)q(t)>0,
\qquad
\ell(t)=\log(e+N(t)).
\]
Define the restricted model class and its minimax risk by
\[
\mathcal M^{(0)}_{n,d,q}
:=
\left\{
P\in\mathcal M_{n,d,q}:
P(S(0)=0,S(1)=0)=1
\right\},
\qquad
\mathfrak R^{(0)}_{n,d,q}
:=
\inf_{\mathsf T\in\mathcal T_n}
\sup_{P\in\mathcal M^{(0)}_{n,d,q}}
E_P[(\mathsf T-\tau(P))^2].
\]
Here the estimator class and joint sample-and-kernel expectation are those of
\cref{obj:def:minimax-risk}.

\begin{enumerate}
\item Since \(N(t)\to\infty\), eventually
\[
N(t)\ge \max\{2,d^2\}.
\]
Also \(\ell(t)\ge1\), because \(e+N(t)\ge e\). Consequently, on this eventual range,
\[
\left[\frac{d}{N(t)\ell(t)}\right]^2
=
\frac{d^2}{N(t)^2\ell(t)^2}
\le
\frac{1}{N(t)\ell(t)^2}
\le
N(t)^{-1}.
\]
The expression in \cref{obj:def:frontier-rate} is nonnegative for every \(t\).
Eventually \(N(t)^{-1}\le1\), and both entries of its minimum are at least
\(N(t)^{-1}\). Thus
\[
N(t)^{-1}
\le
r_{n(t),d,q(t)}
=
\min\left\{
1,\,
N(t)^{-1}+
\left[\frac{d}{N(t)\ell(t)}\right]^2
\right\}
\le
2N(t)^{-1}.
\]
These inequalities prove the first order comparison.
% lean: CausalSmith.Stat.MarRareqLogfrontier.fixed_d_frontier_rate_theta

\item Let
\[
0<\underline c_{\mathrm{match}}
\le \overline C_{\mathrm{match}}<\infty
\]
be the constants supplied by \cref{obj:thm:matched-minimax-frontier}.
Its hypotheses hold at every \(t\): the sample and alphabet sizes are positive,
\(0<q(t)\le1\), and the rare-arrival restriction is assumed. Hence
\[
\underline c_{\mathrm{match}}\,r_{n(t),d,q(t)}
\le
\mathfrak R_{n(t),d,q(t)}
\le
\overline C_{\mathrm{match}}\,r_{n(t),d,q(t)}
\qquad(t\in\mathbb N).
\]
In particular, both quantities are nonnegative, and the lower inequality also gives
\[
r_{n(t),d,q(t)}
\le
\underline c_{\mathrm{match}}^{-1}
\mathfrak R_{n(t),d,q(t)}.
\]
Combining these bounds with the eventual frontier bounds from the preceding step yields
\[
\underline c_{\mathrm{match}}N(t)^{-1}
\le
\mathfrak R_{n(t),d,q(t)}
\le
2\overline C_{\mathrm{match}}N(t)^{-1}.
\]
% lean: CausalSmith.Stat.MarRareqLogfrontier.matched_minimax_frontier

\item We next obtain a pointwise upper bound for the restricted risk.
For each \(t\), use the estimator of \cref{obj:def:mixed-count-estimator}
with parameters \(n(t),d,q(t)\). Its auxiliary output is zero or a value clipped
to \([-1,1]\), so its conditional average satisfies
\[
-1
=
E_\omega[-1\mid\mathbf O_{n(t)}]
\le
\widehat\tau^{\mathrm{mix}}_{n(t),d,q(t)}(\mathbf O_{n(t)})
\le
E_\omega[1\mid\mathbf O_{n(t)}]
=
1.
\]
The average is measurable because it is a finite sum over the admissible prefix
lengths and assignments. Therefore the deterministic kernel
\[
\mathsf T^{\mathrm{mix}}_t(o,dh)
:=
\operatorname{Dirac}_{\widehat\tau^{\mathrm{mix}}_{n(t),d,q(t)}(o)}(dh)
\]
belongs to \(\mathcal T_{n(t)}\). Its joint risk equals its deterministic risk:
\[
E_P[(\mathsf T^{\mathrm{mix}}_t-\tau(P))^2]
=
E_P\!\left[
\left(
\widehat\tau^{\mathrm{mix}}_{n(t),d,q(t)}(\mathbf O_{n(t)})
-\tau(P)
\right)^2
\right].
\]
Let
\[
0<\overline C_{\mathrm{mix}}<\infty
\]
be supplied by \cref{obj:thm:mixed-count-upper}. For every
\(P\in\mathcal M^{(0)}_{n(t),d,q(t)}\), membership in the full model class
allows us to apply its risk-envelope bound:
\[
E_P[(\mathsf T^{\mathrm{mix}}_t-\tau(P))^2]
\le
\mathfrak U_{n(t),d,q(t)}
\le
\overline C_{\mathrm{mix}}\,r_{n(t),d,q(t)}.
\]
All squared risks are nonnegative, so the minimax value is bounded above by
the worst-case risk of this admissible kernel. If the restricted class is
nonempty, taking its supremum gives
\begin{equation}\label{eq:fixed-d-restricted-risk-upper}
\mathfrak R^{(0)}_{n(t),d,q(t)}
\le
\sup_{P\in\mathcal M^{(0)}_{n(t),d,q(t)}}
E_P[(\mathsf T^{\mathrm{mix}}_t-\tau(P))^2]
\le
\overline C_{\mathrm{mix}}\,r_{n(t),d,q(t)}.
\end{equation}
If the restricted class is empty, its real-valued worst-case risk is zero
under the empty-supremum convention. The same inequality follows from
\[
0\le \overline C_{\mathrm{mix}}\,r_{n(t),d,q(t)}.
\]
Thus the upper bound holds for every \(t\).
% lean: CausalSmith.Stat.MarRareqLogfrontier.mixed_count_upper
% lean: CausalSmith.Stat.MarRareqLogfrontier.fixed_d_effective_sample_reduction

\item To derive the restricted lower bound, fix \(t\). For
\(\mu_{\mathrm{test}}\in[1/4,3/4]\), define \(P_{\mu_{\mathrm{test}}}^{(t)}\) by
\[
X=0,\qquad S(0)=S(1)=Y(0)=0,
\qquad
A\sim\operatorname{Bernoulli}(1/2),\quad
Y(1)\sim\operatorname{Bernoulli}(\mu_{\mathrm{test}}),\quad
R\sim\operatorname{Bernoulli}(q(t)),
\]
where \(A,Y(1),R\) are mutually independent, and set
\(S=S(A)\), \(Y=Y(A)\). The category \(0\) exists because \(d\ge1\).
Take the observed sample independently from the induced observed-record law.

This construction satisfies randomization and consistency, and independence
of \(R\) from \((X,A,S,Y)\) implies the required conditional arrival independence.
Its occupied cells satisfy
\[
p_{a,0,0}(P_{\mu_{\mathrm{test}}}^{(t)})=\frac12,
\qquad
\rho_{a,0,0}(P_{\mu_{\mathrm{test}}}^{(t)})=q(t)
\qquad(a\in\{0,1\}),
\]
while every other cell has zero membership probability. Together with the
assumed parameter restrictions and rare-arrival slice, these facts establish
\[
P_{\mu_{\mathrm{test}}}^{(t)}\in\mathcal M^{(0)}_{n(t),d,q(t)},
\qquad
\tau(P_{\mu_{\mathrm{test}}}^{(t)})
=
E_{P_{\mu_{\mathrm{test}}}^{(t)}}[Y(1)-Y(0)]
=
\mu_{\mathrm{test}}.
\]

Define the separation and the two outcome means by
\[
\delta(t)
:=
\min\left\{\frac14,\frac{1}{4\sqrt{N(t)}}\right\},
\qquad
u_0(t):=\frac12,
\qquad
u_1(t):=\frac12+\delta(t).
\]
Since \(N(t)>0\),
\[
0<\delta(t)\le\frac14,
\qquad
u_0(t),u_1(t)\in[1/4,3/4].
\]
For \(i\in\{0,1\}\), define the single-record and sample laws by
\[
Q_i^{\mathrm{obs}}(t)
:=
\mathcal L_{P_{u_i(t)}^{(t)}}(O_1),
\qquad
Q_i^{\mathrm{sam}}(t)
:=
\mathcal L_{P_{u_i(t)}^{(t)}}(\mathbf O_{n(t)})
=
\bigl(Q_i^{\mathrm{obs}}(t)\bigr)^{\otimes n(t)}.
\]

For finite probability laws \(Q_1\ll Q_0\), use the divergence and
total-variation conventions
\[
\chi^2(Q_1\Vert Q_0)
:=
\sum_{o:\,Q_0(\{o\})>0}
\frac{\bigl(Q_1(\{o\})-Q_0(\{o\})\bigr)^2}{Q_0(\{o\})},
\qquad
\operatorname{TV}(Q_1,Q_0)
:=
\sup_{\mathcal A}|Q_1(\mathcal A)-Q_0(\mathcal A)|
=
\frac12\sum_o|Q_1(\{o\})-Q_0(\{o\})|.
\]
Under \(P_{\mu_{\mathrm{test}}}^{(t)}\), every observed record has \(X=S=0\). The remaining
coordinates, with arrived outcome \(RY=R\,Y\), have probabilities
\[
\begin{array}{c|ccccc}
(A,R,RY)
&(0,0,0)&(1,0,0)&(0,1,0)&(1,1,0)&(1,1,1)\\ \hline
\text{Probability}
&\dfrac{1-q(t)}2
&\dfrac{1-q(t)}2
&\dfrac{q(t)}2
&\dfrac{q(t)(1-\mu_{\mathrm{test}})}2
&\dfrac{q(t)\mu_{\mathrm{test}}}2
\end{array}
\]
and all other records have probability zero. At \(q(t)=1\), the first two
atoms have zero probability under both laws and are omitted from the
divergence sum. Thus \(Q_1^{\mathrm{obs}}(t)\ll Q_0^{\mathrm{obs}}(t)\), and summing
the atom probabilities gives
\[
\begin{aligned}
1+\chi^2\!\left(Q_1^{\mathrm{obs}}(t)\Vert Q_0^{\mathrm{obs}}(t)\right)
&=
1-q(t)+\frac{q(t)}2
+q(t)\bigl((1-u_1(t))^2+u_1(t)^2\bigr)\\
&=
1+2q(t)\delta(t)^2.
\end{aligned}
\]
The definition of \(\delta(t)\) gives
\[
16N(t)\delta(t)^2\le1,
\qquad
\chi^2\!\left(Q_1^{\mathrm{obs}}(t)\Vert Q_0^{\mathrm{obs}}(t)\right)
=
2q(t)\delta(t)^2
\le
\frac{1}{8n(t)}.
\]
Factoring the finite product sums and using \(1+x\le e^x\) now yields
\[
\begin{aligned}
1+\chi^2\!\left(Q_1^{\mathrm{sam}}(t)\Vert Q_0^{\mathrm{sam}}(t)\right)
&=
\left(
1+\chi^2\!\left(Q_1^{\mathrm{obs}}(t)\Vert Q_0^{\mathrm{obs}}(t)\right)
\right)^{n(t)}\\
&\le
\exp\!\left(
n(t)\chi^2\!\left(Q_1^{\mathrm{obs}}(t)\Vert Q_0^{\mathrm{obs}}(t)\right)
\right)
\le e^{1/8}.
\end{aligned}
\]
Since
\[
e^{1/2}e^{1/2}=e<3,
\qquad
e^{1/8}\le e^{1/2}<\sqrt3<2,
\]
the sample divergence is less than \(1\). Cauchy--Schwarz applied to the
finite total-variation sum therefore gives
\[
\operatorname{TV}\!\left(Q_1^{\mathrm{sam}}(t),Q_0^{\mathrm{sam}}(t)\right)
\le
\frac12
\sqrt{\chi^2\!\left(Q_1^{\mathrm{sam}}(t)\Vert Q_0^{\mathrm{sam}}(t)\right)}
\le\frac12.
\]
Consequently, every measurable sample event \(\mathcal A\) satisfies
\[
Q_1^{\mathrm{sam}}(t)(\mathcal A^c)
+
Q_0^{\mathrm{sam}}(t)(\mathcal A)
\ge
1-\operatorname{TV}\!\left(Q_1^{\mathrm{sam}}(t),Q_0^{\mathrm{sam}}(t)\right)
\ge\frac12.
\]
% lean: CausalSmith.Stat.MarRareqLogfrontier.oneCell_product_testing_half

\item Let \(\mathsf T\in\mathcal T_{n(t)}\) be arbitrary, and define its
barycenter and the relevant sample events by
\[
b_{\mathsf T}(o):=\int_{[-1,1]}h\,\mathsf T(o,dh),
\]
\[
\mathcal E_{\mathsf T,i}
:=
\left\{o:\frac{\delta(t)}2\le
|b_{\mathsf T}(o)-u_i(t)|\right\}
\quad(i\in\{0,1\}),
\qquad
\mathcal A_{\mathsf T}
:=
\left\{o:
|b_{\mathsf T}(o)-u_1(t)|
\le
|b_{\mathsf T}(o)-u_0(t)|
\right\}.
\]
These events are measurable. The triangle inequality gives, at every sample,
\[
\delta(t)
=
|u_1(t)-u_0(t)|
\le
|b_{\mathsf T}(o)-u_0(t)|
+
|b_{\mathsf T}(o)-u_1(t)|.
\]
It follows that
\[
\mathcal A_{\mathsf T}\subseteq\mathcal E_{\mathsf T,0},
\qquad
\mathcal A_{\mathsf T}^{c}\subseteq\mathcal E_{\mathsf T,1}.
\]
Applying the testing bound from the preceding step and then these inclusions,
\[
\frac12
\le
Q_1^{\mathrm{sam}}(t)(\mathcal A_{\mathsf T}^{c})
+
Q_0^{\mathrm{sam}}(t)(\mathcal A_{\mathsf T})
\le
Q_1^{\mathrm{sam}}(t)(\mathcal E_{\mathsf T,1})
+
Q_0^{\mathrm{sam}}(t)(\mathcal E_{\mathsf T,0}).
\]
Thus
\[
\frac14
\le
\max_{i\in\{0,1\}}
Q_i^{\mathrm{sam}}(t)(\mathcal E_{\mathsf T,i}).
\]
On \(\mathcal E_{\mathsf T,i}\), the barycenter's squared error is at least
\((\delta(t)/2)^2\). The barycenter inequality of
\cref{obj:lem:randomized-kernel-barycenter} therefore supplies
\[
\left(\frac{\delta(t)}2\right)^2
Q_i^{\mathrm{sam}}(t)(\mathcal E_{\mathsf T,i})
\le
\int
(b_{\mathsf T}(o)-u_i(t))^2\,Q_i^{\mathrm{sam}}(t)(do)
\le
E_{P_{u_i(t)}^{(t)}}[(\mathsf T-u_i(t))^2].
\]
Whether the probability for \(i=0\) is at most the probability for \(i=1\),
or conversely, multiplication by the nonnegative factor
\((\delta(t)/2)^2\) preserves their order. Hence
\[
\begin{aligned}
\frac{\delta(t)^2}{16}
&\le
\left(\frac{\delta(t)}2\right)^2
\max_{i\in\{0,1\}}
Q_i^{\mathrm{sam}}(t)(\mathcal E_{\mathsf T,i})\\
&=
\max_{i\in\{0,1\}}
\left\{
\left(\frac{\delta(t)}2\right)^2
Q_i^{\mathrm{sam}}(t)(\mathcal E_{\mathsf T,i})
\right\}\\
&\le
\max_{i\in\{0,1\}}
E_{P_{u_i(t)}^{(t)}}[(\mathsf T-u_i(t))^2].
\end{aligned}
\]
% lean: CausalSmith.Stat.MarRareqLogfrontier.complete_arrival_two_point_decision_reduction

\item The separation also has the required lower calibration.
If \(1/4\le1/(4\sqrt{N(t)})\), then
\[
\delta(t)^2=\frac1{16}
\ge
\frac1{16}\min\{1,N(t)^{-1}\}.
\]
In the other case,
\[
\delta(t)^2
=
\frac{1}{16N(t)}
\ge
\frac1{16}\min\{1,N(t)^{-1}\}.
\]
Combining either inequality with the preceding risk bound gives, for every
\(\mathsf T\in\mathcal T_{n(t)}\),
\[
\frac1{256}\min\{1,N(t)^{-1}\}
\le
\frac{\delta(t)^2}{16}
\le
\max_{i\in\{0,1\}}
E_{P_{u_i(t)}^{(t)}}[(\mathsf T-u_i(t))^2].
\]

The estimator class is nonempty: the constant kernel
\[
\mathsf T^{0}(o,dh):=\operatorname{Dirac}_{0}(dh)
\]
is admissible. Moreover, binary potential outcomes imply
\(\tau(P)\in[-1,1]\), and the kernel outputs lie in \([-1,1]\), so
\[
0\le E_P[(\mathsf T-\tau(P))^2]\le4
\]
for every full-data law \(P\) and admissible \(\mathsf T\).
Both testing laws belong to the full model and to its restricted subclass.
Their maximum risk therefore bounds each corresponding worst-case risk
from below; the uniform bound by \(4\) ensures these suprema are finite.
Taking the infimum over estimators gives
\[
\frac1{256}\min\{1,N(t)^{-1}\}
\le
\mathfrak R_{n(t),d,q(t)},
\qquad
\frac1{256}\min\{1,N(t)^{-1}\}
\le
\mathfrak R^{(0)}_{n(t),d,q(t)}.
\]
% lean: CausalSmith.Stat.MarRareqLogfrontier.oneCell_minimax_risk_floors

\item The last lower bound establishes nonnegativity of the restricted risk
for every \(t\). Its upper bound in \cref{eq:fixed-d-restricted-risk-upper}, combined with the eventual
frontier bound, gives
\[
0\le
\mathfrak R^{(0)}_{n(t),d,q(t)}
\le
2\overline C_{\mathrm{mix}}N(t)^{-1}
\qquad\text{eventually}.
\]
Also \(N(t)\to\infty\) implies \(N(t)\ge1\) eventually, so
\[
\min\{1,N(t)^{-1}\}=N(t)^{-1}.
\]
The lower bound consequently becomes
\[
\frac1{256}N(t)^{-1}
\le
\mathfrak R^{(0)}_{n(t),d,q(t)},
\qquad\text{equivalently}\qquad
N(t)^{-1}\le256\,\mathfrak R^{(0)}_{n(t),d,q(t)}.
\]
Together these inequalities prove the restricted-risk order comparison,
and hence all three assertions.
% lean: CausalSmith.Stat.MarRareqLogfrontier.fixed_d_effective_sample_reduction
\end{enumerate}
\end{proof}

\begin{proof}[Proof of \cref{obj:prop:phase-boundaries}]
\begin{enumerate}
\item\label{proof:phase-boundaries:comparison} For every \(t\in\mathbb N\), write
\[
\ell_t:=\log(e+N_t),\qquad
r_t:=r_{n_t,d_t,q_t}
=\min\left\{1,N_t^{-1}
+\left(\frac{d_t}{N_t\ell_t}\right)^2\right\},
\qquad
\mathfrak R_t:=\mathfrak R_{n_t,d_t,q_t}.
\]
The expression for \(r_t\) is supplied by \cref{obj:def:frontier-rate}. The size and arrival assumptions give
\[
N_t=n_tq_t>0,\qquad
N_t^{-1}+\left(\frac{d_t}{N_t\ell_t}\right)^2\ge0,
\qquad r_t\ge0.
\]
All hypotheses of \cref{obj:thm:matched-minimax-frontier} hold at every \(t\). Consequently, there are constants
\(0<\underline c\le\overline C<\infty\), independent of \(t\), such that
\[
\underline c\,r_t\le\mathfrak R_t\le\overline C\,r_t.
\]
In particular,
\[
0\le r_t\le\underline c^{-1}\mathfrak R_t,
\qquad
0\le\mathfrak R_t\le\overline C\,r_t.
\]
The first inequality and the squeeze theorem show that
\(\mathfrak R_t\to0\) implies \(r_t\to0\); the second gives the converse. Thus
\[
\mathfrak R_t\longrightarrow0
\quad\Longleftrightarrow\quad
r_t\longrightarrow0.
\]
% lean: CausalSmith.Stat.MarRareqLogfrontier.phase_boundaries

\item\label{proof:phase-boundaries:convergence} Define, for every \(t\in\mathbb N\),
\[
\nu_t:=N_t^{-1},
\qquad
\xi_t:=\frac{d_t}{N_t\ell_t}.
\]
Since \(e>1\) and \(N_t>0\),
\[
e+N_t>1,\qquad
\ell_t>0,\qquad
N_t\ell_t>0,\qquad
\nu_t\ge0,\qquad \xi_t\ge0.
\]
Effective sample growth gives
\[
\nu_t\longrightarrow0.
\]
Continuity of the square root, together with
\[
\sqrt{\xi_t^2}=|\xi_t|=\xi_t,
\]
shows that \(\xi_t^2\to0\) implies \(\xi_t\to0\). Continuity of squaring gives the reverse implication. Hence
\[
\xi_t^2\longrightarrow0
\quad\Longleftrightarrow\quad
\xi_t\longrightarrow0.
\]
Because \(N_t\ell_t>0\) for every \(t\), the quotient characterization of little \(o\) gives
\[
d_t=o(N_t\ell_t)
\quad\Longleftrightarrow\quad
\xi_t\longrightarrow0.
\]

We next establish the truncation equivalence
\[
r_t=\min\{1,\nu_t+\xi_t^2\}\longrightarrow0
\quad\Longleftrightarrow\quad
\xi_t^2\longrightarrow0.
\]
For the forward direction, \(r_t\to0\) implies \(r_t<1\) eventually. At each such index, \(\nu_t+\xi_t^2\ge1\) would give \(r_t=1\). Therefore, eventually,
\[
\nu_t+\xi_t^2<1,\qquad
r_t=\nu_t+\xi_t^2.
\]
The sum consequently tends to zero, and
\[
0\le\xi_t^2\le\nu_t+\xi_t^2
\]
implies \(\xi_t^2\to0\). Conversely, if \(\xi_t^2\to0\), then
\[
\nu_t+\xi_t^2\longrightarrow0,
\qquad
\min\{1,\nu_t+\xi_t^2\}\longrightarrow\min\{1,0\}=0,
\]
by continuity of addition and minimum. Combining these equivalences with \cref{proof:phase-boundaries:comparison} yields
\[
\mathfrak R_t\longrightarrow0
\quad\Longleftrightarrow\quad
d_t=o(N_t\ell_t)
=o\!\left(N_t\log(e+N_t)\right).
\]
% lean: CausalSmith.Stat.MarRareqLogfrontier.phase_min_rate_tendsto
% lean: CausalSmith.Stat.MarRareqLogfrontier.phase_frontier_rate_tendsto

\item\label{proof:phase-boundaries:theta-comparison} The comparison in \cref{proof:phase-boundaries:comparison} also gives a two-sided asymptotic comparison. Indeed, nonnegativity allows us to write, for every \(t\),
\[
|\mathfrak R_t|
\le\overline C\,|r_t|,
\qquad
|r_t|
\le\underline c^{-1}|\mathfrak R_t|.
\]
These are the defining bounds for
\[
\mathfrak R_t=O(r_t),
\qquad
r_t=O(\mathfrak R_t),
\qquad\text{and hence}\qquad
\mathfrak R_t=\Theta(r_t).
\]
% lean: CausalSmith.Stat.MarRareqLogfrontier.phase_boundaries

\item\label{proof:phase-boundaries:baseline} Since \(N_t\to\infty\), eventually \(N_t\ge1\). At those indices,
\[
0\le\nu_t\le1,
\qquad
\nu_t\le\nu_t+\xi_t^2.
\]
Thus both entries in the minimum defining \(r_t\) are at least \(\nu_t\), and
\[
0\le\nu_t\le\min\{1,\nu_t+\xi_t^2\}=r_t.
\]
This proves the baseline bound
\[
\nu_t=O(r_t).
\]
% lean: CausalSmith.Stat.MarRareqLogfrontier.phase_rate_lower_bigO

\item\label{proof:phase-boundaries:dimension-necessity} Define the dimension scale, for every \(t\in\mathbb N\), by
\[
B_{\mathrm{phase},t}:=\sqrt{N_t}\,\ell_t.
\]
The positivity established in \cref{proof:phase-boundaries:convergence} gives
\[
B_{\mathrm{phase},t}>0,
\qquad
N_t\ell_t>0.
\]
Moreover, using \((\sqrt{N_t})^2=N_t\),
\[
\nu_t(N_t\ell_t)^2
=N_t^{-1}N_t^2\ell_t^2
=N_t\ell_t^2
=B_{\mathrm{phase},t}^2.
\]

Suppose first that \(r_t=O(\nu_t)\). Choose a real constant
\(\beta_{\mathrm{rate}}\) such that, eventually,
\[
0\le r_t\le\beta_{\mathrm{rate}}\nu_t.
\]
Since \(\beta_{\mathrm{rate}}\nu_t\to0\), the squeeze theorem gives \(r_t\to0\), and hence \(r_t<1\) eventually. As in \cref{proof:phase-boundaries:convergence}, this forces
\[
\nu_t+\xi_t^2<1,
\qquad
r_t=\nu_t+\xi_t^2
\]
eventually. Intersecting these eventual conditions with the assumed bound gives
\[
\nu_t+\xi_t^2\le\beta_{\mathrm{rate}}\nu_t,
\qquad
\xi_t^2\le\beta_{\mathrm{rate}}\nu_t,
\]
where the latter inequality uses \(\nu_t\ge0\). Multiplication by the positive quantity \((N_t\ell_t)^2\) then yields
\[
d_t^2
=\xi_t^2(N_t\ell_t)^2
\le\beta_{\mathrm{rate}}\nu_t(N_t\ell_t)^2
=\beta_{\mathrm{rate}}B_{\mathrm{phase},t}^2.
\]
Set
\[
\beta_{\mathrm{bound}}:=\max\{1,\beta_{\mathrm{rate}}\}\in\mathbb R.
\]
Then
\[
\beta_{\mathrm{bound}}\ge1,\qquad
\beta_{\mathrm{rate}}\le\beta_{\mathrm{bound}}
\le\beta_{\mathrm{bound}}^2.
\]
Consequently, eventually,
\[
d_t^2
\le\beta_{\mathrm{bound}}^2B_{\mathrm{phase},t}^2.
\]
Both \(d_t\) and \(\beta_{\mathrm{bound}}B_{\mathrm{phase},t}\) are nonnegative, so this squared inequality gives
\[
d_t\le\beta_{\mathrm{bound}}B_{\mathrm{phase},t}.
\]
Therefore
\[
r_t=O(\nu_t)
\quad\Longrightarrow\quad
d_t=O(B_{\mathrm{phase},t}).
\]
% lean: CausalSmith.Stat.MarRareqLogfrontier.phase_dimension_rate_bigO

\item\label{proof:phase-boundaries:dimension-sufficiency} Conversely, suppose that \(d_t=O(B_{\mathrm{phase},t})\). Choose a real constant
\(\beta_{\mathrm{dim}}\) such that, eventually,
\[
0\le d_t\le\beta_{\mathrm{dim}}B_{\mathrm{phase},t}.
\]
The right-hand side is consequently nonnegative, so squaring gives
\[
d_t^2\le\beta_{\mathrm{dim}}^2B_{\mathrm{phase},t}^2.
\]
Dividing by \((N_t\ell_t)^2>0\) and using the identity from \cref{proof:phase-boundaries:dimension-necessity} yields
\[
\xi_t^2
=\frac{d_t^2}{(N_t\ell_t)^2}
\le
\beta_{\mathrm{dim}}^2
\frac{B_{\mathrm{phase},t}^2}{(N_t\ell_t)^2}
=\beta_{\mathrm{dim}}^2\nu_t.
\]
It follows that, eventually,
\[
0\le r_t
=\min\{1,\nu_t+\xi_t^2\}
\le\nu_t+\xi_t^2
\le(1+\beta_{\mathrm{dim}}^2)\nu_t.
\]
Thus \(r_t=O(\nu_t)\). Together with \cref{proof:phase-boundaries:dimension-necessity}, this establishes
\[
r_t=O(\nu_t)
\quad\Longleftrightarrow\quad
d_t=O(B_{\mathrm{phase},t})
=O\!\left(\sqrt{N_t}\log(e+N_t)\right).
\]
% lean: CausalSmith.Stat.MarRareqLogfrontier.phase_dimension_rate_bigO

\item If \(\mathfrak R_t=\Theta(\nu_t)\), then
\[
r_t=O(\mathfrak R_t),
\qquad
\mathfrak R_t=O(\nu_t)
\quad\Longrightarrow\quad
r_t=O(\nu_t),
\]
by transitivity of big \(O\). The equivalence in \cref{proof:phase-boundaries:dimension-sufficiency} therefore gives
\(d_t=O(B_{\mathrm{phase},t})\).

Conversely, if \(d_t=O(B_{\mathrm{phase},t})\), then \cref{proof:phase-boundaries:dimension-sufficiency} gives \(r_t=O(\nu_t)\). Combining this with \cref{proof:phase-boundaries:theta-comparison,proof:phase-boundaries:baseline} gives
\[
\mathfrak R_t=O(r_t),\quad r_t=O(\nu_t)
\quad\Longrightarrow\quad
\mathfrak R_t=O(\nu_t),
\]
and
\[
\nu_t=O(r_t),\quad r_t=O(\mathfrak R_t)
\quad\Longrightarrow\quad
\nu_t=O(\mathfrak R_t).
\]
These two bounds prove
\[
\mathfrak R_t=\Theta(N_t^{-1})
\quad\Longleftrightarrow\quad
d_t=O\!\left(\sqrt{N_t}\log(e+N_t)\right),
\]
which is the second claimed equivalence.
% lean: CausalSmith.Stat.MarRareqLogfrontier.phase_boundaries
\end{enumerate}
\end{proof}

\begin{proof}[Proof of \cref{obj:thm:unrestricted-fixed-q-minimax-frontier}]
Fix \(0<q<1/2\). For \(n,d\ge1\), put
\[
\mathfrak h_{n,d}^{\circ}
:=
\frac1n+\left(\frac{d}{n\log(e+n)}\right)^2,
\qquad
\mathfrak h_{n,d}:=\min\{1,\mathfrak h_{n,d}^{\circ}\}.
\]
For a bounded estimator kernel \(\mathsf T\) and a full-data law \(P\), write
\[
\mathcal R_P(\mathsf T)
:=
\int\!\int_{[-1,1]}
(h-\tau(P))^2\,\mathsf T(o,dh)\,
\mathcal L_P(\mathbf O_n)(do).
\]
Thus \(\mathcal R_P(\mathsf T)\) integrates both the sample and the estimator draw.

\begin{enumerate}
\item
The deterministic estimator in \cref{obj:def:mixed-count-estimator} defines an admissible decision kernel
\[
\mathsf T_{\mathrm{mix}}(o,dh)
:=
\delta_{\widehat\tau^{\mathrm{mix}}_{n,d,q}(o)}(dh).
\]
By \cref{obj:thm:unrestricted-mixed-count-upper}, there is a universal constant \(\overline C_{\mathrm{mix}}>0\) such that
\[
0\le \mathcal R_P(\mathsf T_{\mathrm{mix}})
\le \overline C_{\mathrm{mix}}\,r_{n,d,q}
\qquad
\text{for every }P\in\mathcal M^+_{n,d,q}.
\]
Taking the supremum over the model and then using this kernel in the defining infimum gives
\[
\mathfrak R^+_{n,d,q}
\le
\sup_{P\in\mathcal M^+_{n,d,q}}
\mathcal R_P(\mathsf T_{\mathrm{mix}})
\le
\overline C_{\mathrm{mix}}\,r_{n,d,q}.
\]
If the model class is empty, its real-valued worst-case risk is taken to be zero, and the same inequality follows from
\[
r_{n,d,q}
=
\min\left\{1,\frac1{nq}
+\left(\frac{d}{nq\log(e+nq)}\right)^2\right\}
\ge0.
\]
% lean: CausalSmith.Stat.MarRareqLogfrontier.unrestricted_fixed_q_frontier_upper

\item
We next convert this upper bound to the fixed-\(q\) rate. Define
\[
\beta_q:=\frac{1-\log q}{q},
\qquad
N:=nq,
\qquad
\ell:=\log(e+N),
\qquad
\ell_{n,1}:=\log(e+n).
\]
Throughout the present domain,
\[
N>0,\qquad \ell\ge1,\qquad \ell_{n,1}>0,
\qquad \log q\le0,\qquad \beta_q\ge1.
\]
Since
\[
e+n\le\frac{e+N}{q},
\]
monotonicity of the logarithm yields
\[
\ell_{n,1}\le\ell-\log q
\le(1-\log q)\ell.
\]
The second inequality follows from
\[
(-\log q)(\ell-1)\ge0.
\]
Consequently,
\[
n\ell_{n,1}
\le n(1-\log q)\ell
=\beta_q N\ell.
\]
Moreover,
\[
q\beta_q=1-\log q\ge1,
\qquad
q\beta_q^2=(1-\log q)\beta_q\ge1,
\]
so
\[
\frac1N\le\frac{\beta_q^2}{n}.
\]
The positive denominator comparison also gives
\[
0\le\frac{d}{N\ell}
\le\beta_q\frac{d}{n\ell_{n,1}},
\qquad
\left(\frac{d}{N\ell}\right)^2
\le\beta_q^2
\left(\frac{d}{n\ell_{n,1}}\right)^2.
\]
Adding these bounds,
\[
\frac1N+\left(\frac{d}{N\ell}\right)^2
\le\beta_q^2\mathfrak h_{n,d}^{\circ}.
\]
Truncation therefore gives
\[
r_{n,d,q}
\le\min\{1,\beta_q^2\mathfrak h_{n,d}^{\circ}\}
\le\beta_q^2\mathfrak h_{n,d}.
\]
Indeed, if \(\mathfrak h_{n,d}^{\circ}\le1\), the last inequality is
\[
\min\{1,\beta_q^2\mathfrak h_{n,d}^{\circ}\}
\le\beta_q^2\mathfrak h_{n,d}^{\circ}
=\beta_q^2\mathfrak h_{n,d};
\]
if \(\mathfrak h_{n,d}^{\circ}\ge1\), it is
\[
\min\{1,\beta_q^2\mathfrak h_{n,d}^{\circ}\}
\le1\le\beta_q^2
=\beta_q^2\mathfrak h_{n,d}.
\]
Define the positive fixed-\(q\) constant
\[
\overline C_{\mathrm{poly}}(q)
:=\overline C_{\mathrm{mix}}\beta_q^2.
\]
The upper bound from the preceding step now implies
\[
\mathfrak R^+_{n,d,q}
\le\overline C_{\mathrm{mix}}r_{n,d,q}
\le\overline C_{\mathrm{poly}}(q)\mathfrak h_{n,d}.
\]
% lean: CausalSmith.Stat.MarRareqLogfrontier.frontierRate_le_fixedQRate_mul; CausalSmith.Stat.MarRareqLogfrontier.unrestricted_fixed_q_polynomial_upper

\item
For the lower bound, invoke the published fixed-positivity lower bound on the control-zero observational subclass
\citep[Theorem~2, Section~3.2, p.~9, with its proof in Section~C.5, pp.~39--44]{ZengBalakrishnanHanKennedy2026}.
Its truncated, all-\(n,d\) comparison is supplied by
\cref{obj:thm:zeng-fixed-positivity-polynomial-frontier}: at \(\epsilon=q\), there is a constant \(\underline c_Z(q)>0\) such that
\[
\underline c_Z(q)\mathfrak h_{n,d}
\le\mathfrak R^Z_{n,d,q}
\qquad(n,d\ge1).
\]

We transfer this lower bound to the unrestricted experiment. For an observational estimator kernel \(\widehat\theta\), define
\[
\mathcal R^Z_{P_Z}(\widehat\theta)
:=
\int\!\int_{[-1,1]}
(h-\theta_Z(P_Z))^2\,
\widehat\theta(\mathbf z,dh)\,
\mathcal L_{P_Z}(\mathbf Z_n)(d\mathbf z).
\]
Given any admissible unrestricted kernel \(\mathsf T\), define its observational pullback by
\[
\widehat\theta_{\mathsf T}(\mathbf z,\mathcal V)
:=
2^{-n}
\sum_{\mathbf u\in\{0,1\}^n}
\mathsf T\!\left(
\bigl(o(x_i,b_i,y_i;u_i)\bigr)_{i=1}^n,\mathcal V
\right),
\]
where
\[
\mathbf z=((x_i,b_i,y_i))_{i=1}^n,
\qquad
\mathbf u=(u_i)_{i=1}^n,
\qquad
\mathcal V\in\mathscr B([-1,1]),
\]
and \(o(x,b,y;u)\) is the synthetic observation map of
\cref{obj:lem:synthetic-randomization-kernel}. This finite mixture composes the parameter-independent fair-coin transformation with \(\mathsf T\), and hence is an admissible observational kernel.

For every \(P_Z\in\mathcal D_{d,q}\),
\cref{obj:lem:synthetic-randomization-kernel} supplies a completion
\(P_Z^*\in\mathcal M^+_{n,d,q}\) whose observed product law is exactly the transformed sample law and whose target satisfies
\[
\tau(P_Z^*)=\theta_Z(P_Z).
\]
Thus
\[
\mathcal R^Z_{P_Z}(\widehat\theta_{\mathsf T})
=
\mathcal R_{P_Z^*}(\mathsf T).
\]
Both targets and all kernel outputs lie in \([-1,1]\), giving
\[
0\le\mathcal R^Z_{P_Z}(\widehat\theta_{\mathsf T})\le4,
\qquad
0\le\mathcal R_P(\mathsf T)\le4.
\]
The supremum comparison and the observational minimax definition therefore give, for every \(\mathsf T\),
\[
\mathfrak R^Z_{n,d,q}
\le
\sup_{P_Z\in\mathcal D_{d,q}}
\mathcal R^Z_{P_Z}(\widehat\theta_{\mathsf T})
\le
\sup_{P\in\mathcal M^+_{n,d,q}}
\mathcal R_P(\mathsf T).
\]
Taking the infimum over \(\mathsf T\) yields
\[
\mathfrak R^Z_{n,d,q}\le\mathfrak R^+_{n,d,q}.
\]

Now define
\[
\underline c(q):=\frac{\underline c_Z(q)}{\beta_q^2}>0.
\]
Using the rate comparison established in the preceding step,
\[
\underline c(q)r_{n,d,q}
\le
\frac{\underline c_Z(q)}{\beta_q^2}
\,\beta_q^2\mathfrak h_{n,d}
=
\underline c_Z(q)\mathfrak h_{n,d}
\le
\mathfrak R^Z_{n,d,q}
\le
\mathfrak R^+_{n,d,q}.
\]
% lean: CausalSmith.Stat.MarRareqLogfrontier.zeng_fixed_positivity_polynomial_frontier; CausalSmith.Stat.MarRareqLogfrontier.zengMinimaxRisk_le_unrestrictedMinimaxRisk; CausalSmith.Stat.MarRareqLogfrontier.unrestricted_fixed_q_minimax_frontier

\item
Choose the common upper constant
\[
\overline C(q)
:=
\max\{\overline C_{\mathrm{mix}},
       \overline C_{\mathrm{poly}}(q)\}
+\underline c(q).
\]
These are finite real constants depending only on \(q\), and
\[
0<\underline c(q)\le\overline C(q).
\]
The rates are nonnegative:
\[
\mathfrak h_{n,d}\ge0,
\qquad r_{n,d,q}\ge0,
\]
and the chosen constant satisfies
\[
\overline C_{\mathrm{mix}}\le\overline C(q),
\qquad
\overline C_{\mathrm{poly}}(q)\le\overline C(q).
\]

To obtain the fixed-rate lower bound with the same constant, observe that
\[
0<N\le n,
\qquad
0<\ell\le\ell_{n,1},
\qquad
0<N\ell\le n\ell_{n,1}.
\]
It follows that
\[
\frac1n\le\frac1N,
\qquad
0\le\frac{d}{n\ell_{n,1}}\le\frac{d}{N\ell},
\]
and hence
\[
\left(\frac{d}{n\ell_{n,1}}\right)^2
\le
\left(\frac{d}{N\ell}\right)^2.
\]
Adding and truncating at one proves
\[
\mathfrak h_{n,d}
=
\min\left\{1,\frac1n+
\left(\frac{d}{n\ell_{n,1}}\right)^2\right\}
\le
\min\left\{1,\frac1N+
\left(\frac{d}{N\ell}\right)^2\right\}
=r_{n,d,q}.
\]
Combining this comparison with the preceding bounds gives
\[
\begin{aligned}
\underline c(q)\mathfrak h_{n,d}
&\le \underline c(q)r_{n,d,q}
 \le\mathfrak R^+_{n,d,q},\\
\mathfrak R^+_{n,d,q}
&\le\overline C_{\mathrm{poly}}(q)\mathfrak h_{n,d}
 \le\overline C(q)\mathfrak h_{n,d},\\
\underline c(q)r_{n,d,q}
&\le\mathfrak R^+_{n,d,q},\\
\mathfrak R^+_{n,d,q}
&\le\overline C_{\mathrm{mix}}r_{n,d,q}
 \le\overline C(q)r_{n,d,q}.
\end{aligned}
\]
These are the two asserted comparisons.
% lean: CausalSmith.Stat.MarRareqLogfrontier.fixedQRate_le_frontierRate; CausalSmith.Stat.MarRareqLogfrontier.unrestricted_fixed_q_minimax_frontier
\end{enumerate}
\end{proof}

\begin{proof}[Proof of \cref{obj:prop:unrestricted-fixed-q-phase-boundaries}]
Fix \(q\), the arrival floor, with \(0<q<1/2\). Throughout the proof, \(n\) ranges over positive integers, and all limits are taken as \(n\to\infty\).

\begin{enumerate}
\item For a sequence of integers \(d_n\ge1\), define
\[
\ell_n^{\mathrm{fix}}:=\log(e+n),
\qquad
r_n^{\mathrm{fix}}
:=
\min\left\{1,\frac1n+
\left(\frac{d_n}{n\ell_n^{\mathrm{fix}}}\right)^2\right\}.
\]
Apply \cref{obj:thm:unrestricted-fixed-q-minimax-frontier}, whose lower comparison uses the published zero-control-regression lower bound
\citep[Theorem~2, Section~3.2, p.~9, and its proof in Section~C.5, pp.~39--44]{ZengBalakrishnanHanKennedy2026}.
It supplies constants \(0<\underline c(q)\le\overline C(q)\), depending only on \(q\), such that
\[
\underline c(q)\,r_n^{\mathrm{fix}}
\le \mathfrak R^{+}_{n,d_n,q}
\le \overline C(q)\,r_n^{\mathrm{fix}}
\qquad(n\ge1).
\]
These constants are common to every dimension sequence. Moreover,
\[
0\le r_n^{\mathrm{fix}},
\qquad
0\le\overline C(q).
\]
% lean: CausalSmith.Stat.MarRareqLogfrontier.unrestricted_fixed_q_minimax_frontier

\item We first establish the consistency equivalence. The comparison above gives
\[
0\le r_n^{\mathrm{fix}}
\le \underline c(q)^{-1}\mathfrak R^{+}_{n,d_n,q},
\qquad
0\le\mathfrak R^{+}_{n,d_n,q}
\le\overline C(q)\,r_n^{\mathrm{fix}}.
\]
Applying the squeeze theorem in each direction yields
\[
\mathfrak R^{+}_{n,d_n,q}\longrightarrow0
\quad\Longleftrightarrow\quad
r_n^{\mathrm{fix}}\longrightarrow0.
\]

To characterize the latter limit, define
\[
\nu_n:=\frac1n,
\qquad
\xi_n:=\frac{d_n}{n\ell_n^{\mathrm{fix}}}.
\]
Since \(e+n>1\), we have
\[
\ell_n^{\mathrm{fix}}>0,
\qquad
n\ell_n^{\mathrm{fix}}>0,
\qquad
\xi_n\ge0,
\qquad
\nu_n\ge0,
\qquad
\nu_n\longrightarrow0.
\]
The last limit follows from \(n\to\infty\) and continuity of inversion at infinity. Continuity of the square root and nonnegativity of \(\xi_n\) give
\[
\xi_n^2\longrightarrow0
\quad\Longrightarrow\quad
\xi_n=\sqrt{\xi_n^2}\longrightarrow0.
\]
Conversely, continuity of squaring gives
\[
\xi_n\longrightarrow0
\quad\Longrightarrow\quad
\xi_n^2\longrightarrow0.
\]
The positive denominator also gives the usual quotient characterization of little-\(o\):
\[
d_n=o\!\left(n\ell_n^{\mathrm{fix}}\right)
\quad\Longleftrightarrow\quad
\xi_n\longrightarrow0.
\]

It remains to handle the truncation. If
\(r_n^{\mathrm{fix}}\to0\), then eventually
\[
r_n^{\mathrm{fix}}=\min\{1,\nu_n+\xi_n^2\}<1.
\]
At every such index,
\[
\nu_n+\xi_n^2<1,
\qquad
r_n^{\mathrm{fix}}=\nu_n+\xi_n^2:
\]
indeed, \(\nu_n+\xi_n^2\ge1\) would make the minimum equal to \(1\).
Thus \(\nu_n+\xi_n^2\to0\), and
\[
0\le\xi_n^2\le\nu_n+\xi_n^2
\]
implies \(\xi_n^2\to0\). Conversely, if \(\xi_n^2\to0\), then
\[
\nu_n+\xi_n^2\longrightarrow0,
\qquad
\min\{1,\nu_n+\xi_n^2\}\longrightarrow\min\{1,0\}=0,
\]
by continuity of addition and the minimum.
% lean: CausalSmith.Stat.MarRareqLogfrontier.phase_min_rate_tendsto

Consequently,
\[
r_n^{\mathrm{fix}}\longrightarrow0
\quad\Longleftrightarrow\quad
\xi_n^2\longrightarrow0
\quad\Longleftrightarrow\quad
\xi_n\longrightarrow0
\quad\Longleftrightarrow\quad
d_n=o\!\left(n\log(e+n)\right).
\]
Together with the risk comparison, this proves the first assertion.
% lean: CausalSmith.Stat.MarRareqLogfrontier.fixedQRate_tendsto_iff_dimension_littleO

\item For the sufficient condition for the parametric rate, choose
\[
c:=\underline c(q)>0.
\]
Given \(M>0\), define
\[
C:=\overline C(q)(1+M^2).
\]
These constants satisfy
\[
c=\underline c(q)
\le\overline C(q)
\le\overline C(q)(1+M^2)=C.
\]
Suppose that eventually
\[
d_n\le M\sqrt n\,\ell_n^{\mathrm{fix}}.
\]
Both sides are nonnegative, so squaring this inequality and using
\((\sqrt n)^2=n\) gives
\[
d_n^2
\le\left(M\sqrt n\,\ell_n^{\mathrm{fix}}\right)^2
=M^2n\left(\ell_n^{\mathrm{fix}}\right)^2.
\]
Dividing by the positive squared denominator yields
\[
\xi_n^2
=\frac{d_n^2}{(n\ell_n^{\mathrm{fix}})^2}
\le\frac{M^2}{n}.
\]
Here the cancellation uses
\[
\frac1n(n\ell_n^{\mathrm{fix}})^2
=n\left(\ell_n^{\mathrm{fix}}\right)^2.
\]
Since \(1/n\le1\) for \(n\ge1\), both entries of the minimum defining the rate are at least \(1/n\). Hence
\[
\frac1n
\le r_n^{\mathrm{fix}}
=\min\left\{1,\frac1n+\xi_n^2\right\}
\le\frac1n+\xi_n^2
\le\frac{1+M^2}{n}.
\]
Multiplying the lower and upper bounds by the corresponding frontier constants gives, eventually,
\[
\frac{c}{n}
\le\underline c(q)\,r_n^{\mathrm{fix}}
\le\mathfrak R^{+}_{n,d_n,q}
\le\overline C(q)\,r_n^{\mathrm{fix}}
\le\frac{C}{n}.
\]
The choice of \(c\) depends only on \(q\), and the choice of \(C\) depends only on \(q\) and \(M\). The eventual dimension bound determines the threshold along each sequence.
% lean: CausalSmith.Stat.MarRareqLogfrontier.unrestricted_fixed_q_phase_boundaries

\item For the converse, suppose there exist constants \(0<c_*\le C_*\) such that eventually
\[
\frac{c_*}{n}
\le\mathfrak R^{+}_{n,d_n,q}
\le\frac{C_*}{n}.
\]
In particular, \(C_*\ge0\). Combining the upper risk bound with the lower frontier comparison gives
\[
\underline c(q)\,r_n^{\mathrm{fix}}
\le\mathfrak R^{+}_{n,d_n,q}
\le\frac{C_*}{n},
\]
and therefore
\[
0\le r_n^{\mathrm{fix}}
\le\frac{C_*}{\underline c(q)}\,\frac1n
\longrightarrow0.
\]
Thus eventually \(r_n^{\mathrm{fix}}<1\).

Define the dimension-bound constant
\[
K_{\mathrm{dim}}
:=\max\left\{1,\frac{C_*}{\underline c(q)}\right\}.
\]
Its calibration gives
\[
0\le K_{\mathrm{dim}},
\qquad
1\le K_{\mathrm{dim}},
\qquad
\frac{C_*}{\underline c(q)}
\le K_{\mathrm{dim}}
\le K_{\mathrm{dim}}^2.
\]
At all sufficiently large indices, \(r_n^{\mathrm{fix}}<1\) implies
\[
\frac1n+\xi_n^2<1,
\qquad
r_n^{\mathrm{fix}}=\frac1n+\xi_n^2,
\]
because a sum at least \(1\) would make the truncated rate equal to \(1\).
The preceding rate bound and \(1/n\ge0\) now yield
\[
\xi_n^2
\le\frac1n+\xi_n^2
=r_n^{\mathrm{fix}}
\le\frac{C_*}{\underline c(q)}\,\frac1n.
\]
Multiplying by \((n\ell_n^{\mathrm{fix}})^2>0\) gives
\[
d_n^2
\le\frac{C_*}{\underline c(q)}
\,n\left(\ell_n^{\mathrm{fix}}\right)^2
\le K_{\mathrm{dim}}^2
\,n\left(\ell_n^{\mathrm{fix}}\right)^2.
\]
The second inequality uses
\(n(\ell_n^{\mathrm{fix}})^2\ge0\). Also,
\[
K_{\mathrm{dim}}\sqrt n\,\ell_n^{\mathrm{fix}}\ge0,
\qquad
\left(K_{\mathrm{dim}}\sqrt n\,\ell_n^{\mathrm{fix}}\right)^2
=K_{\mathrm{dim}}^2n\left(\ell_n^{\mathrm{fix}}\right)^2.
\]
Since \(d_n\ge0\), comparison of these squares implies
\[
d_n\le K_{\mathrm{dim}}\sqrt n\,\ell_n^{\mathrm{fix}}
\qquad\text{eventually}.
\]
Both the dimension and the comparison scale are nonnegative, so this is precisely the required bound
\[
d_n=O\!\left(\sqrt n\,\log(e+n)\right).
\]
% lean: CausalSmith.Stat.MarRareqLogfrontier.unrestricted_fixed_q_phase_boundaries
\end{enumerate}
\end{proof}

\begin{proof}[Proof of \cref{obj:thm:unrestricted-mixed-count-upper}]
\begin{enumerate}
\item By \cref{obj:thm:unrestricted-uniform-mixed-count-certificate}, choose a universal constant \(0<\overline C<\infty\). For every \(n,d,q,P\) satisfying the stated domain, model, and sampling conditions, this constant gives
\[
E_P\!\left[
 \bigl(\widehat\tau^{\mathrm{mix}}_{n,d,q}(\mathbf O_n)-\Phi(P)\bigr)^2
\right]
\le \mathfrak U_{n,d,q}
\le \overline C\,r_{n,d,q},
\]
where \(\Phi(P)\) is the signed cell functional of
\cref{obj:def:cell-functional}. The envelope here has precisely the branch formulas in the statement.
% lean: CausalSmith.Stat.MarRareqLogfrontier.unrestricted_uniform_mixed_count_certificate

\item The model hypothesis permits application of
\cref{obj:lem:unrestricted-cell-identification}, which supplies
\[
\tau(P)=\Phi(P).
\]
Substituting this identity into the preceding risk bound yields
\[
E_P\!\left[
 \bigl(\widehat\tau^{\mathrm{mix}}_{n,d,q}(\mathbf O_n)-\tau(P)\bigr)^2
\right]
\le \mathfrak U_{n,d,q}
\le \overline C\,r_{n,d,q}.
\]
% lean: CausalSmith.Stat.MarRareqLogfrontier.unrestricted_cell_identification

\item To establish the auxiliary-risk comparison, fix an observed sample
\[
\mathbf o=(o_i)_{i=1}^n
\in
\bigl(\{0,\ldots,d-1\}\times\{0,1\}^{4}\bigr)^n,
\qquad
o_i=(X_i,A_i,S_i,R_i,R_iY_i).
\]
Define the prefix probabilities and remaining probability by
\[
w_{k_0}:=P_\omega(K_0=k_0)
=e^{-n/2}\frac{(n/2)^{k_0}}{k_0!}
\quad(k_0\in\mathbb Z_{\ge0}),
\qquad
\delta_{\mathrm{tail}}
:=1-\sum_{k_0=0}^{n}w_{k_0}
=P_\omega(K_0>n).
\]
The disjoint prefix events imply
\[
w_{k_0}\ge0,\qquad
\sum_{k_0=0}^{n}w_{k_0}\le1,\qquad
\delta_{\mathrm{tail}}\ge0.
\]
For each \(0\le k_0\le n\), define the assignment set and its output average by
\[
\Sigma_{k_0}
:=\bigl\{\sigma:\{1,\ldots,k_0\}\to\{1,2,3\}\bigr\},
\qquad
\overline\tau_{k_0}(\mathbf o)
:=3^{-k_0}\sum_{\sigma\in\Sigma_{k_0}}
 \widetilde\tau(\mathbf o;(k_0,\sigma)),
\]
where the three assignment values denote the membership, pilot, and estimation streams. The notation \((k_0,\sigma)\) specifies the auxiliary prefix and assignment in \cref{obj:def:mixed-count-estimator}.

The weights \(w_0,\ldots,w_n,\delta_{\mathrm{tail}}\) sum to one, and
\[
\sum_{k_0=0}^{n}w_{k_0}\overline\tau_{k_0}(\mathbf o)-\tau(P)
=
\delta_{\mathrm{tail}}\bigl(-\tau(P)\bigr)
+\sum_{k_0=0}^{n}w_{k_0}
 \bigl(\overline\tau_{k_0}(\mathbf o)-\tau(P)\bigr).
\]
Convexity of the square therefore gives
\[
\left(
 \sum_{k_0=0}^{n}w_{k_0}\overline\tau_{k_0}(\mathbf o)-\tau(P)
\right)^2
\le
\delta_{\mathrm{tail}}\tau(P)^2
+\sum_{k_0=0}^{n}w_{k_0}
 \bigl(\overline\tau_{k_0}(\mathbf o)-\tau(P)\bigr)^2.
\]
% lean: CausalSmith.Stat.MarRareqLogfrontier.finite_weighted_square_with_remainder

\item Since \(|\Sigma_{k_0}|=3^{k_0}\), centering the assignment average gives
\[
\overline\tau_{k_0}(\mathbf o)-\tau(P)
=
3^{-k_0}\sum_{\sigma\in\Sigma_{k_0}}
 \bigl(\widetilde\tau(\mathbf o;(k_0,\sigma))-\tau(P)\bigr).
\]
The finite mean-square inequality, obtained by Cauchy--Schwarz, consequently yields
\[
\bigl(\overline\tau_{k_0}(\mathbf o)-\tau(P)\bigr)^2
\le
3^{-k_0}\sum_{\sigma\in\Sigma_{k_0}}
 \bigl(\widetilde\tau(\mathbf o;(k_0,\sigma))-\tau(P)\bigr)^2.
\]
This also covers \(k_0=0\): the empty prefix has one assignment, so the averaging denominator equals one.
% lean: CausalSmith.Stat.MarRareqLogfrontier.finite_uniform_square

\item Multiplying these inequalities by the nonnegative prefix probabilities and summing gives
\[
\begin{aligned}
&\sum_{k_0=0}^{n}w_{k_0}
 \bigl(\overline\tau_{k_0}(\mathbf o)-\tau(P)\bigr)^2\\
&\qquad\le
\sum_{k_0=0}^{n}w_{k_0}3^{-k_0}
 \sum_{\sigma\in\Sigma_{k_0}}
 \bigl(\widetilde\tau(\mathbf o;(k_0,\sigma))-\tau(P)\bigr)^2.
\end{aligned}
\]
The estimator's finite averaging formula in
\cref{obj:def:mixed-count-estimator} is
\[
\widehat\tau^{\mathrm{mix}}_{n,d,q}(\mathbf o)
=
\sum_{k_0=0}^{n}w_{k_0}\overline\tau_{k_0}(\mathbf o).
\]
Combining this identity with the preceding two bounds establishes, for every \(\mathbf o\),
\[
\begin{aligned}
&\bigl(\widehat\tau^{\mathrm{mix}}_{n,d,q}(\mathbf o)-\tau(P)\bigr)^2\\
&\quad\le
\sum_{k_0=0}^{n}w_{k_0}3^{-k_0}
 \sum_{\sigma\in\Sigma_{k_0}}
 \bigl(\widetilde\tau(\mathbf o;(k_0,\sigma))-\tau(P)\bigr)^2
+\delta_{\mathrm{tail}}\tau(P)^2\\
&\quad=
E_\omega\!\left[
 \bigl(\widetilde\tau(\mathbf O_n;\omega)-\tau(P)\bigr)^2
 \,\middle|\,\mathbf O_n=\mathbf o
\right].
\end{aligned}
\]
The last term in the sum is the squared loss of the zero output on \(K_0>n\). All functions in this inequality are measurable and integrable on the finite observed-sample space. Integrating under the canonical product law yields
\[
E_P\!\left[
 \bigl(\widehat\tau^{\mathrm{mix}}_{n,d,q}(\mathbf O_n)-\tau(P)\bigr)^2
\right]
\le
E_P\!\left[
 E_\omega\!\left[
  \bigl(\widetilde\tau(\mathbf O_n;\omega)-\tau(P)\bigr)^2
  \,\middle|\,\mathbf O_n
 \right]
\right],
\]
as required.
% lean: CausalSmith.Stat.MarRareqLogfrontier.mixed_count_derandomization
\end{enumerate}
\end{proof}

\begin{proof}[Proof of \cref{obj:thm:connected-interval-frontier}]
Fix \(\alpha\in(0,1)\). Throughout the proof, write
\[
\beta:=1-\alpha>0,\qquad
N:=nq>0,\qquad
\ell:=\log(e+N)\ge1,
\]
where \(n,d\ge1\), \(0<q\le1\), and \(q\ell^2\le1/64\).
Define the two length scales by
\[
s_{\mathrm{par}}:=\min\{1,N^{-1/2}\},
\qquad
s_{\mathrm{many}}:=\min\left\{1,\frac{d}{N\ell}\right\}.
\]

\begin{enumerate}
\item \label{proof:connected-interval-upper}\textbf{Upper bound for the risk-envelope interval.}
By \cref{obj:thm:mixed-count-upper}, there is a universal
\(\overline C>0\) such that, for every \(P\in\mathcal M_{n,d,q}\),
\[
E_P\!\left[
  \bigl(\widehat\tau^{\mathrm{mix}}_{n,d,q}-\tau(P)\bigr)^2
\right]
\le \mathfrak U_{n,d,q}
\le \overline C\,r_{n,d,q}.
\]
For every observed sample, define the envelope radius by
\[
\rho_{\mathrm{env}}
:=\sqrt{\frac{\mathfrak U_{n,d,q}}{\alpha}}.
\]
The endpoints in \cref{obj:def:risk-envelope-interval} satisfy
\[
\operatorname{length}(I^{\mathrm{mix}}_\alpha)
=
\min\{1,\widehat\tau^{\mathrm{mix}}_{n,d,q}+\rho_{\mathrm{env}}\}
-
\max\{-1,\widehat\tau^{\mathrm{mix}}_{n,d,q}-\rho_{\mathrm{env}}\}
\le 2\rho_{\mathrm{env}}.
\]
Consequently, with
\[
C_\alpha:=2\sqrt{\frac{\overline C}{\alpha}}>0,
\]
we have
\[
E_P[\operatorname{length}(I^{\mathrm{mix}}_\alpha)]
\le
2\sqrt{\frac{\mathfrak U_{n,d,q}}{\alpha}}
\le C_\alpha\sqrt{r_{n,d,q}}.
\]
Taking the supremum gives
\[
\sup_{P\in\mathcal M_{n,d,q}}
E_P[\operatorname{length}(I^{\mathrm{mix}}_\alpha)]
\le C_\alpha\sqrt{r_{n,d,q}}.
\]
For an empty model class, the supremum is zero by the convention in
\cref{obj:def:interval-length-risk}, and the same inequality follows from
\(r_{n,d,q}\ge0\).
% lean: CausalSmith.Stat.MarRareqLogfrontier.riskEnvelope_interval_worst_upper

\item \label{proof:connected-interval-calibration}\textbf{Calibration of the activated regime.}
Define constants depending only on \(\alpha\) by
\[
\eta:=e^{-32},\qquad
M_{\mathrm{act},\alpha}:=\left\lceil\frac8\beta\right\rceil+1,
\qquad
D_{\mathrm{act},\alpha}
:=
\left\lceil
\frac{18432M_{\mathrm{act},\alpha}^2}{\beta}
\right\rceil+1,
\]
\[
T_{\mathrm{act},\alpha}
:=
\max\left\{
1,e^2,4\eta,2\eta D_{\mathrm{act},\alpha},
\exp\left(\frac{32-\log(\beta/16)}{28}\right)
\right\},
\qquad
c_{\mathrm{act},\alpha}
:=\frac{21\beta}{32}\frac{\eta}{24}>0.
\]
The ceiling inequalities give
\[
M_{\mathrm{act},\alpha}\ge1,\qquad
\frac1{M_{\mathrm{act},\alpha}}\le\frac\beta8,\qquad
D_{\mathrm{act},\alpha}\ge1,\qquad
\frac{18432M_{\mathrm{act},\alpha}^2}{\beta}
\le D_{\mathrm{act},\alpha}.
\]
Use the quantities from \cref{obj:def:activation-lower-handle}:
\[
K:=\lceil\ell\rceil,\qquad H:=K^2,\qquad
b_0:=\frac{\eta}{N\ell},\qquad
J:=\min\left\{d-1,\left\lfloor\frac1{2b_0}\right\rfloor\right\}.
\]
They are defined for all the parameters under consideration, including
\(d=1\), when \(J=0\).

Suppose first that \(N\ge T_{\mathrm{act},\alpha}\). Since \(\ell\ge1\),
\[
D_{\mathrm{act},\alpha}
\le \frac{N}{2\eta}
\le \frac{N\ell}{2\eta}
=\frac1{2b_0}.
\]
Also,
\[
b_0\le\frac{\eta}{N}\le\frac14,\qquad
\ell\ge2,\qquad K\ge2.
\]
The ceiling bound \(K\le\ell+1\le2\ell\), together with the slice restriction,
gives
\[
qH=qK^2\le4q\ell^2\le\frac1{16}.
\]
By \cref{obj:lem:moment-matched-reciprocal-priors}, choose finitely supported
probability laws \(\Pi_0,\Pi_1\) on \([1,H]\) with matching moments through
degree \(K\). Interchange their names if necessary and define
\[
\mu_{\mathrm{rec},i}:=\int z^{-1}\,d\Pi_i(z)
\quad(i=0,1).
\]
Then
\[
\int z^v\,d\Pi_0(z)=\int z^v\,d\Pi_1(z)
\quad(0\le v\le K),
\qquad
\mu_{\mathrm{rec},1}-\mu_{\mathrm{rec},0}\ge\frac1{12}.
\]

We will show that whenever \(J\ge D_{\mathrm{act},\alpha}\),
\[
c_{\mathrm{act},\alpha}s_{\mathrm{many}}
\le\mathfrak L_{n,d,q,\alpha}.
\]
% lean: CausalSmith.Stat.MarRareqLogfrontier.interval_activated_large_regime

\item \textbf{The activated grid and its likelihood comparison.}
Continue under
\(N\ge T_{\mathrm{act},\alpha}\) and \(J\ge D_{\mathrm{act},\alpha}\).
For each integer \(0\le\iota\le M_{\mathrm{act},\alpha}\), define
\[
\Pi^{(\iota)}
:=
\left(1-\frac{\iota}{M_{\mathrm{act},\alpha}}\right)\Pi_0
+\frac{\iota}{M_{\mathrm{act},\alpha}}\Pi_1.
\]
These are finitely supported probability laws on \([1,H]\), with the same
moments through degree \(K\).

For completeness, the full-data laws underlying the activated experiment
can be specified as follows. For a fixed
\(\mathbf z=(z_0,\ldots,z_{d-1})\in[1,H]^d\), let
\(P^{\mathrm{act}}_{\mathbf z}\) have baseline probabilities
\[
P^{\mathrm{act}}_{\mathbf z}(X=x)
=
\begin{cases}
b_0,&0\le x<J,\\
1-Jb_0,&x=J,\\
0,&J<x<d.
\end{cases}
\]
Let treatment be an independent fair coin, and set
\[
S(0)=S(1)=Y(0)=0,\qquad S=0,\qquad Y=AY(1).
\]
Conditional on \(X=x\), generate the treated potential outcome and arrival
independently, with
\[
P^{\mathrm{act}}_{\mathbf z}(Y(1)=1\mid X=x)
=
\begin{cases}
z_x^{-1},&x<J,\\
0,&x\ge J,
\end{cases}
\qquad
P^{\mathrm{act}}_{\mathbf z}(R=1\mid X=x,A)
=
\begin{cases}
qz_x,&x<J,\\
qH,&x\ge J.
\end{cases}
\]
The inequalities \(Jb_0\le1/2\) and \(qH\le1/16\) ensure valid
probabilities. Independence of treatment, consistency, conditional
independence of arrival and outcome, and the arrival floor follow from
this construction. Thus
\[
P^{\mathrm{act}}_{\mathbf z}\in\mathcal M_{n,d,q},
\qquad
\tau(P^{\mathrm{act}}_{\mathbf z})
=b_0\sum_{x=0}^{J-1}z_x^{-1}.
\]

Let \(\widetilde{\mathbb Q}_{\mathrm{act},\iota}\) be the augmented
sample law obtained from the rule in
\cref{obj:def:activation-lower-handle}, using
\((\Pi^{(\iota)})^{\otimes d}\) for the latent vector. Let
\(\mathbb Q_{\mathrm{act},\iota}\) be its projection onto the observed
sample. The conditional probabilities in that construction give
\[
\mathbb Q_{\mathrm{act},\iota}
=
\int
\mathcal L_{P^{\mathrm{act}}_{\mathbf z}}(\mathbf O_n)\,
d(\Pi^{(\iota)})^{\otimes d}(\mathbf z).
\]
Define the activation count in each rare category and the common cutoff
event by
\[
V_x:=\sum_{i=1}^n
\mathbf1\{X_i=x,\ \mathsf B_{\mathrm{act},i}=1\}
\quad(0\le x<J),
\qquad
G_{\mathrm{act}}:=\{V_x\le K\text{ for every }x<J\}.
\]
Conditional on the membership, treatment, and activation pattern, every
activated observation in category \(x\) contributes an affine factor in
\(z_x\), and every inactive observation contributes a constant factor.
The conditional likelihood therefore has degree at most \(V_x\) in each
\(z_x\). On \(G_{\mathrm{act}}\), integration against the product priors
and moment matching imply equality of all augmented sample atoms:
\[
\widetilde{\mathbb Q}_{\mathrm{act},\iota}(\{\mathbf o_{\mathrm{aug}}\})
=
\widetilde{\mathbb Q}_{\mathrm{act},0}(\{\mathbf o_{\mathrm{aug}}\})
\qquad(\mathbf o_{\mathrm{aug}}\in G_{\mathrm{act}}).
\]
Hence the two laws also assign the same mass to \(G_{\mathrm{act}}\)
and its complement. Splitting any event into its intersections with
these two sets, and then projecting, yields
\[
\operatorname{TV}
\bigl(\mathbb Q_{\mathrm{act},0},
      \mathbb Q_{\mathrm{act},\iota}\bigr)
\le
\widetilde{\mathbb Q}_{\mathrm{act},0}(G_{\mathrm{act}}^c).
\]

To bound this last probability, define
\[
K_{\mathrm{tail}}:=K+1,\qquad
\lambda_{\mathrm{act}}:=nb_0qH=\frac{\eta K^2}{\ell}.
\]
The membership and activation pattern has a common law under all the
priors, and each \(V_x\) has distribution
\(\operatorname{Binomial}(n,b_0qH)\). Counting subsets of
\(K_{\mathrm{tail}}\) successful observations gives
\[
E\binom{V_x}{K_{\mathrm{tail}}}
=
\binom{n}{K_{\mathrm{tail}}}(b_0qH)^{K_{\mathrm{tail}}},
\]
with both sides zero when \(K_{\mathrm{tail}}>n\). Since
\(\mathbf1\{V_x>K\}\le\binom{V_x}{K_{\mathrm{tail}}}\),
\[
\Pr(V_x>K)
\le\frac{\lambda_{\mathrm{act}}^{K_{\mathrm{tail}}}}
         {K_{\mathrm{tail}}!}.
\]
The ceiling inequalities imply
\[
\lambda_{\mathrm{act}}
=\frac{\eta K^2}{\ell}
\le2\eta K_{\mathrm{tail}}
\le e^{-31}K_{\mathrm{tail}},
\qquad
K_{\mathrm{tail}}\ge\ell.
\]
Using the \(K_{\mathrm{tail}}\)-th term of the exponential series,
\[
\frac{\lambda_{\mathrm{act}}^{K_{\mathrm{tail}}}}
     {K_{\mathrm{tail}}!}
\le
e^{-31K_{\mathrm{tail}}}
\frac{K_{\mathrm{tail}}^{K_{\mathrm{tail}}}}
     {K_{\mathrm{tail}}!}
\le e^{-30K_{\mathrm{tail}}}
\le e^{-30\ell}.
\]
A union bound, followed by the mass cap and
\(N\le e^\ell\), \(\ell\le e^\ell\), now gives
\[
\widetilde{\mathbb Q}_{\mathrm{act},0}(G_{\mathrm{act}}^c)
\le Je^{-30\ell}
\le \frac{N\ell}{2e^{-32}}e^{-30\ell}
\le e^{32-28\ell}.
\]
Finally, the definition of \(T_{\mathrm{act},\alpha}\) implies
\[
\ell\ge\frac{32-\log(\beta/16)}{28},
\qquad
e^{32-28\ell}\le\frac\beta{16}.
\]
Thus, for every grid index,
\[
\operatorname{TV}
\bigl(\mathbb Q_{\mathrm{act},0},
      \mathbb Q_{\mathrm{act},\iota}\bigr)
\le\frac\beta{16}.
\]
% lean: CausalSmith.Stat.MarRareqLogfrontier.activatedAugmentedMixture_interval_minimax_calibrated

\item \textbf{Coverage of the prior neighborhoods and their length cost.}
Define the target separation, centers, and neighborhood radius by
\[
\Delta_{\mathrm{act}}
:=Jb_0(\mu_{\mathrm{rec},1}-\mu_{\mathrm{rec},0})
\ge\frac{Jb_0}{12}>0,
\]
\[
\theta_{\mathrm{act},\iota}
:=
Jb_0\left[
\left(1-\frac{\iota}{M_{\mathrm{act},\alpha}}\right)
\mu_{\mathrm{rec},0}
+\frac{\iota}{M_{\mathrm{act},\alpha}}\mu_{\mathrm{rec},1}
\right]
=
\theta_{\mathrm{act},0}
+\frac{\iota\Delta_{\mathrm{act}}}{M_{\mathrm{act},\alpha}},
\qquad
w_{\mathrm{act}}
:=\frac{\Delta_{\mathrm{act}}}{8M_{\mathrm{act},\alpha}}.
\]
Under the \(\iota\)-th product prior, the actual target has mean
\(\theta_{\mathrm{act},\iota}\). Independence of the latent coordinates
and \(z_x^{-1}\in[0,1]\) give
\[
\operatorname{Var}\!\left(
b_0\sum_{x=0}^{J-1}z_x^{-1}
\right)\le\frac{Jb_0^2}{4}.
\]
Indeed, a variable in \([0,1]\) has variance at most \(1/4\), and the
variances of these independent summands add. Chebyshev's inequality
therefore gives
\[
\Pr_{\Pi^{(\iota)\otimes d}}\!\left(
\left|\tau(P^{\mathrm{act}}_{\mathbf z})
-\theta_{\mathrm{act},\iota}\right|>w_{\mathrm{act}}
\right)
\le
\frac{Jb_0^2/4}{w_{\mathrm{act}}^2}
\le
\frac{2304M_{\mathrm{act},\alpha}^2}{J}
\le\frac\beta8.
\]

Fix any \(\mathsf I\in\mathcal C_{n,d,q,\alpha}\), and let
\(\nu_{\mathrm{act},\iota}\) denote its returned-interval law when the
sample has law \(\mathbb Q_{\mathrm{act},\iota}\):
\[
\nu_{\mathrm{act},\iota}
:=\mathbb Q_{\mathrm{act},\iota}\mathsf I.
\]
For each grid index, define the interval event
\[
F_{\mathrm{act},\iota}
:=
\left\{
[l,u]:
[l,u]\cap
[\theta_{\mathrm{act},\iota}-w_{\mathrm{act}},
 \theta_{\mathrm{act},\iota}+w_{\mathrm{act}}]
\ne\varnothing
\right\}.
\]
Honesty holds for every fixed completion. Averaging that coverage over
the finite prior and subtracting the target-concentration failure
probability yields
\[
\nu_{\mathrm{act},\iota}(F_{\mathrm{act},\iota})
\ge\beta-\frac\beta8=\frac{7\beta}{8}.
\]
By \cref{obj:lem:randomized-kernel-tv-comparison} and the preceding
likelihood comparison,
\[
\operatorname{TV}
(\nu_{\mathrm{act},0},\nu_{\mathrm{act},\iota})
\le\frac\beta{16},
\qquad
\nu_{\mathrm{act},0}(F_{\mathrm{act},\iota})
\ge\frac{13\beta}{16}.
\]

Define the number of neighborhoods met by a returned interval by
\[
\mathcal N_{\mathrm{hit}}([l,u])
:=
\sum_{\iota=0}^{M_{\mathrm{act},\alpha}}
\mathbf1_{F_{\mathrm{act},\iota}}([l,u]).
\]
Meeting the \(\iota\)-th neighborhood is equivalent to
\(\theta_{\mathrm{act},\iota}\in[l-w_{\mathrm{act}},u+w_{\mathrm{act}}]\).
The centers are spaced by
\(\Delta_{\mathrm{act}}/M_{\mathrm{act},\alpha}\). Thus the distance
between the extreme centers contained in this expanded interval gives
\[
(u-l)+2w_{\mathrm{act}}
\ge
\frac{\Delta_{\mathrm{act}}}{M_{\mathrm{act},\alpha}}
\bigl(\mathcal N_{\mathrm{hit}}([l,u])-1\bigr).
\]
For zero or one contained center, the right side is nonpositive or
zero, so the inequality covers those cases as well. Taking expectations,
\[
\begin{aligned}
E_{\nu_{\mathrm{act},0}}[u-l]
&\ge
\frac{\Delta_{\mathrm{act}}}{M_{\mathrm{act},\alpha}}
\left[
(M_{\mathrm{act},\alpha}+1)\frac{13\beta}{16}-1
\right]
-2w_{\mathrm{act}}\\
&=
\Delta_{\mathrm{act}}
\left[
\frac{13\beta}{16}
+\frac{13\beta}{16M_{\mathrm{act},\alpha}}
-\frac{5}{4M_{\mathrm{act},\alpha}}
\right]\\
&\ge\frac{21\beta}{32}\Delta_{\mathrm{act}},
\end{aligned}
\]
where the last step uses \(1/M_{\mathrm{act},\alpha}\le\beta/8\).
The left side is a prior average of the expected lengths of this same
procedure over laws in \(\mathcal M_{n,d,q}\). Therefore
\[
\frac{21\beta}{32}\Delta_{\mathrm{act}}
\le
\sup_{P\in\mathcal M_{n,d,q}}E_P[\operatorname{length}(\mathsf I)].
\]
The constant procedure returning \([-1,1]\) is honest, so the procedure
class is nonempty. Taking its defining infimum gives
\[
\frac{21\beta}{32}\Delta_{\mathrm{act}}
\le\mathfrak L_{n,d,q,\alpha}.
\]

To express the separation in terms of \(d\), split according to the
minimum defining \(J\). If
\(d-1\le\lfloor(2b_0)^{-1}\rfloor\), then \(J=d-1\), and
\(J\ge D_{\mathrm{act},\alpha}\ge1\) implies \(d\ge2\) and
\(J\ge d/2\). Otherwise,
\(J=\lfloor(2b_0)^{-1}\rfloor\), and
\[
Jb_0\ge\frac12-b_0\ge\frac14.
\]
In both cases,
\[
Jb_0\ge\min\left\{\frac{db_0}{2},\frac14\right\}.
\]
Since \(\eta\le1/2\), the two inequalities
\[
\frac{\eta}{24}s_{\mathrm{many}}
\le\frac{db_0}{24},
\qquad
\frac{\eta}{24}s_{\mathrm{many}}
\le\frac1{48}
\]
give
\[
\frac{\eta}{24}s_{\mathrm{many}}
\le
\frac1{12}\min\left\{\frac{db_0}{2},\frac14\right\}
\le\frac{Jb_0}{12}
\le\Delta_{\mathrm{act}}.
\]
Consequently,
\[
c_{\mathrm{act},\alpha}s_{\mathrm{many}}
\le\mathfrak L_{n,d,q,\alpha}
\qquad
(N\ge T_{\mathrm{act},\alpha},\ J\ge D_{\mathrm{act},\alpha}).
\]
% lean: CausalSmith.Stat.MarRareqLogfrontier.interval_activated_large_regime

\item \textbf{An explicit one-cell floor.}
Use the family \(P_u\) supplied by
\cref{obj:lem:one-cell-testing-family}. We record the calibration needed
here:
\[
c_{\mathrm{par},\alpha}:=\frac{3\beta^2}{512}>0,
\qquad
c_{\mathrm{par},\alpha}s_{\mathrm{par}}
\le\mathfrak L_{n,d,q,\alpha}.
\]
To obtain this explicit constant, define
\[
\eta_{\mathrm{cell}}:=\frac\beta8,\qquad
w_{\mathrm{cell}}
:=\min\left\{\frac14,\frac{\eta_{\mathrm{cell}}}{4\sqrt N}\right\}.
\]
Choose an integer \(M_{\mathrm{cell},\alpha}\ge8/\beta\), and set
\[
\delta_{\mathrm{cell}}
:=\frac{w_{\mathrm{cell}}}{M_{\mathrm{cell},\alpha}}>0,
\qquad
u_{\mathrm{cell},\iota}
:=\frac12+\iota\delta_{\mathrm{cell}}
\quad(0\le\iota\le M_{\mathrm{cell},\alpha}).
\]
These grid points are distinct and belong to \([1/4,3/4]\), and
\[
2Nw_{\mathrm{cell}}^2\le\frac{\eta_{\mathrm{cell}}^2}{8},
\qquad
w_{\mathrm{cell}}\ge
\frac{\eta_{\mathrm{cell}}}{4}s_{\mathrm{par}}.
\]
Indeed, \(N>0\) and \(0<\eta_{\mathrm{cell}}=\beta/8<1\), so
\(w_{\mathrm{cell}}>0\). The defining minimum gives
\[
w_{\mathrm{cell}}\le\frac{\eta_{\mathrm{cell}}}{4\sqrt N},
\qquad
2Nw_{\mathrm{cell}}^2
\le 2N\left(\frac{\eta_{\mathrm{cell}}}{4\sqrt N}\right)^2
=\frac{\eta_{\mathrm{cell}}^2}{8},
\]
where the equality uses \((\sqrt N)^2=N\). For the lower bound, if
\(1/4\le\eta_{\mathrm{cell}}/(4\sqrt N)\), then
\(w_{\mathrm{cell}}=1/4\), and \(s_{\mathrm{par}}\le1\) gives
\[
\frac{\eta_{\mathrm{cell}}}{4}s_{\mathrm{par}}
\le\frac{\eta_{\mathrm{cell}}}{4}
\le\frac14=w_{\mathrm{cell}}.
\]
Otherwise, \(w_{\mathrm{cell}}=\eta_{\mathrm{cell}}/(4\sqrt N)\),
and \(s_{\mathrm{par}}\le N^{-1/2}\) gives
\[
\frac{\eta_{\mathrm{cell}}}{4}s_{\mathrm{par}}
\le\frac{\eta_{\mathrm{cell}}}{4\sqrt N}
=w_{\mathrm{cell}}.
\]

Write the one-record and sample laws of this family as
\[
\mathbb Q_{\mathrm{cell},u}^{(1)}:=\mathcal L_{P_u}(O_1),
\qquad
\mathbb Q_{\mathrm{cell},u}^{(n)}
:=(\mathbb Q_{\mathrm{cell},u}^{(1)})^{\otimes n}.
\]
In this family, the arrival floor and \(P_u(R=1)=q\), together with
balanced assignment, force arrival probability \(q\) in each treatment
arm. Conditional independence of arrival and outcome then gives the
two treated arrived-outcome atom probabilities \(qu/2\) and
\(q(1-u)/2\); all other observed atom probabilities are constant in \(u\).
Consequently, defining the one-record chi-squared divergence by its
finite sum,
\[
\mathcal X_{\mathrm{cell},\iota}^{(1)}
:=
\sum_{\substack{o:\\
\mathbb Q_{\mathrm{cell},1/2}^{(1)}(\{o\})>0}}
\frac{
\left[
\mathbb Q_{\mathrm{cell},u_{\mathrm{cell},\iota}}^{(1)}(\{o\})
-\mathbb Q_{\mathrm{cell},1/2}^{(1)}(\{o\})
\right]^2}
{\mathbb Q_{\mathrm{cell},1/2}^{(1)}(\{o\})}
=
2q(\iota\delta_{\mathrm{cell}})^2.
\]
This computation also covers \(q=1\), since the vanished nonarrival
atoms are omitted from the sum. Define
\(\mathcal X_{\mathrm{cell},\iota}^{(n)}\) by the analogous sum for the
sample laws, and put
\[
t_{\mathrm{cell},\iota}
:=2N(\iota\delta_{\mathrm{cell}})^2
\le\frac{\eta_{\mathrm{cell}}^2}{8}.
\]
Independence of the observations gives
\[
1+\mathcal X_{\mathrm{cell},\iota}^{(n)}
=
\bigl(1+\mathcal X_{\mathrm{cell},\iota}^{(1)}\bigr)^n
\le e^{t_{\mathrm{cell},\iota}}.
\]
For \(0\le t_{\mathrm{cell},\iota}<1\), comparison of the exponential
and geometric series gives
\[
e^{t_{\mathrm{cell},\iota}}
\le\frac1{1-t_{\mathrm{cell},\iota}},
\qquad
\frac{t_{\mathrm{cell},\iota}}{1-t_{\mathrm{cell},\iota}}
\le\eta_{\mathrm{cell}}^2,
\]
using \(t_{\mathrm{cell},\iota}\le\eta_{\mathrm{cell}}^2/8\) and
\(\eta_{\mathrm{cell}}\le1\). Thus Cauchy--Schwarz applied to the
finite likelihood-ratio sum yields
\[
\operatorname{TV}
\bigl(\mathbb Q_{\mathrm{cell},1/2}^{(n)},
      \mathbb Q_{\mathrm{cell},u_{\mathrm{cell},\iota}}^{(n)}\bigr)
\le\frac12\sqrt{\mathcal X_{\mathrm{cell},\iota}^{(n)}}
\le\eta_{\mathrm{cell}}.
\]

For any honest procedure \(\mathsf I\), coverage of each grid point is
at least \(\beta\) under its corresponding sample law. By
\cref{obj:lem:randomized-kernel-tv-comparison}, under the reference law
\(P_{1/2}\) each point is contained in the returned interval with
probability at least
\(\beta-\eta_{\mathrm{cell}}=7\beta/8\). A connected interval
containing a given number of grid points has length at least their
spacing times that number minus one. Summing the inclusion
probabilities therefore gives
\[
\begin{aligned}
E_{P_{1/2}}[\operatorname{length}(\mathsf I)]
&\ge
\delta_{\mathrm{cell}}
\left[(M_{\mathrm{cell},\alpha}+1)\frac{7\beta}{8}-1\right]\\
&\ge\frac{3\beta}{4}w_{\mathrm{cell}},
\end{aligned}
\]
because \(1/M_{\mathrm{cell},\alpha}\le\beta/8\).
All interval risks lie in \([0,2]\), and the constant interval
\([-1,1]\) is honest. Hence taking the supremum over models and the
infimum over honest procedures preserves this lower bound. Finally,
\[
c_{\mathrm{par},\alpha}s_{\mathrm{par}}
\le
\frac{3\beta}{4}\frac{\eta_{\mathrm{cell}}}{4}s_{\mathrm{par}}
\le
\frac{3\beta}{4}w_{\mathrm{cell}}
\le\mathfrak L_{n,d,q,\alpha},
\]
which proves the asserted explicit floor.
% lean: CausalSmith.Stat.MarRareqLogfrontier.oneCell_interval_minimax_floor

\item \textbf{Extension of the many-cell scale to all regimes.}
Define
\[
c_{\mathrm{many},\alpha}
:=
\min\left\{
c_{\mathrm{act},\alpha},
\ c_{\mathrm{par},\alpha}
       \min\{1,T_{\mathrm{act},\alpha}^{-1/2}\},
\ \frac{c_{\mathrm{par},\alpha}}{D_{\mathrm{act},\alpha}}
\right\}>0.
\]
We claim
\[
c_{\mathrm{many},\alpha}s_{\mathrm{many}}
\le\mathfrak L_{n,d,q,\alpha}
\]
throughout the parameter range.

First suppose \(N\ge T_{\mathrm{act},\alpha}\). If
\(J\ge D_{\mathrm{act},\alpha}\), the activated bound gives
\[
c_{\mathrm{many},\alpha}s_{\mathrm{many}}
\le c_{\mathrm{act},\alpha}s_{\mathrm{many}}
\le\mathfrak L_{n,d,q,\alpha}.
\]
If \(J<D_{\mathrm{act},\alpha}\), the capacity inequality from \cref{proof:connected-interval-calibration}
implies
\[
\left\lfloor\frac1{2b_0}\right\rfloor
\ge D_{\mathrm{act},\alpha}.
\]
Thus the minimum defining \(J\) forces
\(d-1<D_{\mathrm{act},\alpha}\), and hence
\(d\le D_{\mathrm{act},\alpha}\). Since \(N\ge1\) and \(\ell\ge1\),
\[
s_{\mathrm{many}}
\le\frac{d}{N\ell}
\le\frac{D_{\mathrm{act},\alpha}}{\sqrt N}
=
D_{\mathrm{act},\alpha}s_{\mathrm{par}}.
\]
It follows that
\[
c_{\mathrm{many},\alpha}s_{\mathrm{many}}
\le
c_{\mathrm{many},\alpha}D_{\mathrm{act},\alpha}s_{\mathrm{par}}
\le c_{\mathrm{par},\alpha}s_{\mathrm{par}}
\le\mathfrak L_{n,d,q,\alpha}.
\]
This case includes \(d=1\) and \(J=0\).

Finally, if \(N<T_{\mathrm{act},\alpha}\), monotonicity of the square
root and reciprocal gives
\[
\min\{1,T_{\mathrm{act},\alpha}^{-1/2}\}
\le s_{\mathrm{par}}.
\]
Using \(s_{\mathrm{many}}\le1\), we obtain
\[
c_{\mathrm{many},\alpha}s_{\mathrm{many}}
\le c_{\mathrm{many},\alpha}
\le
c_{\mathrm{par},\alpha}
\min\{1,T_{\mathrm{act},\alpha}^{-1/2}\}
\le\mathfrak L_{n,d,q,\alpha}.
\]
% lean: CausalSmith.Stat.MarRareqLogfrontier.interval_many_scale_of_large_regime

\item \label{proof:connected-interval-lower-scales}\textbf{Combination of the lower scales.}
The identities
\[
\sqrt{\min\{1,N^{-1}\}}=s_{\mathrm{par}},
\qquad
\sqrt{\min\left\{1,\left(\frac{d}{N\ell}\right)^2\right\}}
=s_{\mathrm{many}}
\]
follow by splitting at \(N^{-1}=1\) and \(d/(N\ell)=1\), respectively.
Define
\[
c_{\mathrm{both},\alpha}
:=\min\{c_{\mathrm{par},\alpha},c_{\mathrm{many},\alpha}\}>0.
\]
The preceding bounds give
\[
c_{\mathrm{both},\alpha}
\max\{s_{\mathrm{par}},s_{\mathrm{many}}\}
\le\mathfrak L_{n,d,q,\alpha}.
\]
Moreover,
\[
r_{n,d,q}
\le2\max\{s_{\mathrm{par}}^2,s_{\mathrm{many}}^2\}.
\]
Indeed, when both untruncated terms \(N^{-1}\) and
\((d/(N\ell))^2\) are at most one, their sum is at most twice their
maximum. If either exceeds one, the maximum of the truncated terms
equals one, while \(r_{n,d,q}\le1\). Since the maximum is nonnegative,
\[
\sqrt{r_{n,d,q}}
\le2\max\{s_{\mathrm{par}},s_{\mathrm{many}}\}.
\]
Therefore, defining
\[
c_{\mathrm{low},\alpha}:=\frac{c_{\mathrm{both},\alpha}}2>0,
\]
we obtain
\[
c_{\mathrm{low},\alpha}\sqrt{r_{n,d,q}}
\le\mathfrak L_{n,d,q,\alpha}.
\]
% lean: CausalSmith.Stat.MarRareqLogfrontier.frontier_interval_lower_of_two_scales

\item \textbf{Honesty and the minimax comparison.}
Set
\[
\underline c_\alpha
:=\min\{c_{\mathrm{low},\alpha},C_\alpha\},
\qquad
\overline C_\alpha:=C_\alpha.
\]
Then \(0<\underline c_\alpha\le\overline C_\alpha<\infty\), and
\cref{proof:connected-interval-lower-scales} supplies the required lower inequality.

Fix \(P\in\mathcal M_{n,d,q}\). Define the squared error by
\[
V_{\mathrm{env}}
:=
\bigl(\widehat\tau^{\mathrm{mix}}_{n,d,q}(\mathbf O_n)-\tau(P)\bigr)^2.
\]
The risk bound implies
\[
0\le E_P[V_{\mathrm{env}}]\le\mathfrak U_{n,d,q}.
\]
If \(\mathfrak U_{n,d,q}=0\), nonnegativity gives
\(V_{\mathrm{env}}=0\) almost surely, so the interval contains
\(\tau(P)\) almost surely. If \(\mathfrak U_{n,d,q}>0\), define
\[
t_{\mathrm{env}}:=\frac{\mathfrak U_{n,d,q}}{\alpha}>0.
\]
Integrating the pointwise inequality
\(V_{\mathrm{env}}\ge
t_{\mathrm{env}}\mathbf1\{V_{\mathrm{env}}\ge t_{\mathrm{env}}\}\)
gives
\[
t_{\mathrm{env}}\Pr_P(V_{\mathrm{env}}\ge t_{\mathrm{env}})
\le E_P[V_{\mathrm{env}}]\le\mathfrak U_{n,d,q},
\qquad
\Pr_P(V_{\mathrm{env}}\ge t_{\mathrm{env}})\le\alpha.
\]
Because \(\tau(P)\in[-1,1]\), squared error at most
\(t_{\mathrm{env}}\) implies inclusion in the clipped interval.
Consequently,
\[
\{\tau(P)\notin I^{\mathrm{mix}}_\alpha\}
\subseteq\{V_{\mathrm{env}}\ge t_{\mathrm{env}}\},
\qquad
\Pr_P\{\tau(P)\in I^{\mathrm{mix}}_\alpha\}\ge1-\alpha.
\]

The mixed-count estimator takes values in \([-1,1]\) by
\cref{obj:def:mixed-count-estimator}. Its envelope interval therefore
has ordered endpoints in \([-1,1]\). The finite observed-sample space
makes its endpoint map measurable, and its Dirac kernel is an honest
member of \(\mathcal C_{n,d,q,\alpha}\). For this kernel, joint
sample-and-kernel expected length equals
\(E_P[\operatorname{length}(I^{\mathrm{mix}}_\alpha)]\).
The definition in \cref{obj:def:interval-length-risk} hence gives
\[
\mathfrak L_{n,d,q,\alpha}
\le
\sup_{P\in\mathcal M_{n,d,q}}
E_P[\operatorname{length}(I^{\mathrm{mix}}_\alpha)].
\]
Combining this comparison with \cref{proof:connected-interval-upper,proof:connected-interval-lower-scales} proves
\[
\underline c_\alpha\sqrt{r_{n,d,q}}
\le\mathfrak L_{n,d,q,\alpha}
\le
\sup_{P\in\mathcal M_{n,d,q}}
E_P[\operatorname{length}(I^{\mathrm{mix}}_\alpha)]
\le\overline C_\alpha\sqrt{r_{n,d,q}},
\]
together with the asserted coverage.
% lean: CausalSmith.Stat.MarRareqLogfrontier.connected_interval_frontier
\end{enumerate}
\end{proof}

\begin{proof}[Proof of \cref{obj:thm:zeng-fixed-positivity-polynomial-frontier}]
\begin{enumerate}
\item \textit{Uniform risk bound for the deterministic estimator.}
Let \(\overline C>0\) be the universal constant supplied by
\cref{obj:thm:unrestricted-mixed-count-upper}. Fix \(n,d\ge1\),
\(0<\epsilon<1/2\), and \(P_Z\in\mathcal D_{d,\epsilon}\).
Define the auxiliary assignment vector and synthetic sample by
\[
\mathbf u\sim\operatorname{Unif}(\{0,1\}^n),
\qquad
\mathbf O_n^{\mathrm{syn}}
:=(o_i(u_i))_{i=1}^n,
\]
with \(\mathbf u\) independent of \(\mathbf Z_n\).
By \cref{obj:lem:synthetic-randomization-kernel}, there exists
\(P_Z^*\in\mathcal M^+_{n,d,\epsilon}\) such that
\[
\mathcal L_{P_Z,\mathbf u}(\mathbf O_n^{\mathrm{syn}})
=
\bigl(\mathcal L_{P_Z^*}(X,A,S,R,RY)\bigr)^{\otimes n},
\qquad
\tau(P_Z^*)=\theta_Z(P_Z).
\]
Thus the synthetic sample satisfies the sampling and model hypotheses of
\cref{obj:thm:unrestricted-mixed-count-upper} at \(q=\epsilon\).

For each fixed sample realization
\(\mathbf z=((x_i,b_i,y_i))_{i=1}^n\), Cauchy--Schwarz on the
\(2^n\) assignment vectors gives
\[
\begin{split}
&\left[
\sum_{\mathbf u\in\{0,1\}^n}
\left\{
\widehat\tau^{\mathrm{mix}}_{n,d,\epsilon}
\bigl(o_1(u_1),\ldots,o_n(u_n)\bigr)
-\theta_Z(P_Z)
\right\}
\right]^2\\
&\qquad\le
2^n\sum_{\mathbf u\in\{0,1\}^n}
\left\{
\widehat\tau^{\mathrm{mix}}_{n,d,\epsilon}
\bigl(o_1(u_1),\ldots,o_n(u_n)\bigr)
-\theta_Z(P_Z)
\right\}^2.
\end{split}
\]
Dividing by \(2^{2n}\) and using the defining finite average yields
\[
\begin{split}
&\left(
\widehat\theta^{\mathrm{poly}}_{n,d,\epsilon}(\mathbf z)
-\theta_Z(P_Z)
\right)^2\\
&\qquad\le
\frac1{2^n}
\sum_{\mathbf u\in\{0,1\}^n}
\left\{
\widehat\tau^{\mathrm{mix}}_{n,d,\epsilon}
\bigl(o_1(u_1),\ldots,o_n(u_n)\bigr)
-\theta_Z(P_Z)
\right\}^2.
\end{split}
\]
Integrating this inequality, identifying the uniform coin average with
the synthetic sample law, and then applying
\cref{obj:thm:unrestricted-mixed-count-upper}, we obtain
\[
\begin{split}
&E_{P_Z}\!\left[
\left(
\widehat\theta^{\mathrm{poly}}_{n,d,\epsilon}(\mathbf Z_n)
-\theta_Z(P_Z)
\right)^2
\right]\\
&\quad\le
E_{P_Z,\mathbf u}\!\left[
\left(
\widehat\tau^{\mathrm{mix}}_{n,d,\epsilon}
(\mathbf O_n^{\mathrm{syn}})
-\theta_Z(P_Z)
\right)^2
\right]\\
&\quad=
E_{P_Z^*}\!\left[
\left(
\widehat\tau^{\mathrm{mix}}_{n,d,\epsilon}(\mathbf O_n)
-\tau(P_Z^*)
\right)^2
\right]
\le \mathfrak U_{n,d,\epsilon}
\le \overline C\,r_{n,d,\epsilon}.
\end{split}
\]
Taking the supremum over a nonempty observational class proves the
first assertion. If that class is empty, its real-valued supremum is
zero, and the same bound follows from
\[
0\le r_{n,d,\epsilon}
=
\min\left\{
1,\frac1{n\epsilon}
+\left[\frac d{n\epsilon\log(e+n\epsilon)}\right]^2
\right\},
\qquad
0<\overline C.
\]
% lean: CausalSmith.Stat.MarRareqLogfrontier.thetaPoly_uniform_upper

\item \textit{Upper bound at fixed positivity.}
Fix \(0<\epsilon<1/2\). For \(n,d\ge1\), define
\[
r^{\mathrm{raw}}_{n,d}
:=
\frac1n+\left[\frac d{n\log(e+n)}\right]^2,
\qquad
r^{\mathrm{fix}}_{n,d}:=\min\{1,r^{\mathrm{raw}}_{n,d}\},
\]
and set
\[
\ell_{\epsilon,n}:=\log(e+n\epsilon),
\qquad
\ell_{1,n}:=\log(e+n),
\qquad
\gamma_{\epsilon,\mathrm{cmp}}
:=\frac{1-\log\epsilon}{\epsilon}.
\]
All denominators below are positive. In particular,
\[
\ell_{\epsilon,n}\ge1,
\qquad
\ell_{1,n}>0,
\qquad
\log\epsilon\le0,
\qquad
\gamma_{\epsilon,\mathrm{cmp}}\ge1.
\]
Since \(\epsilon\le1\),
\[
e+n\le\frac{e+n\epsilon}{\epsilon}.
\]
Monotonicity of the logarithm therefore gives
\[
\ell_{1,n}
\le \ell_{\epsilon,n}-\log\epsilon
\le (1-\log\epsilon)\ell_{\epsilon,n},
\]
where the second inequality follows from
\[
(-\log\epsilon)(\ell_{\epsilon,n}-1)\ge0.
\]
Consequently,
\[
n\ell_{1,n}
\le
\gamma_{\epsilon,\mathrm{cmp}}\,
n\epsilon\ell_{\epsilon,n}.
\]
Also,
\[
\epsilon\gamma_{\epsilon,\mathrm{cmp}}^2
\ge
\epsilon\gamma_{\epsilon,\mathrm{cmp}}
=
1-\log\epsilon
\ge1.
\]
These two comparisons imply, respectively,
\[
\left[\frac d{n\epsilon\ell_{\epsilon,n}}\right]^2
\le
\gamma_{\epsilon,\mathrm{cmp}}^2
\left[\frac d{n\ell_{1,n}}\right]^2,
\qquad
\frac1{n\epsilon}
\le
\gamma_{\epsilon,\mathrm{cmp}}^2\frac1n.
\]
Adding gives
\[
\frac1{n\epsilon}
+\left[\frac d{n\epsilon\ell_{\epsilon,n}}\right]^2
\le
\gamma_{\epsilon,\mathrm{cmp}}^2 r^{\mathrm{raw}}_{n,d}.
\]
Truncation preserves the needed comparison. Indeed, if
\(r^{\mathrm{raw}}_{n,d}\le1\), then
\[
\min\{1,\gamma_{\epsilon,\mathrm{cmp}}^2r^{\mathrm{raw}}_{n,d}\}
\le
\gamma_{\epsilon,\mathrm{cmp}}^2r^{\mathrm{raw}}_{n,d}
=
\gamma_{\epsilon,\mathrm{cmp}}^2r^{\mathrm{fix}}_{n,d};
\]
if \(r^{\mathrm{raw}}_{n,d}\ge1\), then
\[
\min\{1,\gamma_{\epsilon,\mathrm{cmp}}^2r^{\mathrm{raw}}_{n,d}\}
\le1
\le
\gamma_{\epsilon,\mathrm{cmp}}^2
=
\gamma_{\epsilon,\mathrm{cmp}}^2r^{\mathrm{fix}}_{n,d}.
\]
Thus
\[
r_{n,d,\epsilon}
\le
\gamma_{\epsilon,\mathrm{cmp}}^2r^{\mathrm{fix}}_{n,d}.
\]

The mixed-count estimator takes values in \([-1,1]\), by its clipped
output and averaging construction in \cref{obj:def:mixed-count-estimator}.
Hence, for every sample,
\[
-2^n
\le
\sum_{\mathbf u\in\{0,1\}^n}
\widehat\tau^{\mathrm{mix}}_{n,d,\epsilon}
\bigl(o_1(u_1),\ldots,o_n(u_n)\bigr)
\le2^n,
\]
so
\[
-1\le
\widehat\theta^{\mathrm{poly}}_{n,d,\epsilon}(\mathbf z)
\le1.
\]
This finite-sample map is measurable and is an admissible deterministic
procedure in \cref{obj:def:zeng-minimax-risk}. The defining infimum and
the preceding bounds therefore give
\[
\begin{split}
\mathfrak R^Z_{n,d,\epsilon}
&\le
\sup_{P_Z\in\mathcal D_{d,\epsilon}}
E_{P_Z}\!\left[
\left(
\widehat\theta^{\mathrm{poly}}_{n,d,\epsilon}(\mathbf Z_n)
-\theta_Z(P_Z)
\right)^2
\right]\\
&\le \overline C\,r_{n,d,\epsilon}
\le C_{\epsilon,\mathrm{up}}\,r^{\mathrm{fix}}_{n,d},
\qquad
C_{\epsilon,\mathrm{up}}
:=
\overline C\,\gamma_{\epsilon,\mathrm{cmp}}^2>0.
\end{split}
\]
% lean: CausalSmith.Stat.MarRareqLogfrontier.zeng_minimax_upper_fixedQ

\item \textit{Published lower bound on the zero-control subclass.}
Define
\[
\mathcal D^{(0)}_{d,\epsilon}
:=
\left\{
P_Z\in\mathcal D_{d,\epsilon}:
\mu_{0x}=0\ \text{for every }x\text{ with }p_x>0
\right\},
\]
and its minimax risk by
\[
\mathfrak R^{Z,0}_{n,d,\epsilon}
:=
\inf_{\widehat\theta}
\sup_{P_Z\in\mathcal D^{(0)}_{d,\epsilon}}
E_{P_Z}\!\left[
\bigl(\widehat\theta-\theta_Z(P_Z)\bigr)^2
\right],
\]
using the same procedure class as in
\cref{obj:def:zeng-minimax-risk}. The subclass inclusion gives
\[
\mathfrak R^{Z,0}_{n,d,\epsilon}
\le
\mathfrak R^Z_{n,d,\epsilon}.
\]
The published lower bound
\citep[Theorem~2, Section~3.2, p.~9, and its zero-control construction
in Section~C.5, pp.~39--44]{ZengBalakrishnanHanKennedy2026}
supplies constants
\[
c_{\epsilon,\mathrm{lit}}>0,
\qquad
c_{\epsilon,\mathrm{dim}}>0,
\qquad
n_{\epsilon,\mathrm{lit}}\in\mathbb N
\]
such that
\[
\mathfrak R^{Z,0}_{n,d,\epsilon}
\ge
c_{\epsilon,\mathrm{lit}}
\left\{
\frac1n+\left[\frac d{n\log n}\right]^2
\right\}
\]
whenever
\[
n\ge n_{\epsilon,\mathrm{lit}},
\qquad
1\le d\le c_{\epsilon,\mathrm{dim}}\,n\log n.
\]
The target belongs to \([-1,1]\), since it is a probability-weighted
average of differences of two means in \([0,1]\). By
\cref{obj:lem:randomized-kernel-barycenter}, clipping and then taking
conditional barycenters preserve or reduce squared-error risk at
every law. Thus the deterministic and randomized minimax values agree,
and the published bound applies to the procedure class above.
% lean: CausalSmith.Stat.MarRareqLogfrontier.ZengTheoremTwoLower

\item \textit{A uniform parametric lower bound.}
We next establish a universal constant \(c_{\mathrm{par}}>0\) such that
\[
\frac{c_{\mathrm{par}}}{n}
\le
\mathfrak R^{Z,0}_{n,d,\epsilon}
\le
\mathfrak R^Z_{n,d,\epsilon}
\qquad
(n,d\ge1,\ 0<\epsilon<1/2).
\]
For \(n\ge1\), define
\[
\delta_n:=\frac1{4\sqrt n},
\qquad
\vartheta^{\mathrm{test}}_0:=\frac12,
\qquad
\vartheta^{\mathrm{test}}_1:=\frac12+\delta_n.
\]
Then \(0<\delta_n\le1/4\), so both target values lie in \([0,1]\).
For \(\iota\in\{0,1\}\), define the one-record law
\(P^{\mathrm{test}}_{Z,\iota}\) on
\(\{1,\ldots,d\}\times\{0,1\}^2\) by
\[
P^{\mathrm{test}}_{Z,\iota}\{(x,b,y)\}
=
\begin{cases}
1/2,&(x,b,y)=(1,0,0),\\
(1-\vartheta^{\mathrm{test}}_\iota)/2,
 &(x,b,y)=(1,1,0),\\
\vartheta^{\mathrm{test}}_\iota/2,
 &(x,b,y)=(1,1,1),\\
0,&\text{otherwise}.
\end{cases}
\]
Its sole occupied baseline category has treatment probability \(1/2\),
control mean zero, and treated mean \(\vartheta^{\mathrm{test}}_\iota\).
Thus
\[
P^{\mathrm{test}}_{Z,\iota}\in\mathcal D^{(0)}_{d,\epsilon},
\qquad
\theta_Z(P^{\mathrm{test}}_{Z,\iota})
=\vartheta^{\mathrm{test}}_\iota.
\]

The two laws have common support. Define their chi-square divergence by
\[
\chi^2\!\left(
P^{\mathrm{test}}_{Z,1}\Vert P^{\mathrm{test}}_{Z,0}
\right)
:=
\sum_{\substack{z\in\{1,\ldots,d\}\times\{0,1\}^2:\\
P^{\mathrm{test}}_{Z,0}\{z\}>0}}
\frac{
\bigl(P^{\mathrm{test}}_{Z,1}\{z\}
-P^{\mathrm{test}}_{Z,0}\{z\}\bigr)^2
}{
P^{\mathrm{test}}_{Z,0}\{z\}
}.
\]
Only the two treated-outcome probabilities change. Each changes in
absolute value by \(\delta_n/2\), and each has reference probability
\(1/4\). Hence
\[
\chi^2\!\left(
P^{\mathrm{test}}_{Z,1}\Vert P^{\mathrm{test}}_{Z,0}
\right)
=2\delta_n^2
=\frac1{8n}.
\]
Write the sample laws and their likelihood ratio as
\[
Q^{\mathrm{test}}_\iota
:=
\bigl(P^{\mathrm{test}}_{Z,\iota}\bigr)^{\otimes n},
\qquad \iota\in\{0,1\},
\]
\[
L_{\mathrm{test}}(\mathbf z)
:=
\begin{cases}
Q^{\mathrm{test}}_1\{\mathbf z\}/
Q^{\mathrm{test}}_0\{\mathbf z\},
 &Q^{\mathrm{test}}_0\{\mathbf z\}>0,\\
0,&Q^{\mathrm{test}}_0\{\mathbf z\}=0.
\end{cases}
\]
Independence gives the second-moment identity and bound
\[
E_{Q^{\mathrm{test}}_0}[L_{\mathrm{test}}^2]
=
\left[
1+\chi^2\!\left(
P^{\mathrm{test}}_{Z,1}\Vert P^{\mathrm{test}}_{Z,0}
\right)
\right]^n
=
\left(1+\frac1{8n}\right)^n
\le e^{1/8}.
\]

Define the universal testing constant by
\[
c_{\mathrm{test}}:=\frac1{4e^{1/8}}>0.
\]
For any event \(\mathcal A\) in the finite sample space, set
\[
\eta_{\mathcal A}
:=
Q^{\mathrm{test}}_1(\mathcal A^c)
+
Q^{\mathrm{test}}_0(\mathcal A).
\]
Cauchy--Schwarz yields
\[
1-\eta_{\mathcal A}
\le Q^{\mathrm{test}}_1(\mathcal A)
=
E_{Q^{\mathrm{test}}_0}
[L_{\mathrm{test}}\mathbf1_{\mathcal A}]
\le
\sqrt{e^{1/8}Q^{\mathrm{test}}_0(\mathcal A)}.
\]
If \(\eta_{\mathcal A}<c_{\mathrm{test}}\), then
\(Q^{\mathrm{test}}_0(\mathcal A)\le\eta_{\mathcal A}\) implies
\[
\sqrt{e^{1/8}Q^{\mathrm{test}}_0(\mathcal A)}
<\frac12,
\qquad
1-\eta_{\mathcal A}>\frac34,
\]
because \(c_{\mathrm{test}}\le1/4\). This contradicts the preceding
inequality. Therefore every event satisfies
\[
Q^{\mathrm{test}}_1(\mathcal A^c)
+
Q^{\mathrm{test}}_0(\mathcal A)
\ge c_{\mathrm{test}}.
\]

Consider an arbitrary estimator. By
\cref{obj:lem:randomized-kernel-barycenter}, it suffices to consider its
clipped conditional action kernel \(\kappa\), supported on \([-1,1]\).
Define its barycenter by
\[
b_\kappa(\mathbf z)
:=
\int_{[-1,1]}h\,\kappa(\mathbf z,dh).
\]
The same cited result gives
\[
b_\kappa(\mathbf z)\in[-1,1],
\qquad
E_{Q^{\mathrm{test}}_\iota}
\left[
(b_\kappa-\vartheta^{\mathrm{test}}_\iota)^2
\right]
\le
E_{P^{\mathrm{test}}_{Z,\iota}}
\left[
(\widehat\theta-\vartheta^{\mathrm{test}}_\iota)^2
\right],
\qquad \iota\in\{0,1\},
\]
where the right-hand expectation includes the estimator's
randomization.

Define the nearest-target test and its large-error events by
\[
\mathcal A_{\mathrm{test}}
:=
\left\{
\mathbf z:
|b_\kappa(\mathbf z)-\vartheta^{\mathrm{test}}_1|
\le
|b_\kappa(\mathbf z)-\vartheta^{\mathrm{test}}_0|
\right\},
\]
\[
\mathcal A_{\mathrm{bad},\iota}
:=
\left\{
\mathbf z:
|b_\kappa(\mathbf z)-\vartheta^{\mathrm{test}}_\iota|
\ge\frac{\delta_n}{2}
\right\},
\qquad \iota\in\{0,1\}.
\]
The triangle inequality gives
\[
\delta_n
=
|\vartheta^{\mathrm{test}}_1-\vartheta^{\mathrm{test}}_0|
\le
|b_\kappa-\vartheta^{\mathrm{test}}_0|
+
|b_\kappa-\vartheta^{\mathrm{test}}_1|.
\]
Consequently,
\[
\mathcal A_{\mathrm{test}}\subseteq\mathcal A_{\mathrm{bad},0},
\qquad
\mathcal A_{\mathrm{test}}^c\subseteq\mathcal A_{\mathrm{bad},1}.
\]
Applying the testing bound and then these inclusions yields
\[
c_{\mathrm{test}}
\le
Q^{\mathrm{test}}_1(\mathcal A_{\mathrm{bad},1})
+
Q^{\mathrm{test}}_0(\mathcal A_{\mathrm{bad},0}),
\]
and hence
\[
\frac{c_{\mathrm{test}}}{2}
\le
\max_{\iota\in\{0,1\}}
Q^{\mathrm{test}}_\iota(\mathcal A_{\mathrm{bad},\iota}).
\]
On each large-error event, squared loss is at least
\((\delta_n/2)^2\). Therefore
\[
\begin{split}
\frac{c_{\mathrm{test}}\delta_n^2}{8}
&\le
\left(\frac{\delta_n}{2}\right)^2
\max_{\iota\in\{0,1\}}
Q^{\mathrm{test}}_\iota(\mathcal A_{\mathrm{bad},\iota})\\
&\le
\max_{\iota\in\{0,1\}}
E_{Q^{\mathrm{test}}_\iota}
\left[
(b_\kappa-\vartheta^{\mathrm{test}}_\iota)^2
\right]\\
&\le
\max_{\iota\in\{0,1\}}
E_{P^{\mathrm{test}}_{Z,\iota}}
\left[
(\widehat\theta-\vartheta^{\mathrm{test}}_\iota)^2
\right].
\end{split}
\]
Both testing laws belong to the zero-control subclass, so taking the
minimax infimum gives
\[
\frac{c_{\mathrm{par}}}{n}
=
\frac{c_{\mathrm{test}}\delta_n^2}{8}
\le
\mathfrak R^{Z,0}_{n,d,\epsilon}
\le
\mathfrak R^Z_{n,d,\epsilon},
\qquad
c_{\mathrm{par}}:=\frac{c_{\mathrm{test}}}{128}>0.
\]
% lean: CausalSmith.Stat.MarRareqLogfrontier.zeng_parametric_minimax_lower

\item \textit{Calibration and the large-sample bound within the published range.}
Choose an integer \(n_{\epsilon,1}\) satisfying
\[
n_{\epsilon,1}\ge\frac2{c_{\epsilon,\mathrm{dim}}},
\]
and define
\[
n_{\epsilon,*}
:=
\max\{n_{\epsilon,\mathrm{lit}},3,n_{\epsilon,1}\},
\qquad
c_{\epsilon,\mathrm{sat}}
:=
c_{\epsilon,\mathrm{lit}}
\left(\frac{c_{\epsilon,\mathrm{dim}}}{2}\right)^2,
\]
\[
c_\epsilon
:=
\min\left\{
\frac{c_{\mathrm{par}}}{n_{\epsilon,*}},
\min\{c_{\epsilon,\mathrm{lit}},c_{\epsilon,\mathrm{sat}}\}
\right\}>0.
\]
For all \(n,d\ge1\),
\[
0\le r^{\mathrm{fix}}_{n,d}\le1.
\]

First suppose \(n\ge n_{\epsilon,*}\). Then
\[
n\ge n_{\epsilon,\mathrm{lit}},
\qquad
n\ge3,
\qquad
\log n>1,
\qquad
n\log n>0.
\]
Furthermore, the choice of \(n_{\epsilon,1}\) gives
\[
2\le c_{\epsilon,\mathrm{dim}}\,n
\le c_{\epsilon,\mathrm{dim}}\,n\log n.
\]
If
\[
d\le c_{\epsilon,\mathrm{dim}}\,n\log n,
\]
the published bound and subclass inclusion imply
\[
c_{\epsilon,\mathrm{lit}}
\left\{
\frac1n+\left[\frac d{n\log n}\right]^2
\right\}
\le
\mathfrak R^{Z,0}_{n,d,\epsilon}
\le
\mathfrak R^Z_{n,d,\epsilon}.
\]
Since
\[
0<\log n\le\log(e+n),
\qquad
0\le\frac d{n\log(e+n)}
\le\frac d{n\log n},
\]
we have
\[
r^{\mathrm{fix}}_{n,d}
\le
\frac1n+\left[\frac d{n\log(e+n)}\right]^2
\le
\frac1n+\left[\frac d{n\log n}\right]^2.
\]
Using \(c_\epsilon\le c_{\epsilon,\mathrm{lit}}\) gives
\[
\begin{split}
c_\epsilon r^{\mathrm{fix}}_{n,d}
&\le
c_\epsilon
\left\{
\frac1n+\left[\frac d{n\log n}\right]^2
\right\}\\
&\le
c_{\epsilon,\mathrm{lit}}
\left\{
\frac1n+\left[\frac d{n\log n}\right]^2
\right\}
\le
\mathfrak R^Z_{n,d,\epsilon}.
\end{split}
\]
% lean: CausalSmith.Stat.MarRareqLogfrontier.zeng_minimax_lower_fixedQ

\item \textit{The large-alphabet case.}
Continue to assume \(n\ge n_{\epsilon,*}\), and now suppose
\[
d>c_{\epsilon,\mathrm{dim}}\,n\log n.
\]
Define the smaller alphabet size by
\[
d_{\epsilon,n}^{\mathrm{pad}}
:=
\left\lfloor c_{\epsilon,\mathrm{dim}}\,n\log n\right\rfloor.
\]
The preceding lower bound of two on the quantity inside the floor gives
\[
1\le d_{\epsilon,n}^{\mathrm{pad}}\le d,
\qquad
d_{\epsilon,n}^{\mathrm{pad}}
\le c_{\epsilon,\mathrm{dim}}\,n\log n.
\]
Also, the elementary floor bound gives
\[
\begin{split}
\frac{c_{\epsilon,\mathrm{dim}}}{2}\,n\log n
&\le c_{\epsilon,\mathrm{dim}}\,n\log n-1\\
&<d_{\epsilon,n}^{\mathrm{pad}}.
\end{split}
\]
The published lower bound therefore applies at this smaller alphabet:
\[
c_{\epsilon,\mathrm{lit}}
\left\{
\frac1n+
\left[
\frac{d_{\epsilon,n}^{\mathrm{pad}}}{n\log n}
\right]^2
\right\}
\le
\mathfrak R^{Z,0}_{n,d_{\epsilon,n}^{\mathrm{pad}},\epsilon}
\le
\mathfrak R^Z_{n,d_{\epsilon,n}^{\mathrm{pad}},\epsilon}.
\]

To compare the two alphabet sizes, extend each law
\(P_Z\in\mathcal D_{d_{\epsilon,n}^{\mathrm{pad}},\epsilon}\)
by zero masses. Explicitly, define the \(d\)-category law
\(\operatorname{pad}_{d_{\epsilon,n}^{\mathrm{pad}},d}P_Z\) by
\[
\bigl(\operatorname{pad}_{d_{\epsilon,n}^{\mathrm{pad}},d}P_Z\bigr)
\{(x,b,y)\}
:=
\begin{cases}
P_Z\{(x,b,y)\},&1\le x\le d_{\epsilon,n}^{\mathrm{pad}},\\
0,&d_{\epsilon,n}^{\mathrm{pad}}<x\le d,
\end{cases}
\qquad b,y\in\{0,1\}.
\]
On the original categories, all baseline, treatment, and outcome
probabilities are preserved; the added categories have zero mass.
Thus
\[
\operatorname{pad}_{d_{\epsilon,n}^{\mathrm{pad}},d}P_Z
\in\mathcal D_{d,\epsilon},
\qquad
\theta_Z\!\left(
\operatorname{pad}_{d_{\epsilon,n}^{\mathrm{pad}},d}P_Z
\right)
=\theta_Z(P_Z).
\]
For a \(d\)-category estimator \(\widehat\theta\), define
\(\widehat\theta^{\mathrm{res}}\) on the smaller sample space by
\[
\widehat\theta^{\mathrm{res}}
\bigl((x_i,b_i,y_i)_{i=1}^n\bigr)
:=
\widehat\theta
\bigl((x_i,b_i,y_i)_{i=1}^n\bigr),
\qquad
x_i\in\{1,\ldots,d_{\epsilon,n}^{\mathrm{pad}}\},
\]
preserving its conditional action distribution if it is randomized.
The product sample law is preserved under this inclusion, so
\[
\begin{split}
&E_{P_Z}\!\left[
\left(\widehat\theta^{\mathrm{res}}-\theta_Z(P_Z)\right)^2
\right]\\
&\qquad=
E_{\operatorname{pad}_{d_{\epsilon,n}^{\mathrm{pad}},d}P_Z}
\!\left[
\left(
\widehat\theta-
\theta_Z\!\left(
\operatorname{pad}_{d_{\epsilon,n}^{\mathrm{pad}},d}P_Z
\right)
\right)^2
\right].
\end{split}
\]
Taking suprema and then infima proves
\[
\mathfrak R^Z_{n,d_{\epsilon,n}^{\mathrm{pad}},\epsilon}
\le
\mathfrak R^Z_{n,d,\epsilon}.
\]

Since \(n\log n>0\), the floor bound also gives
\[
\frac{c_{\epsilon,\mathrm{dim}}}{2}
\le
\frac{d_{\epsilon,n}^{\mathrm{pad}}}{n\log n}.
\]
Squaring and using \(1/n\ge0\), we obtain
\[
c_{\epsilon,\mathrm{sat}}
\le
c_{\epsilon,\mathrm{lit}}
\left\{
\frac1n+
\left[
\frac{d_{\epsilon,n}^{\mathrm{pad}}}{n\log n}
\right]^2
\right\}.
\]
Finally, \(r^{\mathrm{fix}}_{n,d}\le1\) and
\(c_\epsilon\le c_{\epsilon,\mathrm{sat}}\) give
\[
\begin{split}
c_\epsilon r^{\mathrm{fix}}_{n,d}
&\le c_\epsilon
\le c_{\epsilon,\mathrm{sat}}\\
&\le
c_{\epsilon,\mathrm{lit}}
\left\{
\frac1n+
\left[
\frac{d_{\epsilon,n}^{\mathrm{pad}}}{n\log n}
\right]^2
\right\}\\
&\le
\mathfrak R^{Z,0}_{n,d_{\epsilon,n}^{\mathrm{pad}},\epsilon}
\le
\mathfrak R^Z_{n,d_{\epsilon,n}^{\mathrm{pad}},\epsilon}
\le
\mathfrak R^Z_{n,d,\epsilon}.
\end{split}
\]
% lean: CausalSmith.Stat.MarRareqLogfrontier.zeng_minimax_lower_fixedQ

\item \textit{Small samples and the final constants.}
If \(1\le n<n_{\epsilon,*}\), the parametric bound and the calibration
of \(c_\epsilon\) imply
\[
c_\epsilon r^{\mathrm{fix}}_{n,d}
\le c_\epsilon
\le\frac{c_{\mathrm{par}}}{n_{\epsilon,*}}
\le\frac{c_{\mathrm{par}}}{n}
\le\mathfrak R^Z_{n,d,\epsilon}.
\]
Together with the two large-sample cases, this establishes the lower
bound for every \(n,d\ge1\). Define
\[
C_\epsilon:=\max\{C_{\epsilon,\mathrm{up}},c_\epsilon\}.
\]
Then \(0<c_\epsilon\le C_\epsilon<\infty\). Since
\(r^{\mathrm{fix}}_{n,d}\ge0\), the upper bound already proved gives
\[
\mathfrak R^Z_{n,d,\epsilon}
\le
C_{\epsilon,\mathrm{up}}r^{\mathrm{fix}}_{n,d}
\le
C_\epsilon r^{\mathrm{fix}}_{n,d}.
\]
Thus, for every \(n,d\ge1\),
\[
c_\epsilon
\min\left\{
1,\frac1n+\left[\frac d{n\log(e+n)}\right]^2
\right\}
\le
\mathfrak R^Z_{n,d,\epsilon}
\le
C_\epsilon
\min\left\{
1,\frac1n+\left[\frac d{n\log(e+n)}\right]^2
\right\},
\]
as claimed.
% lean: CausalSmith.Stat.MarRareqLogfrontier.zeng_fixed_positivity_polynomial_frontier
\end{enumerate}
\end{proof}

\begin{proof}[Proof of \cref{obj:thm:unrestricted-uniform-mixed-count-certificate}]
We take \(\overline C=10^{30}\). We first compare the deterministic envelope with the rate, and then establish the risk bound.

\begin{enumerate}
\item Define the untruncated rate expression by
\[
r^\circ_{n,d,q}
:=N^{-1}+\left(\frac{d}{N\ell}\right)^2,
\qquad
N=nq,\quad m=\frac n6,\quad N_0=\frac N6,
\quad \ell=\log(e+N).
\]
The domain assumptions give
\[
0<N\le n,\qquad m>0,\qquad N_0>0,
\qquad r_{n,d,q}=\min\{1,r^\circ_{n,d,q}\}.
\]
If \(N\le1\), then \(N^{-1}\ge1\), so
\[
\mathfrak U_{n,d,q}=r_{n,d,q}=1
\le10^{30}r_{n,d,q}.
\]
For the envelope comparison, henceforth suppose \(N>1\).
% lean: CausalSmith.Stat.MarRareqLogfrontier.riskEnvelope_le_frontierRate

\item The logarithmic scale and stream sizes satisfy
\[
1<\ell,\qquad e^\ell=e+N,\qquad e^{-\ell}\le N^{-1},
\qquad
\frac1{N_0}=6N^{-1},\qquad
\frac1m=\frac6n\le6N^{-1}.
\]
Indeed, \(e+N>e\), and inversion of \(e^\ell=e+N>N\) gives the exponential bound.

The elementary numerical bound for \(\log2\) gives
\[
\gamma_0=\log2-\frac12\ge\frac18.
\]
For every real \(t_{\mathrm{dec}}>0\), the inequality
\(t_{\mathrm{dec}}\le e^{t_{\mathrm{dec}}}\) implies
\(e^{-t_{\mathrm{dec}}}\le t_{\mathrm{dec}}^{-1}\). Applying this with
\(t_{\mathrm{dec}}=n\gamma_0\) yields
\[
e^{-n\gamma_0}
\le\frac1{n\gamma_0}
\le\frac8n
\le8N^{-1}.
\]
% lean: CausalSmith.Stat.MarRareqLogfrontier.poisson_tail_rate

\item Suppose first that \(\ell\ge128\) and \(d\ge\sqrt N\). On this branch,
\[
k=\lfloor\ell/64\rfloor,\qquad B=256\ell,
\qquad
\frac{\ell}{128}\le k\le\frac{\ell}{64}.
\]
For completeness, for every real \(t_{\mathrm{floor}}\ge1\),
\[
\lfloor t_{\mathrm{floor}}\rfloor\ge\frac{t_{\mathrm{floor}}}{2}.
\]
If \(t_{\mathrm{floor}}<2\), its floor is at least one; if
\(t_{\mathrm{floor}}\ge2\), use
\(\lfloor t_{\mathrm{floor}}\rfloor>t_{\mathrm{floor}}-1
\ge t_{\mathrm{floor}}/2\). This proves the lower bound on \(k\).

The inequality \(2t_{\mathrm{exp}}\le e^{t_{\mathrm{exp}}}\), for real
\(t_{\mathrm{exp}}\), gives
\[
\ell\le16e^{\ell/32},
\qquad
\ell^4\le65536e^{\ell/8}.
\]
Since \(e<3\) and \(N>1\),
\[
e^\ell=e+N\le4N\le4d^2,
\qquad e^{\ell/2}\le2d.
\]
Consequently,
\[
\begin{aligned}
(e^{\ell/8}+1)\ell^4
&\le2e^{\ell/8}\,65536e^{\ell/8}\\
&=131072e^{\ell/4}
\le131072e^{\ell/2}
\le262144d,
\end{aligned}
\qquad
\ell\le e^{\ell/2}\le2d.
\]
% lean: CausalSmith.Stat.MarRareqLogfrontier.needle_growth

\item On the same branch, the preceding bounds imply
\[
e^{-16\ell}\le N^{-1},
\qquad e^{-32\ell}\le N^{-1},
\qquad N^{-1}\le\frac{2d}{N\ell}.
\]
Moreover,
\[
\begin{aligned}
\frac{4dB}{N_0k^2}
&=\frac{6144d\ell}{Nk^2}\\
&\le\frac{6144d\ell}{N(\ell/128)^2}
=100663296\,\frac{d}{N\ell}.
\end{aligned}
\]
Thus
\[
\frac{4dB}{N_0k^2}+3e^{-16\ell}
\le100663302\,\frac{d}{N\ell},
\]
and, because both sides are nonnegative,
\[
b_{n,d,q}^2
\le40532401478172816
\left(\frac{d}{N\ell}\right)^2.
\]
% lean: CausalSmith.Stat.MarRareqLogfrontier.riskEnvelope_le_frontierRate

\item Since \(N_0\le m\) and \(B\ge1\),
\[
\frac{B}{mN_0}
\le\frac{B}{N_0^2}
\le\frac{B^2}{N_0^2}.
\]
Therefore
\[
\frac{B^2}{N_0^2}+\frac{B}{mN_0}
\le\frac{2B^2}{N_0^2}
=4718592\,\frac{\ell^2}{N^2}.
\]
The growth estimate above also gives
\[
\frac{d(e^{\ell/8}+1)\ell^2}{N^2}
=
\frac{d(e^{\ell/8}+1)\ell^4}{N^2\ell^2}
\le262144\left(\frac{d}{N\ell}\right)^2.
\]
Combining these inequalities,
\[
64d(e^{\ell/8}+1)
\left(\frac{B^2}{N_0^2}+\frac{B}{mN_0}\right)
\le79164837199872
\left(\frac{d}{N\ell}\right)^2.
\]
The remaining variance terms satisfy
\[
8\left(\frac4{N_0}+\frac1m\right)\le240N^{-1},
\]
and, using \(N^{-1}\le1\) and \(1/m\le6\),
\[
32e^{-32\ell}\left(1+\frac1m\right)\le224N^{-1}.
\]
Together with the prefix-tail bound, this proves
\[
b_{n,d,q}^2+v_{n,d,q}+4e^{-n\gamma_0}
\le40611566315372688
\left(\frac{d}{N\ell}\right)^2+496N^{-1}
\le10^{30}r^\circ_{n,d,q}.
\]
If \(r^\circ_{n,d,q}\le1\), take the minimum with four on the left.
If \(r^\circ_{n,d,q}>1\), use
\(\mathfrak U_{n,d,q}\le4\le10^{30}\). In either case,
\[
\mathfrak U_{n,d,q}\le10^{30}r_{n,d,q}.
\]
% lean: CausalSmith.Stat.MarRareqLogfrontier.riskEnvelope_le_frontierRate

\item On the remaining branch, \(\ell<128\) or \(d<\sqrt N\). Since \(e\ge1\),
\[
0\le\frac{8d}{eN_0}\le\frac{48d}{N},
\qquad
\left(\frac{8d}{eN_0}\right)^2
\le2304\left(\frac dN\right)^2.
\]
If \(\ell<128\), then
\[
\left(\frac dN\right)^2
=\ell^2\left(\frac{d}{N\ell}\right)^2
\le16384\left(\frac{d}{N\ell}\right)^2.
\]
If \(d<\sqrt N\), then \(d^2\le N\), so
\[
\left(\frac dN\right)^2\le N^{-1}.
\]
Hence throughout this branch,
\[
\left(\frac dN\right)^2
\le N^{-1}+16384\left(\frac{d}{N\ell}\right)^2.
\]
Also,
\[
4\left(\frac4{N_0}+\frac1m\right)\le120N^{-1}.
\]
It follows that
\[
\begin{aligned}
\left(\frac{8d}{eN_0}\right)^2
+4\left(\frac4{N_0}+\frac1m\right)+4e^{-n\gamma_0}
&\le2456N^{-1}
+37748736\left(\frac{d}{N\ell}\right)^2\\
&\le10^{30}r^\circ_{n,d,q}.
\end{aligned}
\]
The same split at \(r^\circ_{n,d,q}=1\) proves
\[
\mathfrak U_{n,d,q}\le10^{30}r_{n,d,q}.
\]
This completes the deterministic comparison over the stated domain.
% lean: CausalSmith.Stat.MarRareqLogfrontier.riskEnvelope_le_frontierRate

\item Fix now \(P\in\mathcal M^+_{n,d,q}\). By
\cref{obj:lem:unrestricted-cell-identification,obj:def:ate-functional},
\[
\Phi(P)=\tau(P)
=P(Y(1)=1)-P(Y(0)=1)\in[-1,1].
\]
The auxiliary output in \cref{obj:def:mixed-count-estimator} lies in
\([-1,1]\). Its conditional average lies in that interval as well: for each
prefix length \(\nu\le n\), there are \(3^\nu\) equally weighted stream
assignments, and
\[
\Pr(K_0=\nu)\ge0,
\qquad
\sum_{\nu=0}^{n}\Pr(K_0=\nu)\le1,
\]
with the remaining probability assigned output zero. Thus
\[
\widehat\tau^{\mathrm{mix}}_{n,d,q}(\mathbf O_n)\in[-1,1],
\qquad
E_P\!\left[
(\widehat\tau^{\mathrm{mix}}_{n,d,q}-\Phi(P))^2
\right]\le4.
\]
If \(N\le1\), the auxiliary output and its conditional average are zero, so
\[
E_P\!\left[
(\widehat\tau^{\mathrm{mix}}_{n,d,q}-\Phi(P))^2
\right]=\Phi(P)^2\le1=\mathfrak U_{n,d,q}.
\]
If \(N>1\) and \(\mathfrak U_{n,d,q}=4\), the displayed bound by four proves
the required inequality. We henceforth consider
\(N>1\) and \(\mathfrak U_{n,d,q}\ne4\).
% lean: CausalSmith.Stat.MarRareqLogfrontier.unrestricted_auxiliary_risk_envelope

\item We first establish the transfer to independent Poisson streams.
Write
\[
\mathsf Q_{\mathrm{str}}
:=
\bigotimes_{h_{\mathrm{str}}\in\{0,1,2\}}
\left[
\sum_{\nu=0}^{\infty}
e^{-m}\frac{m^\nu}{\nu!}
\bigl(\mathcal L_P(O_1)\bigr)^{\otimes\nu}
\right],
\]
where each sum is a probability law on the disjoint union of finite
observed-record sequences. The stream labels \(0,1,2\) denote membership,
pilot, and estimation, respectively.

An uncapped prefix of length \(K_0\sim\operatorname{Poisson}(n/2)\), marked
independently with these three labels, has law \(\mathsf Q_{\mathrm{str}}\)
after separation into streams. Indeed, for any prescribed nonnegative
integer stream sizes \(\nu_0,\nu_1,\nu_2\), put
\(\nu_{\mathrm{tot}}:=\nu_0+\nu_1+\nu_2\). Their joint probability is
\[
e^{-n/2}\frac{(n/2)^{\nu_{\mathrm{tot}}}}{\nu_{\mathrm{tot}}!}
\frac{\nu_{\mathrm{tot}}!}{\nu_0!\nu_1!\nu_2!}
3^{-\nu_{\mathrm{tot}}}
=
\prod_{h_{\mathrm{str}}=0}^{2}
e^{-m}\frac{m^{\nu_{h_{\mathrm{str}}}}}
{\nu_{h_{\mathrm{str}}}!}.
\]
Conditional on these sizes, the records in the three streams have their
respective product observed-data laws. Equality on all count fibres proves
the asserted stream law.

Let
\[
\widetilde\tau_{\mathrm{Pois}}
:=
\operatorname{clip}_{[-1,1]}
\left(2\sum_{j=(a,x,s)\in\mathcal J_d}g(a)W_jG_j\right)
\]
be the uncapped output, using the statistics of
\cref{obj:synth_9,obj:synth_10,obj:synth_11,obj:synth_12,obj:synth_14,obj:synth_15,obj:synth_17}
on these finite streams. On the ratio branch, \(G_j=D_j\).
Here and below \(E_{\mathrm{Pois}}\) and
\(\Pr_{\mathrm{Pois}}\) denote expectation and probability under
\(\mathsf Q_{\mathrm{str}}\).

For the capped auxiliary output, conditional averaging gives
\[
\begin{aligned}
&E_\omega\!\left[
(\widetilde\tau(\mathbf O_n;\omega)-\Phi(P))^2
\mid\mathbf O_n\right]\\
&\quad=
(\widehat\tau^{\mathrm{mix}}_{n,d,q}(\mathbf O_n)-\Phi(P))^2
+
E_\omega\!\left[
(\widetilde\tau(\mathbf O_n;\omega)
-\widehat\tau^{\mathrm{mix}}_{n,d,q}(\mathbf O_n))^2
\mid\mathbf O_n\right].
\end{aligned}
\]
On each prefix-count fibre \(\nu\le n\), its stream law agrees with the
corresponding uncapped law. Summing these fibres, discarding the
nonnegative uncapped loss on the remaining fibres, and bounding the
overflow loss by four yields
\[
E_P\!\left[
(\widehat\tau^{\mathrm{mix}}_{n,d,q}-\Phi(P))^2
\right]
\le
E_{\mathrm{Pois}}\!\left[
(\widetilde\tau_{\mathrm{Pois}}-\Phi(P))^2
\right]+4\Pr(K_0>n).
\]
The Poisson probability series evaluates the exponential moment as
\[
E[2^{K_0}]
=e^{-n/2}\sum_{\nu=0}^{\infty}\frac{n^\nu}{\nu!}
=e^{n/2}.
\]
Since \(\{K_0>n\}\subseteq\{2^{K_0}\ge2^n\}\), exponential Markov gives
\[
\Pr(K_0>n)
\le\Pr(2^{K_0}\ge2^n)
\le2^{-n}E[2^{K_0}]
=e^{-n\log2+n/2}.
\]
Multiplying by \(e^{-n/2}\) cancels the intensity factor, so
\[
\Pr(K_0>n)e^{-n/2}
\le e^{-n\log2},
\qquad
\Pr(K_0>n)\le e^{-n\gamma_0}.
\]
Consequently,
\[
E_P\!\left[
(\widehat\tau^{\mathrm{mix}}_{n,d,q}-\Phi(P))^2
\right]
\le
E_{\mathrm{Pois}}\!\left[
(\widetilde\tau_{\mathrm{Pois}}-\Phi(P))^2
\right]+4e^{-n\gamma_0}.
\]
All these expectations are well defined; the capped sample space is finite
and the clipped outputs are bounded.
% lean: CausalSmith.Stat.MarRareqLogfrontier.deterministicRisk_mixedCountEstimator_le_streamRisk_add_tail

\item For every \(j=(a,x,s)\in\mathcal J_d\), define
\[
p_j:=p_{axs}(P),\qquad
\mu_j:=\mu_{axs}(P),\qquad
t_j:=mP(A=a,X=x,S=s,R=1).
\]
These definitions include null cells, with the mean convention of
\cref{obj:def:cell-functional}. Probability monotonicity and the
occupied-cell arrival condition give
\[
0\le\mu_j\le1,\qquad
t_j\ge N_0p_j,\qquad
\sum_jp_j=1,\qquad |\mathcal J_d|=4d.
\]
For \(p_j=0\), the membership, arrival, and arrived-one probabilities are
zero, and all the associated counts vanish almost surely.

Applying the same Poisson partition calculation to stream and cell
categories gives independent counts across distinct cells and streams,
with
\[
M_j\sim\operatorname{Poisson}(mp_j),
\qquad
C'_j,C_j\sim\operatorname{Poisson}(t_j).
\]
The estimation arrived-one and arrived-zero intensities are, respectively,
\(t_j\mu_j\) and \(t_j(1-\mu_j)\). Poisson moments therefore yield
\[
E_{\mathrm{Pois}}W_j=p_j,\qquad
E_{\mathrm{Pois}}W_j^2=p_j^2+\frac{p_j}{m}.
\]
Also, the membership statistic is independent of every statistic formed
from the pilot and estimation streams.
% lean: CausalSmith.Stat.MarRareqLogfrontier.integral_poissonMemberWeight_mul_selectedCellEstimate

\item Consider the polynomial branch
\(\ell\ge128\), \(d\ge\sqrt N\). Define the pilot-selection event by
\[
\mathsf L_j:=\{C'_j\le B/4\}.
\]
The coefficient and cell-statistic definitions in
\cref{obj:synth_13,obj:synth_14,obj:synth_15} give
\[
1-Q_k(t)=\sum_{v=1}^{k-1}a_vt^v,
\qquad
G_j=\mathbf1_{\mathsf L_j}H_j+
\mathbf1_{\mathsf L_j^c}D_j
\quad(t\ge0).
\]
To justify the expansion at zero and its degree, the numerator
\(1-T_k(1-2t/B)\) vanishes at zero and has degree at most \(k\).
The recurrence for \(T_k\), differentiated at one, gives
\[
T_0'(1)=0,\qquad T_1'(1)=1,\qquad
T_{h+1}'(1)=2+2T_h'(1)-T_{h-1}'(1)
\quad(h\ge1).
\]
Induction yields \(T_h'(1)=h^2\). Division of the numerator by \(t\)
therefore gives
\[
Q_k(0)=1,\qquad \deg(1-Q_k)\le k-1,
\]
which proves the displayed finite expansion.

For integers \(v\ge1\), write
\[
(C_j)_v:=\prod_{h=0}^{v-1}\max\{C_j-h,0\},
\qquad
\mathsf F_{j,v}:=
U_j\prod_{h=1}^{v-1}\max\{C_j-h,0\}.
\]
Empty products equal one. Counting ordered distinct tuples with the first
record an arrived one and the remaining records arrived gives
\[
E_{\mathrm{Pois}}\mathsf F_{j,v}
=t_j^v\mu_j.
\]
More explicitly, condition on the marked prefix size, count its ordered
distinct \(v\)-tuples, and then use the Poisson factorial moment
\(E(K_0)_v=(n/2)^v\). The one-record probabilities are
\(P(A=a,X=x,S=s,R=1,Y=1)/3\) for the distinguished record and
\(P(A=a,X=x,S=s,R=1)/3\) for the others. Their product gives the stated
identity. When the arrived-cell probability is zero, both sides vanish.
Thus
\[
E_{\mathrm{Pois}}H_j
=\mu_j\{1-Q_k(t_j)\}.
\]

For the ratio statistic, conditioning on the pattern of estimation
arrivals gives mean \(\mu_j\) when \(C_j>0\); its zero-count value is zero.
Consequently,
\[
E_{\mathrm{Pois}}D_j=\mu_j(1-e^{-t_j}),
\qquad 0\le D_j\le1.
\]
Pilot independence now gives the exact selected-statistic bias:
\[
E_{\mathrm{Pois}}G_j-\mu_j
=-\mu_j\left\{
\Pr_{\mathrm{Pois}}(\mathsf L_j)Q_k(t_j)
+\Pr_{\mathrm{Pois}}(\mathsf L_j^c)e^{-t_j}
\right\}.
\]
% lean: CausalSmith.Stat.MarRareqLogfrontier.markedPoissonSelectedCellEstimate_bias

\item We next derive the coefficient bound needed for the polynomial
moments. For a real polynomial \(F_{\mathrm{pol}}\), define
\[
\|F_{\mathrm{pol}}\|_{\mathrm{coef}}
:=\sum_{h\ge0}|[t_{\mathrm{pol}}^h]F_{\mathrm{pol}}|,
\]
where \(t_{\mathrm{pol}}\) is the polynomial indeterminate and the sum is
finite. For a second real polynomial and coefficient indices, take
\[
F_{\mathrm{in}}\in\mathbb R[t_{\mathrm{pol}}],
\qquad
h_{\mathrm{coef}},i_{\mathrm{coef}},j_{\mathrm{coef}}\in\mathbb N_0.
\]
All coefficient sums below are finite, with coefficients beyond the
polynomial degree equal to zero. The coefficientwise triangle inequality
 gives
\[
\begin{aligned}
\|F_{\mathrm{pol}}+F_{\mathrm{in}}\|_{\mathrm{coef}}
&=\sum_{h_{\mathrm{coef}}\ge0}
\left|[t_{\mathrm{pol}}^{h_{\mathrm{coef}}}]F_{\mathrm{pol}}
+[t_{\mathrm{pol}}^{h_{\mathrm{coef}}}]F_{\mathrm{in}}\right|\\
&\le\sum_{h_{\mathrm{coef}}\ge0}
\left|[t_{\mathrm{pol}}^{h_{\mathrm{coef}}}]F_{\mathrm{pol}}\right|
+\sum_{h_{\mathrm{coef}}\ge0}
\left|[t_{\mathrm{pol}}^{h_{\mathrm{coef}}}]F_{\mathrm{in}}\right|\\
&=\|F_{\mathrm{pol}}\|_{\mathrm{coef}}
+\|F_{\mathrm{in}}\|_{\mathrm{coef}}.
\end{aligned}
\]
Applying this inequality successively to the finite monomial expansion
of the product yields
\[
\begin{aligned}
\|F_{\mathrm{pol}}F_{\mathrm{in}}\|_{\mathrm{coef}}
&=\left\|
\sum_{i_{\mathrm{coef}}\ge0}\sum_{j_{\mathrm{coef}}\ge0}
\bigl([t_{\mathrm{pol}}^{i_{\mathrm{coef}}}]F_{\mathrm{pol}}\bigr)
\bigl([t_{\mathrm{pol}}^{j_{\mathrm{coef}}}]F_{\mathrm{in}}\bigr)
 t_{\mathrm{pol}}^{i_{\mathrm{coef}}+j_{\mathrm{coef}}}
\right\|_{\mathrm{coef}}\\
&\le\sum_{i_{\mathrm{coef}}\ge0}\sum_{j_{\mathrm{coef}}\ge0}
\left\|
\bigl([t_{\mathrm{pol}}^{i_{\mathrm{coef}}}]F_{\mathrm{pol}}\bigr)
\bigl([t_{\mathrm{pol}}^{j_{\mathrm{coef}}}]F_{\mathrm{in}}\bigr)
 t_{\mathrm{pol}}^{i_{\mathrm{coef}}+j_{\mathrm{coef}}}
\right\|_{\mathrm{coef}}\\
&=\sum_{i_{\mathrm{coef}}\ge0}\sum_{j_{\mathrm{coef}}\ge0}
\left|[t_{\mathrm{pol}}^{i_{\mathrm{coef}}}]F_{\mathrm{pol}}\right|
\left|[t_{\mathrm{pol}}^{j_{\mathrm{coef}}}]F_{\mathrm{in}}\right|\\
&=\left(\sum_{i_{\mathrm{coef}}\ge0}
\left|[t_{\mathrm{pol}}^{i_{\mathrm{coef}}}]F_{\mathrm{pol}}\right|\right)
\left(\sum_{j_{\mathrm{coef}}\ge0}
\left|[t_{\mathrm{pol}}^{j_{\mathrm{coef}}}]F_{\mathrm{in}}\right|\right)\\
&=\|F_{\mathrm{pol}}\|_{\mathrm{coef}}
\|F_{\mathrm{in}}\|_{\mathrm{coef}}.
\end{aligned}
\]
Here the norm of each monomial is the absolute value of its coefficient,
by the defining coefficient sum. For every real polynomial \(F_{\mathrm{in}}\), submultiplicativity gives by induction
\[
\|F_{\mathrm{in}}^h\|_{\mathrm{coef}}
\le\|F_{\mathrm{in}}\|_{\mathrm{coef}}^h
\qquad(h\in\mathbb N_0).
\]
For \(F_{\mathrm{pol}}\ne0\), expand the composition as the finite sum
\[
F_{\mathrm{pol}}\circ F_{\mathrm{in}}
=
\sum_{h=0}^{\deg F_{\mathrm{pol}}}
\bigl([t_{\mathrm{pol}}^h]F_{\mathrm{pol}}\bigr)F_{\mathrm{in}}^h.
\]
For the composition bound, assume
\(\|F_{\mathrm{in}}\|_{\mathrm{coef}}\ge1\).
Subadditivity and the power bound then give
\[
\begin{aligned}
\|F_{\mathrm{pol}}\circ F_{\mathrm{in}}\|_{\mathrm{coef}}
&\le
\sum_{h=0}^{\deg F_{\mathrm{pol}}}
|[t_{\mathrm{pol}}^h]F_{\mathrm{pol}}|
\|F_{\mathrm{in}}^h\|_{\mathrm{coef}}\\
&\le
\sum_{h=0}^{\deg F_{\mathrm{pol}}}
|[t_{\mathrm{pol}}^h]F_{\mathrm{pol}}|
\|F_{\mathrm{in}}\|_{\mathrm{coef}}^h\\
&\le
\sum_{h=0}^{\deg F_{\mathrm{pol}}}
|[t_{\mathrm{pol}}^h]F_{\mathrm{pol}}|
\|F_{\mathrm{in}}\|_{\mathrm{coef}}^{\deg F_{\mathrm{pol}}}\\
&=
\|F_{\mathrm{pol}}\|_{\mathrm{coef}}
\|F_{\mathrm{in}}\|_{\mathrm{coef}}^{\deg F_{\mathrm{pol}}}.
\end{aligned}
\]
When \(F_{\mathrm{pol}}=0\), the composition has coefficient norm zero.
The Chebyshev recurrence implies
\[
\|T_{h+2}\|_{\mathrm{coef}}
\le2\|T_{h+1}\|_{\mathrm{coef}}+\|T_h\|_{\mathrm{coef}},
\qquad
\|T_0\|_{\mathrm{coef}}=\|T_1\|_{\mathrm{coef}}=1.
\]
Since \((1+\sqrt2)^2=2(1+\sqrt2)+1\), induction gives
\[
\|T_h\|_{\mathrm{coef}}\le(1+\sqrt2)^h
\quad(h\in\mathbb N).
\]
Define
\[
\mathcal P_k(t_{\mathrm{pol}})
:=1-T_k(1-2t_{\mathrm{pol}}).
\]
Then
\[
\|\mathcal P_k\|_{\mathrm{coef}}
\le1+(1+\sqrt2)^k3^k\le1+8^k,
\]
using \(\|1-2t_{\mathrm{pol}}\|_{\mathrm{coef}}=3\) and
\(\sqrt2\le3/2\). The scaled quotient satisfies
\[
Q_k(Bt_{\mathrm{pol}})
=\frac1{2k^2}\,
\frac{\mathcal P_k(t_{\mathrm{pol}})}{t_{\mathrm{pol}}},
\]
where division is exact because \(\mathcal P_k(0)=0\).
Removing the zero constant coefficient cannot increase the coefficient
norm. Hence, using \(k\ge2\) and \(8^k\ge64\),
\[
\begin{aligned}
\|1-Q_k(Bt_{\mathrm{pol}})\|_{\mathrm{coef}}
&\le1+\frac{1+8^k}{2k^2}\\
&\le1+\frac{1+8^k}{8}
\le8^k.
\end{aligned}
\]
In particular,
\[
\sum_{v=1}^{k-1}|a_vB^v|\le8^k.
\]
% lean: CausalSmith.Stat.MarRareqLogfrontier.needle_scaled_coefficient_oneNorm_le

\item For \(t\ge0\) and integers \(v,w\ge1\), define the factorial overlap
\[
\mathcal A_{v,w}(t)
:=
\sum_{h=0}^{\min\{v,w\}}
\binom vh\binom wh h!\,t^{v+w-h}.
\]
Counting the \(h\) shared entries of two ordered distinct tuples gives
\[
(C_j)_v(C_j)_w
=
\sum_{h=0}^{\min\{v,w\}}
\binom vh\binom wh h!(C_j)_{v+w-h}.
\]
Taking Poisson factorial moments yields
\[
E_{\mathrm{Pois}}[(C_j)_v(C_j)_w]
=\mathcal A_{v,w}(t_j).
\]
Since \(0\le\mathsf F_{j,v}\le(C_j)_v\), expansion of \(H_j^2\) gives
\[
E_{\mathrm{Pois}}H_j^2
\le
\sum_{v=1}^{k-1}\sum_{w=1}^{k-1}
|a_va_w|\mathcal A_{v,w}(t_j).
\]

For \(0\le t\le B\) and \(v,w\le k\), the bound
\[
\binom vh\binom wh h!\le\frac{(vw)^h}{h!}
\]
implies
\[
\begin{aligned}
\mathcal A_{v,w}(t)
&\le B^{v+w}
\sum_{h=0}^{\min\{v,w\}}
\frac{(k^2/B)^h}{h!}\\
&\le B^{v+w}e^{k^2/B}.
\end{aligned}
\]
Therefore, when \(t_j\le B\),
\[
E_{\mathrm{Pois}}H_j^2
\le e^{k^2/B}
\left(\sum_{v=1}^{k-1}|a_vB^v|\right)^2
\le e^{k^2/B}64^k.
\]
The tuning bounds give
\[
\frac{k^2}{B}\le\frac{\ell}{64},
\qquad 64^k=2^{6k}\le e^{6k},
\]
so
\[
E_{\mathrm{Pois}}H_j^2
\le e^{k^2/B+6k}\le e^{\ell/8}
\quad(t_j\le B).
\]
% lean: CausalSmith.Stat.MarRareqLogfrontier.integral_sq_markedPoissonFactorialBranch_le_exp_of_light

\item A second overlap estimate controls arbitrarily large intensities.
Using \(\binom whh!\le w^h\) and the binomial theorem,
\[
\begin{aligned}
\mathcal A_{v,w}(t)
&\le
\sum_{h=0}^{v}\binom vh w^h t^{v-h}t^w\\
&=(t+w)^vt^w
\le(t+v+w)^{v+w}.
\end{aligned}
\]
For \(v,w\le k\), split according as \((t+2k)/B\le1\) or exceeds one.
In both cases,
\[
\left(\frac{t+2k}{B}\right)^{v+w}
\le
\max\left\{1,\left(\frac{t+2k}{B}\right)^{2k}\right\}.
\]
The coefficient bound thus gives
\[
E_{\mathrm{Pois}}H_j^2
\le64^k
\max\left\{1,
\left(\frac{t_j+2k}{B}\right)^{2k}\right\}.
\]
% lean: CausalSmith.Stat.MarRareqLogfrontier.integral_sq_markedPoissonFactorialBranch_le_full_growth

\item Suppose \(t_j>B\). Define the integer pilot cutoff by
\[
\nu_{\mathrm{cut}}:=\lfloor B/4\rfloor.
\]
Then \(\nu_{\mathrm{cut}}<t_j/4\). The Poisson lower-tail estimate used here
follows from the centered exponential moment. To see its relevant form,
put
\[
a_{\mathrm{tail}}:=\frac1{\sqrt2}.
\]
The elementary exponential inequality
\[
e^{-a_{\mathrm{tail}}}
\le1-a_{\mathrm{tail}}+\frac{a_{\mathrm{tail}}^2}{2}
\]
follows by multiplying
\(e^{a_{\mathrm{tail}}}\ge
1+a_{\mathrm{tail}}+a_{\mathrm{tail}}^2/2\) by the nonnegative quadratic
on the right and using
\[
\left(1+a_{\mathrm{tail}}+\frac{a_{\mathrm{tail}}^2}{2}\right)
\left(1-a_{\mathrm{tail}}+\frac{a_{\mathrm{tail}}^2}{2}\right)
=1+\frac{a_{\mathrm{tail}}^4}{4}\ge1.
\]
For \(C'_j\sim\operatorname{Poisson}(t_j)\), exponential Markov therefore
gives
\[
\begin{aligned}
\Pr_{\mathrm{Pois}}\left(t_j-C'_j>\frac{t_j}{\sqrt2}\right)
&\le
\exp\left\{-\frac{t_j}{2}
+t_j(e^{-a_{\mathrm{tail}}}-1+a_{\mathrm{tail}})\right\}\\
&\le e^{-t_j/4}.
\end{aligned}
\]
On \(\{C'_j\le\nu_{\mathrm{cut}}\}\),
\(t_j-C'_j>3t_j/4>t_j/\sqrt2\). Consequently,
\[
\Pr_{\mathrm{Pois}}(\mathsf L_j)\le e^{-t_j/4}.
\]

Since \(t_j>B\), the maximum in the preceding moment bound selects its
second argument. The inequalities \(64^k\le e^{6k}\) and
\(t_{\mathrm{grow}}\le e^{t_{\mathrm{grow}}}\) for
\(t_{\mathrm{grow}}:=(t_j+2k)/B\) give
\[
\Pr_{\mathrm{Pois}}(\mathsf L_j)E_{\mathrm{Pois}}H_j^2
\le
\exp\left\{-\frac{t_j}{4}
+6k+\frac{2k(t_j+2k)}{B}\right\}.
\]
Here
\[
k\le\frac{t_j}{16384},
\qquad \frac{2k}{B}\le\frac1{8192},
\]
and hence
\[
6k+\frac{2k(t_j+2k)}B
\le
t_j\left(\frac6{16384}+\frac1{8192}
+\frac1{67108864}\right)
\le\frac{t_j}{8}.
\]
Thus
\[
\Pr_{\mathrm{Pois}}(\mathsf L_j)E_{\mathrm{Pois}}H_j^2
\le e^{-t_j/8}\le e^{-32\ell}.
\]
Also \(\Pr_{\mathrm{Pois}}(\mathsf L_j)\le e^{-32\ell}\).
% lean: CausalSmith.Stat.MarRareqLogfrontier.markedPoisson_pilot_light_mul_factorial_secondMoment_le_intensity

\item We now bound the two terms in the selected-statistic bias.
For \(0\le t\le B\), the polynomial quotient in
\cref{obj:synth_18} satisfies
\[
tQ_k(t)=\frac{B}{2k^2}\{1-T_k(1-2t/B)\}.
\]
Because \(|1-2t/B|\le1\), the cosine characterization of \(T_k\) gives
\(|T_k(1-2t/B)|\le1\), and therefore
\[
t|Q_k(t)|\le\frac B{k^2}.
\]
For \(t_j\le B\), use \(N_0p_j\mu_j\le t_j\) to obtain
\[
p_j\mu_j
\Pr_{\mathrm{Pois}}(\mathsf L_j)|Q_k(t_j)|
\le\frac{B}{N_0k^2}.
\]
This includes \(t_j=0\), because then \(p_j=0\).

For \(t_j>B\), Cauchy--Schwarz and the heavy moment estimate give
\[
\begin{aligned}
\left\{
\Pr_{\mathrm{Pois}}(\mathsf L_j)
|E_{\mathrm{Pois}}H_j|
\right\}^2
&\le
\Pr_{\mathrm{Pois}}(\mathsf L_j)
E_{\mathrm{Pois}}H_j^2\\
&\le e^{-t_j/8}.
\end{aligned}
\]
Using
\(\mu_jQ_k(t_j)=\mu_j-E_{\mathrm{Pois}}H_j\), it follows that
\[
\begin{aligned}
\mu_j\Pr_{\mathrm{Pois}}(\mathsf L_j)|Q_k(t_j)|
&\le
\Pr_{\mathrm{Pois}}(\mathsf L_j)
+\Pr_{\mathrm{Pois}}(\mathsf L_j)|E_{\mathrm{Pois}}H_j|\\
&\le2e^{-t_j/16}\le2e^{-16\ell}.
\end{aligned}
\]
Splitting the cell sum at \(t_j=B\), using \(|g(a)|=1\) from
\cref{obj:synth_19}, and then using \(4d\) cells and \(\sum_jp_j=1\), yields
\[
\left|
2\sum_jg(a)p_j\mu_j
\Pr_{\mathrm{Pois}}(\mathsf L_j)Q_k(t_j)
\right|
\le\frac{8dB}{N_0k^2}+4e^{-16\ell}.
\]
% lean: CausalSmith.Stat.MarRareqLogfrontier.abs_sum_pilot_needle_bias_le

\item The complementary-pilot term uses an upper-tail bound with the same
integer cutoff. The Poisson probability series gives
\[
E_{\mathrm{Pois}}[2^{C'_j}]
=e^{-t_j}\sum_{\nu=0}^{\infty}\frac{(2t_j)^\nu}{\nu!}
=e^{t_j}.
\]
The threshold inclusion and exponential Markov inequality imply
\[
\begin{aligned}
\Pr_{\mathrm{Pois}}(C'_j>\nu_{\mathrm{cut}})
&\le\Pr_{\mathrm{Pois}}(2^{C'_j}\ge2^{\nu_{\mathrm{cut}}})\\
&\le2^{-\nu_{\mathrm{cut}}}E_{\mathrm{Pois}}[2^{C'_j}]
=e^{-\nu_{\mathrm{cut}}\log2+t_j}.
\end{aligned}
\]
Multiplying by \(e^{-t_j}\) cancels the intensity factor and yields
\[
\Pr_{\mathrm{Pois}}(C'_j>\nu_{\mathrm{cut}})e^{-t_j}
\le e^{-\nu_{\mathrm{cut}}\log2}.
\]
Since
\[
64\ell<\nu_{\mathrm{cut}}+1,\qquad
\ell\ge128,\qquad \log2>\frac58,
\]
we have \(\nu_{\mathrm{cut}}\log2\ge16\ell\). Thus
\[
\Pr_{\mathrm{Pois}}(\mathsf L_j^c)e^{-t_j}\le e^{-16\ell},
\]
including \(t_j=0\), when the upper-tail probability is zero.
Consequently,
\[
\left|
2\sum_jg(a)p_j\mu_j
\Pr_{\mathrm{Pois}}(\mathsf L_j^c)e^{-t_j}
\right|
\le2e^{-16\ell}.
\]

Define the ideal raw contrast by
\[
\tau_{\mathrm{raw}}
:=2\sum_jg(a)W_jG_j.
\]
Membership independence gives
\(E_{\mathrm{Pois}}(W_jG_j)=p_jE_{\mathrm{Pois}}G_j\).
Combining the exact bias identity with the last two aggregate bounds,
\[
|E_{\mathrm{Pois}}\tau_{\mathrm{raw}}-\Phi(P)|
\le
\frac{8dB}{N_0k^2}+6e^{-16\ell}
=b_{n,d,q}.
\]
% lean: CausalSmith.Stat.MarRareqLogfrontier.abs_integral_raw_poissonSelected_estimator_bias_le

\item We next establish the ratio variance contribution.
For a fixed marked prefix length \(\nu\), define the estimation-arrival
index set by
\[
\mathsf I_j^{(\nu)}
:=
\{i\in\{1,\ldots,\nu\}:
i\text{ has estimation label and is an arrival in cell }j\}.
\]
On a prescribed pattern \(\mathsf I_j^{(\nu)}=\mathsf I_{\mathrm{arr}}\) with
\(\nu_{\mathrm{arr}}:=|\mathsf I_{\mathrm{arr}}|>0\), the restricted product law makes the
arrived-outcome indicators independent with mean \(\mu_j\). Their centered
first moments vanish, and their centered squares are at most one.
Expanding the square therefore gives, for patterns of positive
probability,
\[
E\!\left[
(U_j-\nu_{\mathrm{arr}}\mu_j)^2
\mid\mathsf I_j^{(\nu)}=\mathsf I_{\mathrm{arr}}
\right]\le\nu_{\mathrm{arr}},
\]
and hence
\[
E\!\left[
(D_j-\mu_j)^2
\mid\mathsf I_j^{(\nu)}=\mathsf I_{\mathrm{arr}}
\right]\le\nu_{\mathrm{arr}}^{-1}.
\]
Zero-probability patterns contribute zero. On the zero-count pattern,
\((D_j-\mu_j)^2=\mu_j^2\le1\). Summing first over arrival patterns and then
over the Poisson prefix size gives
\[
E_{\mathrm{Pois}}(D_j-\mu_j)^2
\le
E_{\mathrm{Pois}}\!\left[
\mathbf1_{\{C_j=0\}}
+\frac{\mathbf1_{\{C_j>0\}}}{C_j}
\right],
\]
where the quotient is zero at \(C_j=0\).

If \(t_j\le1\), the integrand on the right is at most one, so its expectation
is at most \(1\le4/(1+t_j)\). If \(t_j>1\), its value is at most
\(2/(C_j+1)\). For every nonnegative integer \(h_{\mathrm{count}}\), the Poisson probabilities satisfy
\[
\Pr_{\mathrm{Pois}}(C_j=h_{\mathrm{count}})
=e^{-t_j}\frac{t_j^{h_{\mathrm{count}}}}{h_{\mathrm{count}}!},
\qquad
\frac{t_j}{h_{\mathrm{count}}+1}
\Pr_{\mathrm{Pois}}(C_j=h_{\mathrm{count}})
=
\Pr_{\mathrm{Pois}}(C_j=h_{\mathrm{count}}+1).
\]
The probability series sums to one, and the reciprocal series is bounded termwise by it. Summing the recurrence and shifting the index therefore gives
\[
\begin{aligned}
t_jE_{\mathrm{Pois}}\frac1{C_j+1}
&=
\sum_{h_{\mathrm{count}}=0}^{\infty}
\frac{t_j}{h_{\mathrm{count}}+1}
\Pr_{\mathrm{Pois}}(C_j=h_{\mathrm{count}})\\
&=
\sum_{h_{\mathrm{count}}=0}^{\infty}
\Pr_{\mathrm{Pois}}(C_j=h_{\mathrm{count}}+1)\\
&=\Pr_{\mathrm{Pois}}(C_j>0)
=1-e^{-t_j}.
\end{aligned}
\]
Dividing by \(t_j>1\) yields
\[
E_{\mathrm{Pois}}\frac1{C_j+1}
=\frac{1-e^{-t_j}}{t_j}.
\]
Therefore
\[
E_{\mathrm{Pois}}(D_j-\mu_j)^2
\le\frac2{t_j}\le\frac4{1+t_j}
\quad(t_j>1).
\]
When the arrived-cell probability is zero, \(C_j=U_j=\mu_j=0\)
almost surely, so the centered second moment is zero. In all cases,
\[
\operatorname{Var}_{\mathrm{Pois}}(D_j)
=
\operatorname{Var}_{\mathrm{Pois}}(D_j-\mu_j)
\le E_{\mathrm{Pois}}(D_j-\mu_j)^2
\le\frac4{1+t_j}.
\]

Expanding the variance of the independent product \(W_jD_j\) gives
\[
\operatorname{Var}_{\mathrm{Pois}}(W_jD_j)
=p_j^2\operatorname{Var}_{\mathrm{Pois}}(D_j)
+\frac{p_j}{m}E_{\mathrm{Pois}}D_j^2.
\]
Using \(D_j^2\le1\) and \(N_0p_j\le t_j\),
\[
\frac{p_j^2}{1+t_j}\le\frac{p_j}{N_0},
\]
including \(p_j=0\). Summing over cells proves
\[
\sum_j\operatorname{Var}_{\mathrm{Pois}}(W_jD_j)
\le\frac4{N_0}+\frac1m.
\]
% lean: CausalSmith.Stat.MarRareqLogfrontier.sum_variance_poissonMemberWeight_mul_ratioBranch_le

\item On the polynomial branch, define the weighted correction by
\[
\Delta_j:=W_j(G_j-D_j)
=W_j\mathbf1_{\mathsf L_j}(H_j-D_j).
\]
Pilot and membership independence give
\[
E_{\mathrm{Pois}}\Delta_j^2
=
\left(p_j^2+\frac{p_j}{m}\right)
\Pr_{\mathrm{Pois}}(\mathsf L_j)
E_{\mathrm{Pois}}(H_j-D_j)^2.
\]
Since \((H_j-D_j)^2\le2(H_j^2+D_j^2)\), for \(t_j\le B\),
\[
E_{\mathrm{Pois}}\Delta_j^2
\le2\left(p_j^2+\frac{p_j}{m}\right)(e^{\ell/8}+1).
\]
Here \(p_j\le B/N_0\), so
\[
p_j^2+\frac{p_j}{m}
\le\frac{B^2}{N_0^2}+\frac{B}{mN_0}.
\]
Summing over at most \(4d\) cells,
\[
8\sum_{j:t_j\le B}E_{\mathrm{Pois}}\Delta_j^2
\le
64d(e^{\ell/8}+1)
\left(\frac{B^2}{N_0^2}+\frac{B}{mN_0}\right).
\]
For \(t_j>B\), the two heavy-pilot estimates instead yield
\[
\Pr_{\mathrm{Pois}}(\mathsf L_j)
E_{\mathrm{Pois}}(H_j-D_j)^2
\le4e^{-32\ell},
\]
and therefore
\[
8\sum_{j:t_j>B}E_{\mathrm{Pois}}\Delta_j^2
\le32e^{-32\ell}\left(1+\frac1m\right).
\]
In the last step we used
\[
\sum_j\left(p_j^2+\frac{p_j}{m}\right)
\le1+\frac1m,
\]
because \(0\le p_j\le1\) and \(\sum_jp_j=1\).
% lean: CausalSmith.Stat.MarRareqLogfrontier.sum_integral_sq_memberWeight_mul_selectedCorrection_le

\item Distinct cells are independent in the uncapped marked experiment.
Since \(g(a)^2=1\),
\[
\operatorname{Var}_{\mathrm{Pois}}(\tau_{\mathrm{raw}})
=4\sum_j\operatorname{Var}_{\mathrm{Pois}}(W_jG_j).
\]
For each cell, \(W_jG_j=W_jD_j+\Delta_j\). Nonnegativity of
\(\operatorname{Var}_{\mathrm{Pois}}(W_jD_j-\Delta_j)\) implies
\[
\operatorname{Var}_{\mathrm{Pois}}(W_jG_j)
\le
2\operatorname{Var}_{\mathrm{Pois}}(W_jD_j)
+2\operatorname{Var}_{\mathrm{Pois}}(\Delta_j)
\le
2\operatorname{Var}_{\mathrm{Pois}}(W_jD_j)
+2E_{\mathrm{Pois}}\Delta_j^2.
\]
The preceding aggregate estimates thus give
\[
\operatorname{Var}_{\mathrm{Pois}}(\tau_{\mathrm{raw}})
\le v_{n,d,q}.
\]
The raw contrast has finite second moment by the finite polynomial moment
bounds already established, and its squared error decomposes as
\[
E_{\mathrm{Pois}}(\tau_{\mathrm{raw}}-\Phi(P))^2
=
\operatorname{Var}_{\mathrm{Pois}}(\tau_{\mathrm{raw}})
+(E_{\mathrm{Pois}}\tau_{\mathrm{raw}}-\Phi(P))^2
\le v_{n,d,q}+b_{n,d,q}^2.
\]
Clipping decreases squared distance to every point of \([-1,1]\): below
\(-1\) it replaces the value by \(-1\), above \(1\) by \(1\), and inside the
interval it preserves the value. Since \(\Phi(P)\in[-1,1]\),
\[
E_{\mathrm{Pois}}(\widetilde\tau_{\mathrm{Pois}}-\Phi(P))^2
\le b_{n,d,q}^2+v_{n,d,q}.
\]
Combining this with the prefix transfer and the bound by four proves
\[
E_P\!\left[
(\widehat\tau^{\mathrm{mix}}_{n,d,q}-\Phi(P))^2
\right]
\le
\min\{4,b_{n,d,q}^2+v_{n,d,q}+4e^{-n\gamma_0}\}
=\mathfrak U_{n,d,q}.
\]
% lean: CausalSmith.Stat.MarRareqLogfrontier.unrestricted_auxiliary_risk_envelope_active

\item Finally, consider the ratio branch. Its ideal raw contrast is
\[
\tau_{\mathrm{raw}}^{\mathrm{ratio}}
:=2\sum_jg(a)W_jD_j.
\]
The ratio mean and membership independence give
\[
E_{\mathrm{Pois}}\tau_{\mathrm{raw}}^{\mathrm{ratio}}-\Phi(P)
=-2\sum_jg(a)p_j\mu_je^{-t_j}.
\]
Because \(t_j\ge N_0p_j\), \(0\le\mu_j\le1\), and
\(t_{\mathrm{bias}}e^{-t_{\mathrm{bias}}}\le e^{-1}\) for every real
\(t_{\mathrm{bias}}\ge0\),
\[
p_j\mu_je^{-t_j}
\le p_je^{-N_0p_j}
\le\frac1{eN_0}.
\]
The exponential inequality follows from
\(e^{t_{\mathrm{bias}}-1}\ge t_{\mathrm{bias}}\).
Thus
\[
|E_{\mathrm{Pois}}\tau_{\mathrm{raw}}^{\mathrm{ratio}}-\Phi(P)|
\le\frac{8d}{eN_0}.
\]
Independence across cells and the ratio variance aggregate give
\[
\operatorname{Var}_{\mathrm{Pois}}
(\tau_{\mathrm{raw}}^{\mathrm{ratio}})
=
4\sum_j\operatorname{Var}_{\mathrm{Pois}}(W_jD_j)
\le4\left(\frac4{N_0}+\frac1m\right).
\]
The squared-error decomposition, clipping, prefix transfer, and bound by
four therefore imply
\[
\begin{aligned}
&E_P\!\left[
(\widehat\tau^{\mathrm{mix}}_{n,d,q}-\Phi(P))^2
\right]\\
&\quad\le
\min\left\{4,
\left(\frac{8d}{eN_0}\right)^2
+4\left(\frac4{N_0}+\frac1m\right)
+4e^{-n\gamma_0}\right\}
=\mathfrak U_{n,d,q}.
\end{aligned}
\]
Together with the zero and capped-envelope cases, this establishes the
risk bound for every admissible \(P\). Combining it with the deterministic
comparison proved above gives
\[
E_P\!\left[
\left(\widehat\tau^{\mathrm{mix}}_{n,d,q}(\mathbf O_n)-\Phi(P)\right)^2
\right]
\le\mathfrak U_{n,d,q}
\le10^{30}r_{n,d,q}.
\]
The expectation is under the canonical product observed-sample law
specified in the statement, and \(10^{30}\) is positive, finite, and
independent of \(P,n,d,q\).
% lean: CausalSmith.Stat.MarRareqLogfrontier.unrestricted_uniform_mixed_count_certificate
\end{enumerate}
\end{proof}

\begin{proof}[Proof of \cref{obj:thm:uniform-mixed-count-certificate}]
\begin{enumerate}
\item Choose the universal constant supplied by
\cref{obj:thm:unrestricted-uniform-mixed-count-certificate}, so that
\[
0<\overline C<\infty.
\]
Fix \(n,d,q\) and \(P\in\mathcal M_{n,d,q}\) satisfying the hypotheses.
% lean: CausalSmith.Stat.MarRareqLogfrontier.unrestricted_uniform_mixed_count_certificate

\item The parameter, randomization, consistency, and arrival conditions in
\cref{obj:def:model-class} give
\[
P\in\mathcal M^{+}_{n,d,q},
\]
by \cref{obj:def:unrestricted-arrival-model-class}. Moreover,
\cref{obj:def:model-class} supplies
\cref{obj:ass:iid-sampling} for the observed sample. Applying
\cref{obj:thm:unrestricted-uniform-mixed-count-certificate} therefore yields
\[
E_P\!\left[
\left(\widehat\tau^{\mathrm{mix}}_{n,d,q}(\mathbf O_n)-\Phi(P)\right)^2
\right]
\le \mathfrak U_{n,d,q},
\qquad
\mathfrak U_{n,d,q}\le \overline C\,r_{n,d,q},
\]
where \(\Phi(P)\) is the signed cell functional of
\cref{obj:def:cell-functional}. The estimator, rate, and deterministic
envelope in that certificate are precisely those in the present statement,
with the same formulas on each branch.
% lean: CausalSmith.Stat.MarRareqLogfrontier.uniform_mixed_count_certificate

\item Since \(P\in\mathcal M^{+}_{n,d,q}\),
\cref{obj:lem:unrestricted-cell-identification} gives
\[
\tau(P)=\Phi(P).
\]
Substituting this identity into the risk bound proves
\[
E_P\!\left[
\left(\widehat\tau^{\mathrm{mix}}_{n,d,q}(\mathbf O_n)-\tau(P)\right)^2
\right]
=
E_P\!\left[
\left(\widehat\tau^{\mathrm{mix}}_{n,d,q}(\mathbf O_n)-\Phi(P)\right)^2
\right]
\le \mathfrak U_{n,d,q}
\le \overline C\,r_{n,d,q}.
\]
The chosen constant is universal, so these inequalities hold throughout
the stated model class.
% lean: CausalSmith.Stat.MarRareqLogfrontier.unrestricted_cell_identification
\end{enumerate}
\end{proof}

\begin{proof}[Proof of \cref{obj:thm:mixed-count-upper}]
\begin{enumerate}
\item Choose the universal constant \(\overline C>0\) supplied by
\cref{obj:thm:uniform-mixed-count-certificate}. Fix \(n,d,q\) and a full-data law \(P\) satisfying the stated conditions. Membership in \(\mathcal M_{n,d,q}\) supplies all hypotheses of that certificate.
By \cref{obj:lem:cell-identification-background}, the causal target equals the total observed-cell functional on this model class, which is the identification step used by the certificate. The certificate therefore gives
\[
E_P\!\left[
\left(
\widehat\tau^{\mathrm{mix}}_{n,d,q}(\mathbf O_n)-\tau(P)
\right)^2
\right]
\le \mathfrak U_{n,d,q}
\le \overline C\,r_{n,d,q}.
\]
The envelope and branch rules in that certificate are exactly those in the present statement. It remains to establish the comparison with the auxiliary randomized risk.
% lean: CausalSmith.Stat.MarRareqLogfrontier.uniform_mixed_count_certificate

\item Write the prefix probabilities and their remaining mass as
\[
w_{k_0}:=\Pr(K_0=k_0)
=e^{-n/2}\frac{(n/2)^{k_0}}{k_0!},
\qquad k_0\in\mathbb N,
\qquad
w_{\mathrm{tail}}:=1-\sum_{k_0=0}^{n}w_{k_0}.
\]
The singleton prefix events are disjoint, so
\[
w_{k_0}\ge0,
\qquad
\sum_{k_0=0}^{n}w_{k_0}\le1,
\qquad
w_{\mathrm{tail}}=\Pr(K_0>n)\ge0.
\]
For each integer \(0\le k_0\le n\), define the assignment set by
\[
\mathcal A_{\mathrm{str},k_0}
:=
\left\{
\sigma:\{1,\ldots,k_0\}
\longrightarrow
\{\mathrm{membership},\mathrm{pilot},\mathrm{estimation}\}
\right\},
\qquad
|\mathcal A_{\mathrm{str},k_0}|=3^{k_0}.
\]
Let the output at a prescribed prefix and assignment, and its assignment average, be
\[
\widetilde\tau_{k_0,\sigma}(\mathbf O_n)
:=
\widetilde\tau(\mathbf O_n;\omega)
\quad\text{when }\omega\text{ prescribes }K_0=k_0
\text{ and assignment }\sigma,
\]
\[
\overline\tau_{\mathrm{str},k_0}(\mathbf O_n)
:=
\frac{1}{3^{k_0}}
\sum_{\sigma\in\mathcal A_{\mathrm{str},k_0}}
\widetilde\tau_{k_0,\sigma}(\mathbf O_n).
\]
These outputs use the branch rules of
\cref{obj:def:mixed-count-estimator}. For \(k_0=0\), the assignment set contains the unique empty assignment.

The zero output on the prefix tail and the conditional-average definition give
\[
\widehat\tau^{\mathrm{mix}}_{n,d,q}(\mathbf O_n)
=
\sum_{k_0=0}^{n}
w_{k_0}\overline\tau_{\mathrm{str},k_0}(\mathbf O_n).
\]
Pointwise in the observed sample, its centered value therefore satisfies
\[
\widehat\tau^{\mathrm{mix}}_{n,d,q}(\mathbf O_n)-\tau(P)
=
w_{\mathrm{tail}}\bigl(-\tau(P)\bigr)
+
\sum_{k_0=0}^{n}
w_{k_0}
\bigl(\overline\tau_{\mathrm{str},k_0}(\mathbf O_n)-\tau(P)\bigr).
\]
The weights on the right are nonnegative and sum to one. Convexity of the square yields
\[
\left(
\widehat\tau^{\mathrm{mix}}_{n,d,q}(\mathbf O_n)-\tau(P)
\right)^2
\le
w_{\mathrm{tail}}\tau(P)^2
+
\sum_{k_0=0}^{n}
w_{k_0}
\left(
\overline\tau_{\mathrm{str},k_0}(\mathbf O_n)-\tau(P)
\right)^2.
\]
% lean: CausalSmith.Stat.MarRareqLogfrontier.finite_weighted_square_with_remainder

\item For every \(0\le k_0\le n\), centering the assignment average gives
\[
\overline\tau_{\mathrm{str},k_0}(\mathbf O_n)-\tau(P)
=
\frac{1}{3^{k_0}}
\sum_{\sigma\in\mathcal A_{\mathrm{str},k_0}}
\left(
\widetilde\tau_{k_0,\sigma}(\mathbf O_n)-\tau(P)
\right).
\]
The finite Cauchy--Schwarz inequality consequently gives
\[
\left(
\overline\tau_{\mathrm{str},k_0}(\mathbf O_n)-\tau(P)
\right)^2
\le
\frac{1}{3^{k_0}}
\sum_{\sigma\in\mathcal A_{\mathrm{str},k_0}}
\left(
\widetilde\tau_{k_0,\sigma}(\mathbf O_n)-\tau(P)
\right)^2.
\]
Multiplying these inequalities by the nonnegative \(w_{k_0}\), summing, and inserting the result into the preceding bound yields
\[
\begin{aligned}
&\left(
\widehat\tau^{\mathrm{mix}}_{n,d,q}(\mathbf O_n)-\tau(P)
\right)^2\\
&\quad\le
w_{\mathrm{tail}}\tau(P)^2
+
\sum_{k_0=0}^{n}\frac{w_{k_0}}{3^{k_0}}
\sum_{\sigma\in\mathcal A_{\mathrm{str},k_0}}
\left(
\widetilde\tau_{k_0,\sigma}(\mathbf O_n)-\tau(P)
\right)^2.
\end{aligned}
\]
% lean: CausalSmith.Stat.MarRareqLogfrontier.finite_uniform_square

\item Conditional on the observed sample, the prefix has probabilities \(w_{k_0}\), each assignment has probability \(3^{-k_0}\), and the tail output is zero. Thus
\[
\begin{aligned}
&E_\omega\!\left[
\left(\widetilde\tau(\mathbf O_n;\omega)-\tau(P)\right)^2
\mid\mathbf O_n
\right]\\
&\quad=
w_{\mathrm{tail}}\tau(P)^2
+
\sum_{k_0=0}^{n}\frac{w_{k_0}}{3^{k_0}}
\sum_{\sigma\in\mathcal A_{\mathrm{str},k_0}}
\left(
\widetilde\tau_{k_0,\sigma}(\mathbf O_n)-\tau(P)
\right)^2.
\end{aligned}
\]
The observed-sample space is finite, so the displayed functions are measurable and integrable. Integrating the pointwise inequality establishes
\[
E_P\!\left[
\left(
\widehat\tau^{\mathrm{mix}}_{n,d,q}(\mathbf O_n)-\tau(P)
\right)^2
\right]
\le
E_{P,\omega}\!\left[
\left(
\widetilde\tau(\mathbf O_n;\omega)-\tau(P)
\right)^2
\right].
\]
Together with the certificate bound, this proves all the asserted inequalities with the same universal constant \(\overline C\).
% lean: CausalSmith.Stat.MarRareqLogfrontier.mixed_count_derandomization
% lean: CausalSmith.Stat.MarRareqLogfrontier.mixed_count_upper
\end{enumerate}
\end{proof}

\begin{proof}[Proof of \cref{obj:thm:full-experiment-lower}]
Let \(\underline c_{\mathrm{one}}>0\) be the universal constant supplied by
\cref{obj:lem:one-cell-testing-family}. Define the universal constants
\[
\eta=e^{-32},\qquad
N_{\mathrm{act}}
=\max\left\{1,e^2,4\eta,4096\eta,
e^{(32+\log 8)/28}\right\},
\]
\[
\underline c_{\mathrm{act}}
=\min\left\{
\frac{1}{256N_{\mathrm{act}}},
\frac{1}{256\cdot2048^2},
\frac{\eta^2}{16384}
\right\},
\qquad
\underline c_*=\min\{\underline c_{\mathrm{one}},
\underline c_{\mathrm{act}}\},
\qquad
\underline c=\frac{\underline c_*}{4}.
\]
These constants are positive and independent of \(n,d,q\).

Fix parameters satisfying the hypotheses, and define
\[
N=nq,\qquad \ell=\log(e+N),\qquad
u_{\mathrm{act}}=\frac{d}{N\ell},\qquad
a_{\mathrm{eff}}=N^{-1},\qquad
b_{\mathrm{dim}}=u_{\mathrm{act}}^2.
\]
Then
\[
N>0,\qquad
\ell\ge\log e=1,\qquad
a_{\mathrm{eff}}\ge0,\qquad b_{\mathrm{dim}}\ge0.
\]
By \cref{obj:def:frontier-rate},
\[
r_{n,d,q}=\min\{1,a_{\mathrm{eff}}+b_{\mathrm{dim}}\}.
\]

\begin{enumerate}
\item \emph{The one-cell contribution.}
By \cref{obj:lem:one-cell-testing-family},
\[
\underline c_{\mathrm{one}}\min\{1,a_{\mathrm{eff}}\}
\le \mathfrak R_{n,d,q}.
\]
We next establish the complementary bound
\[
\sqrt N\le d
\quad\Longrightarrow\quad
\underline c_{\mathrm{act}}\min\{1,b_{\mathrm{dim}}\}
\le \mathfrak R_{n,d,q}.
\]
% lean: CausalSmith.Stat.MarRareqLogfrontier.one_cell_testing_family

\item \emph{A numerical one-cell estimate for constant calibration.}
We record explicitly the numerical estimate used below:
\[
\frac1{256}\min\{1,N^{-1}\}\le\mathfrak R_{n,d,q}.
\]
For \(\mu_{\mathrm{one}}\in[1/4,3/4]\), define \(P_{\mathrm{one},\mu_{\mathrm{one}}}\) by
\[
X=0,\qquad S(0)=S(1)=Y(0)=0,\qquad
(A,Y(1),R)\sim
\operatorname{Bernoulli}(1/2)\otimes
\operatorname{Bernoulli}(\mu_{\mathrm{one}})\otimes
\operatorname{Bernoulli}(q),
\]
\[
S=0,\qquad Y=AY(1).
\]
Treatment is independent of the potential variables, and arrival is independent
of the outcome conditional on the observed cell. Every occupied cell has arrival
probability \(q\). Consequently,
\[
P_{\mathrm{one},\mu_{\mathrm{one}}}\in\mathcal M_{n,d,q},
\qquad
\tau(P_{\mathrm{one},\mu_{\mathrm{one}}})=\mu_{\mathrm{one}}.
\]
Write
\[
Q_{\mathrm{one},\mu_{\mathrm{one}}}=\mathcal L_{P_{\mathrm{one},\mu_{\mathrm{one}}}}(O_1),
\qquad
\delta_{\mathrm{one}}
=\min\left\{\frac14,\frac{1}{4\sqrt N}\right\}.
\]
Thus \(0<\delta_{\mathrm{one}}\le1/4\), and the two target values
\(1/2\) and \(1/2+\delta_{\mathrm{one}}\) belong to \([1/4,3/4]\).
Considering the two branches of this minimum gives
\[
\delta_{\mathrm{one}}^2
\ge \frac1{16}\min\{1,N^{-1}\},
\qquad
16N\delta_{\mathrm{one}}^2\le1.
\]

For probability laws \(Q'\ll Q\), use the convention
\[
\chi^2(Q'\Vert Q)
=\int\left(\frac{dQ'}{dQ}-1\right)^2\,dQ.
\]
The two one-record laws differ precisely at the treated, arrived records with
outcome zero or one. Each of these atoms has probability \(q/4\) under
\(Q_{\mathrm{one},1/2}\), and its probability changes in absolute value by
\(q\delta_{\mathrm{one}}/2\). Summing their contributions gives
\[
\chi^2\!\left(
Q_{\mathrm{one},1/2+\delta_{\mathrm{one}}}
\Vert Q_{\mathrm{one},1/2}\right)
=2q\delta_{\mathrm{one}}^2
\le\frac{1}{8n}.
\]
Independence of the records factors the second moment of the likelihood ratio,
so
\[
1+\chi^2\!\left(
Q_{\mathrm{one},1/2+\delta_{\mathrm{one}}}^{\otimes n}
\Vert Q_{\mathrm{one},1/2}^{\otimes n}\right)
=
\left[
1+\chi^2\!\left(
Q_{\mathrm{one},1/2+\delta_{\mathrm{one}}}
\Vert Q_{\mathrm{one},1/2}\right)
\right]^n
\le e^{1/8}.
\]
Since \(e^{1/8}\le e^{1/2}<\sqrt3<2\), Cauchy--Schwarz applied to the
likelihood ratio yields
\[
\operatorname{TV}\!\left(
Q_{\mathrm{one},1/2+\delta_{\mathrm{one}}}^{\otimes n},
Q_{\mathrm{one},1/2}^{\otimes n}\right)
\le
\frac12\sqrt{
\chi^2\!\left(
Q_{\mathrm{one},1/2+\delta_{\mathrm{one}}}^{\otimes n}
\Vert Q_{\mathrm{one},1/2}^{\otimes n}\right)}
\le\frac12.
\]

For every \(\mathsf T\in\mathcal T_n\), define its barycenter on the
observed-sample space by
\[
b_{\mathsf T}(o)=\int_{[-1,1]}h\,\mathsf T(o,dh).
\]
By \cref{obj:lem:randomized-kernel-barycenter},
\[
b_{\mathsf T}(o)\in[-1,1],
\qquad
E_P[(b_{\mathsf T}(\mathbf O_n)-\tau(P))^2]
\le E_P[(\mathsf T-\tau(P))^2].
\]
Define the nearer-target test
\[
\psi_{\mathrm{one}}(o)
=
\mathbf1\left\{
\left|b_{\mathsf T}(o)-\left(\frac12+\delta_{\mathrm{one}}\right)\right|
\le
\left|b_{\mathsf T}(o)-\frac12\right|
\right\}.
\]
An error under either target entails absolute estimation error at least
\(\delta_{\mathrm{one}}/2\). The testing inequality in
\cref{obj:lem:randomized-kernel-tv-comparison}, followed by the barycenter
risk comparison, therefore gives
\[
\begin{aligned}
&\max\left\{
E_{P_{\mathrm{one},1/2}}\!\left[(\mathsf T-\tfrac12)^2\right],
E_{P_{\mathrm{one},1/2+\delta_{\mathrm{one}}}}\!
\left[(\mathsf T-\tfrac12-\delta_{\mathrm{one}})^2\right]
\right\}\\
&\qquad\ge
\frac{\delta_{\mathrm{one}}^2}{8}
\left[
1-\operatorname{TV}\!\left(
Q_{\mathrm{one},1/2+\delta_{\mathrm{one}}}^{\otimes n},
Q_{\mathrm{one},1/2}^{\otimes n}\right)
\right]
\ge\frac{\delta_{\mathrm{one}}^2}{16}
\ge\frac1{256}\min\{1,N^{-1}\}.
\end{aligned}
\]
Both laws belong to the model class. Taking the supremum over that class and
then the infimum over \(\mathsf T\) proves the displayed numerical estimate.
% lean: CausalSmith.Stat.MarRareqLogfrontier.oneCell_minimax_risk_floors

\item \emph{Large effective size and many rare cells.}
To prove the complementary bound, suppose \(\sqrt N\le d\).
Use the construction parameters of
\cref{obj:def:activation-lower-handle}:
\[
K=\lceil\ell\rceil,\qquad H=K^2,\qquad
b_0=\frac{\eta}{N\ell},\qquad
J=\min\left\{d-1,\left\lfloor\frac1{2b_0}\right\rfloor\right\}.
\]
Here \(b_0>0\), \(0\le J\le d-1\), and
\[
Jb_0\le\frac12.
\]
The ceiling bound and \(\ell\ge1\) give
\[
\ell\le K<\ell+1\le2\ell,\qquad H\ge1.
\]
The rare-arrival restriction then implies
\[
0<qH\le4q\ell^2\le\frac1{16}.
\]

First suppose \(N_{\mathrm{act}}\le N\). In particular,
\[
N\ge1,\qquad N\ge e^2,\qquad
N\ge4\eta,\qquad N\ge4096\eta,\qquad
N\ge e^{(32+\log8)/28}.
\]
The capacity bound is
\[
2048(2b_0)=\frac{4096\eta}{N\ell}\le1,
\qquad
2048\le\frac1{2b_0}.
\]
Within this branch, first suppose \(J\ge2048\). Since \(N\ge e^2\),
\[
\ell=\log(e+N)\ge2,\qquad K\ge2.
\]
By \cref{obj:lem:moment-matched-reciprocal-priors}, there are finitely
supported probability laws \(\Pi_0,\Pi_1\) on \([1,H]\) satisfying
\[
\int z^v\,d\Pi_0(z)=\int z^v\,d\Pi_1(z)
\qquad(0\le v\le K,\ v\in\mathbb N).
\]
If the reciprocal mean under \(\Pi_0\) exceeds that under \(\Pi_1\), interchange
the priors; otherwise retain their order. Their separation then reads
\[
\int z^{-1}\,d\Pi_1(z)-\int z^{-1}\,d\Pi_0(z)
\ge\frac1{12}.
\]

For each \(\mathbf z=(z_x)_{x=0}^{d-1}\in[1,H]^d\), define a full-data law
\(P_{\mathrm{act},\mathbf z}\) as follows. Its baseline probabilities are
\[
p_{\mathrm{act},x}=
\begin{cases}
b_0,&0\le x<J,\\
1-Jb_0,&x=J,\\
0,&J<x<d.
\end{cases}
\]
Set \(S(0)=S(1)=Y(0)=0\), draw \(A\sim\operatorname{Bernoulli}(1/2)\)
independently of the baseline and potential variables, and impose \(S=0\) and
\(Y=AY(1)\). Conditional on \(X=x\), draw \(Y(1)\) and \(R\) independently with
\[
Y(1)\mid X=x\sim
\begin{cases}
\operatorname{Bernoulli}(z_x^{-1}),&x<J,\\
\operatorname{Bernoulli}(0),&x\ge J,
\end{cases}
\qquad
R\mid X=x\sim
\begin{cases}
\operatorname{Bernoulli}(qz_x),&x<J,\\
\operatorname{Bernoulli}(qH),&x\ge J.
\end{cases}
\]
The baseline masses sum to one and are nonnegative. Since \(1\le z_x\le H\),
the occupied-cell arrival probabilities belong to \([q,1]\). The specified
independences give randomized assignment and conditional independence of
arrival and outcome, and the definitions of \(S,Y\) give consistency.
With independent sampling from this fixed law, all model conditions hold:
\[
P_{\mathrm{act},\mathbf z}\in\mathcal M_{n,d,q},
\qquad
\tau_{\mathrm{act}}(\mathbf z)
:=\tau(P_{\mathrm{act},\mathbf z})
=b_0\sum_{x=0}^{J-1}z_x^{-1}.
\]

For \(\iota\in\{0,1\}\), draw
\[
\mathbf Z=(Z_x)_{x=0}^{d-1}\sim\Pi_\iota^{\otimes d}
\]
once for the entire sample. Define the observed mixture law
\[
Q_{\mathrm{act},\iota}
=\int_{[1,H]^d}
\mathcal L_{P_{\mathrm{act},\mathbf z}}(\mathbf O_n)\,
d\Pi_\iota^{\otimes d}(\mathbf z),
\]
and let \(\widetilde Q_{\mathrm{act},\iota}\) be the corresponding augmented
law constructed in \cref{obj:def:activation-lower-handle}. Thus
\[
\widetilde Q_{\mathrm{act},\iota}
=
\int_{[1,H]^d}
\mathcal L\!\left(
(\mathsf B_{\mathrm{act},t},O_t)_{t=1}^{n}
\,\middle|\,\mathbf Z=\mathbf z\right)
d\Pi_\iota^{\otimes d}(\mathbf z).
\]
The activation flags have probability \(qH\), independently of the latent
vector and the baseline and treatment pattern.

At the record level, forgetting the flag gives the following probabilities
for the three marks \((R,RY)=(1,1),(1,0),(0,0)\):
\[
\begin{array}{c|ccc}
& (1,1)&(1,0)&(0,0)\\ \hline
A=1,\ X=x<J&q&q(z_x-1)&1-qz_x\\
A=0,\ X=x<J&0&qz_x&1-qz_x\\
X\ge J&0&qH&1-qH
\end{array}
\]
These are exactly the probabilities induced by \(P_{\mathrm{act},\mathbf z}\).
Define the forgetting map by
\[
\operatorname{pr}_{\mathrm{obs}}
\bigl((\mathsf B_{\mathrm{act},t},O_t)_{t=1}^{n}\bigr)
=(O_t)_{t=1}^{n}.
\]
Conditional independence of the records and integration over the common
latent vector therefore give
\[
Q_{\mathrm{act},\iota}
=\widetilde Q_{\mathrm{act},\iota}
\circ\operatorname{pr}_{\mathrm{obs}}^{-1}.
\]
% lean: CausalSmith.Stat.MarRareqLogfrontier.interval_ordered_reciprocal_priors

\item \emph{Equality on the moment-matching event.}
Continue in the branch \(N\ge N_{\mathrm{act}}\), \(J\ge2048\).
For an augmented sample, define the activated count in cell \(x\) and the
cutoff event by
\[
V_x=\sum_{t=1}^{n}
\mathbf1\{X_t=x,\ \mathsf B_{\mathrm{act},t}=1\},
\qquad
\mathcal G=\{V_x\le K\text{ for every }0\le x<J\}.
\]

Fix a possible augmented sample \(\mathbf o_{\mathrm{aug}}\), and denote its
coordinates by the corresponding record symbols. Its baseline, treatment,
surrogate, and flag pattern has the common probability
\[
\varpi_{\mathbf o_{\mathrm{aug}}}
=
\mathbf1\{S_t=0\text{ for every }t\}
\prod_{t=1}^{n}
\frac{p_{\mathrm{act},X_t}}2
(qH)^{\mathsf B_{\mathrm{act},t}}
(1-qH)^{1-\mathsf B_{\mathrm{act},t}}.
\]
For each cell \(x\), define \(\mathcal P_{\mathbf o_{\mathrm{aug}},x}(z)\)
as the product, over records with \(X_t=x\), of their conditional mark
probabilities from \cref{obj:def:activation-lower-handle}:
\[
\mathcal P_{\mathbf o_{\mathrm{aug}},x}(z)
=
\prod_{t:X_t=x}
\Pr\!\left(
(R,RY)=(R_t,R_tY_t)
\,\middle|\,
X=x,\ A=A_t,\ \mathsf B_{\mathrm{act}}=\mathsf B_{\mathrm{act},t},
\ Z_x=z
\right).
\]
Empty products equal one. The row entries extend these expressions to
polynomials on \(\mathbb R\); their probability interpretation applies on
\([1,H]\). Impossible marks have weight zero. For a rare cell, an inactive
factor is constant and an activated factor is affine in \(z\). For every
other cell the factors are constant. Hence
\[
\deg\mathcal P_{\mathbf o_{\mathrm{aug}},x}\le V_x
\quad(x<J),\qquad
\deg\mathcal P_{\mathbf o_{\mathrm{aug}},x}\le0
\quad(x\ge J),
\]
where a zero polynomial satisfies each stated degree bound.

On \(\mathcal G\), define the coefficients
\(\beta_{\mathbf o_{\mathrm{aug}},x,v}\) by the polynomial expansion
\[
\mathcal P_{\mathbf o_{\mathrm{aug}},x}(z)
=\sum_{v=0}^{K}
\beta_{\mathbf o_{\mathrm{aug}},x,v}z^v
\qquad(z\in\mathbb R),
\]
padding with zero coefficients when necessary. Product-prior independence
then gives
\[
\widetilde Q_{\mathrm{act},\iota}\{\mathbf o_{\mathrm{aug}}\}
=
\varpi_{\mathbf o_{\mathrm{aug}}}
\prod_{x=0}^{d-1}
\left[
\sum_{v=0}^{K}
\beta_{\mathbf o_{\mathrm{aug}},x,v}
\int z^v\,d\Pi_\iota(z)
\right]
\qquad(\mathbf o_{\mathrm{aug}}\in\mathcal G).
\]
The matched moments make the right-hand side identical for
\(\iota=0,1\). Thus
\[
\widetilde Q_{\mathrm{act},0}\{\mathbf o_{\mathrm{aug}}\}
=
\widetilde Q_{\mathrm{act},1}\{\mathbf o_{\mathrm{aug}}\}
\qquad(\mathbf o_{\mathrm{aug}}\in\mathcal G).
\]
% lean: CausalSmith.Stat.MarRareqLogfrontier.activatedAugmentedSample_prior_singleton_eq

\item \emph{The cutoff probability and total variation.}
The retained baseline, treatment, and flag pattern has the same distribution
under both priors. In particular, for \(x<J\),
\[
V_x\sim\operatorname{Binomial}(n,b_0qH).
\]
Define
\[
k_{\mathrm{tail}}=K+1,\qquad
\nu_{\mathrm{act}}=nb_0qH=\frac{\eta K^2}{\ell}.
\]
For an integer \(v\ge0\), let its falling factorial be
\[
(v)_{k_{\mathrm{tail}}}
=
\begin{cases}
\displaystyle\prod_{t=0}^{k_{\mathrm{tail}}-1}(v-t),
&v\ge k_{\mathrm{tail}},\\
0,&v<k_{\mathrm{tail}}.
\end{cases}
\]
Expanding over ordered distinct record indices gives
\[
E[(V_x)_{k_{\mathrm{tail}}}]
=(n)_{k_{\mathrm{tail}}}(b_0qH)^{k_{\mathrm{tail}}}
\le\nu_{\mathrm{act}}^{k_{\mathrm{tail}}}.
\]
When \(k_{\mathrm{tail}}>n\), both falling-factorial moments are zero.
On \(V_x>K\), the falling factorial is at least \(k_{\mathrm{tail}}!\).
Therefore
\[
\widetilde Q_{\mathrm{act},\iota}(V_x>K)
\le
\frac{\nu_{\mathrm{act}}^{k_{\mathrm{tail}}}}
{k_{\mathrm{tail}}!}.
\]

The bounds \(K\le2\ell\) and \(K\le k_{\mathrm{tail}}\) imply
\[
K^2\le2\ell k_{\mathrm{tail}},
\qquad
\nu_{\mathrm{act}}\le2\eta k_{\mathrm{tail}}
\le e^{-31}k_{\mathrm{tail}},
\qquad
k_{\mathrm{tail}}\ge\ell,
\]
where \(2\le e\) gives the second inequality involving the exponential.
The exponential series gives
\(k_{\mathrm{tail}}^{k_{\mathrm{tail}}}/k_{\mathrm{tail}}!
\le e^{k_{\mathrm{tail}}}\), whence
\[
\frac{\nu_{\mathrm{act}}^{k_{\mathrm{tail}}}}
{k_{\mathrm{tail}}!}
\le
e^{-31k_{\mathrm{tail}}}
\frac{k_{\mathrm{tail}}^{k_{\mathrm{tail}}}}
{k_{\mathrm{tail}}!}
\le e^{-30k_{\mathrm{tail}}}
\le e^{-30\ell}.
\]
A union bound over the rare cells yields
\[
\widetilde Q_{\mathrm{act},\iota}(\mathcal G^c)
\le Je^{-30\ell}.
\]
The mass cap and elementary exponential bounds give
\[
J\le\frac{N\ell}{2\eta},\qquad
N\le e^\ell,\qquad \ell\le e^\ell,\qquad
N\ell\le e^{2\ell}.
\]
Consequently,
\[
\widetilde Q_{\mathrm{act},\iota}(\mathcal G^c)
\le
\frac{N\ell}{2\eta}e^{-30\ell}
\le
\frac{e^{2\ell}}{\eta}e^{-30\ell}
=e^{32-28\ell}.
\]
The threshold \(N\ge e^{(32+\log8)/28}\) implies
\[
\ell\ge\frac{32+\log8}{28},
\qquad
\widetilde Q_{\mathrm{act},\iota}(\mathcal G^c)\le\frac18.
\]

Equality of the atoms on \(\mathcal G\) implies equality of the probabilities
of \(\mathcal G\), and hence of its complement. For every augmented-sample
event \(\mathcal B\),
\[
\begin{aligned}
&\left|
\widetilde Q_{\mathrm{act},0}(\mathcal B)
-\widetilde Q_{\mathrm{act},1}(\mathcal B)
\right|\\
&\quad=
\left|
\widetilde Q_{\mathrm{act},0}(\mathcal B\cap\mathcal G^c)
-\widetilde Q_{\mathrm{act},1}(\mathcal B\cap\mathcal G^c)
\right|
\le\widetilde Q_{\mathrm{act},0}(\mathcal G^c).
\end{aligned}
\]
Applying this inequality to preimages under the forgetting map gives
\[
\operatorname{TV}(Q_{\mathrm{act},0},Q_{\mathrm{act},1})
\le
\operatorname{TV}(\widetilde Q_{\mathrm{act},0},
\widetilde Q_{\mathrm{act},1})
\le\frac18.
\]
% lean: CausalSmith.Stat.MarRareqLogfrontier.activatedAugmentedMixture_cutoff_budget

\item \emph{From center separation to squared-error risk.}
For \(\iota\in\{0,1\}\), define the prior center and its separation by
\[
\theta_{\mathrm{act},\iota}
=E_{\Pi_\iota^{\otimes d}}[\tau_{\mathrm{act}}(\mathbf Z)]
=Jb_0\int z^{-1}\,d\Pi_\iota(z),
\qquad
\Delta_{\mathrm{act}}
=|\theta_{\mathrm{act},1}-\theta_{\mathrm{act},0}|.
\]
The reciprocal-mean separation implies
\[
\Delta_{\mathrm{act}}\ge\frac{Jb_0}{12}.
\]
Since \(Z_x^{-1}\in[0,1]\), its variance satisfies
\[
\operatorname{Var}_{\Pi_\iota}(Z_x^{-1})
=
E_{\Pi_\iota}\!\left[(Z_x^{-1}-\tfrac12)^2\right]
-\left(E_{\Pi_\iota}[Z_x^{-1}]-\tfrac12\right)^2
\le\frac14.
\]
Independence of the latent coordinates therefore yields
\[
E_{\Pi_\iota^{\otimes d}}\!\left[
(\tau_{\mathrm{act}}(\mathbf Z)-\theta_{\mathrm{act},\iota})^2
\right]
=
b_0^2\sum_{x=0}^{J-1}
\operatorname{Var}_{\Pi_\iota}(Z_x^{-1})
\le\frac{Jb_0^2}{4}.
\]

Fix \(\mathsf T\in\mathcal T_n\), and define its maximal model risk and its
centered mixture risks by
\[
\mathcal W(\mathsf T)
=
\sup_{P\in\mathcal M_{n,d,q}}
E_P[(\mathsf T-\tau(P))^2],
\qquad
C_{\mathrm{act},\iota}(\mathsf T)
=
\int
(b_{\mathsf T}(o)-\theta_{\mathrm{act},\iota})^2\,
dQ_{\mathrm{act},\iota}(o).
\]
The bounded action and target make these risks finite. Class membership of
every prior support law gives
\[
\int
E_{P_{\mathrm{act},\mathbf z}}\!\left[
(\mathsf T-\tau_{\mathrm{act}}(\mathbf z))^2
\right]
d\Pi_\iota^{\otimes d}(\mathbf z)
\le\mathcal W(\mathsf T).
\]
Using the barycenter comparison from
\cref{obj:lem:randomized-kernel-barycenter} and
\[
(b_{\mathsf T}(o)-\theta_{\mathrm{act},\iota})^2
\le
2(b_{\mathsf T}(o)-\tau_{\mathrm{act}}(\mathbf z))^2
+
2(\tau_{\mathrm{act}}(\mathbf z)-\theta_{\mathrm{act},\iota})^2
\]
inside the mixture integral gives
\[
C_{\mathrm{act},\iota}(\mathsf T)
\le2\mathcal W(\mathsf T)+\frac{Jb_0^2}{2}.
\]

Define the nearer-center test
\[
\psi_{\mathrm{act}}(o)
=
\mathbf1\left\{
|b_{\mathsf T}(o)-\theta_{\mathrm{act},1}|
\le
|b_{\mathsf T}(o)-\theta_{\mathrm{act},0}|
\right\}.
\]
The triangle inequality shows that a testing error under center \(\iota\)
entails
\[
|b_{\mathsf T}(o)-\theta_{\mathrm{act},\iota}|
\ge\frac{\Delta_{\mathrm{act}}}{2}.
\]
By \cref{obj:lem:randomized-kernel-tv-comparison},
\[
Q_{\mathrm{act},0}(\psi_{\mathrm{act}}=1)
+
Q_{\mathrm{act},1}(\psi_{\mathrm{act}}=0)
\ge
1-\operatorname{TV}(Q_{\mathrm{act},0},Q_{\mathrm{act},1}).
\]
Bounding each error probability by its centered squared risk consequently
gives
\[
\frac{\Delta_{\mathrm{act}}^2}{8}
\left[
1-\operatorname{TV}(Q_{\mathrm{act},0},Q_{\mathrm{act},1})
\right]
\le
\max_{\iota\in\{0,1\}}C_{\mathrm{act},\iota}(\mathsf T).
\]
Together with the total variation and centered-risk bounds, this implies
\[
\frac{7\Delta_{\mathrm{act}}^2}{128}
-\frac{Jb_0^2}{4}
\le\mathcal W(\mathsf T).
\]
Now \(J\ge2048\) and \(\Delta_{\mathrm{act}}\ge Jb_0/12\), so
\[
\mathcal W(\mathsf T)
\ge
(Jb_0)^2
\left(\frac7{18432}-\frac1{4J}\right)
\ge\frac{(Jb_0)^2}{4096}.
\]
Taking the infimum over \(\mathsf T\) gives the calibrated fuzzy bound
\[
\frac{(Jb_0)^2}{4096}\le\mathfrak R_{n,d,q}.
\]
% lean: CausalSmith.Stat.MarRareqLogfrontier.activatedAugmentedMixture_minimaxRisk_calibrated

\item \emph{Converting rare-cell mass to the dimensional contribution.}
Still suppose \(N\ge N_{\mathrm{act}}\) and \(J\ge2048\). We have
\[
b_0=\frac{\eta}{N\ell}\le\frac14,
\qquad
0<\eta=e^{-32}\le\frac12,
\qquad d\ge2.
\]
The first inequality uses \(N\ge4\eta\) and \(\ell\ge1\); the second follows
from \(e^{32}\ge1+32\).

If \(d-1\le\lfloor(2b_0)^{-1}\rfloor\), then
\[
J=d-1\ge\frac d2,\qquad Jb_0\ge\frac{db_0}{2}.
\]
Otherwise the floor determines \(J\), and
\[
Jb_0
=\left\lfloor\frac1{2b_0}\right\rfloor b_0
\ge\frac12-b_0\ge\frac14.
\]
Thus both branches give
\[
\min\left\{\frac{db_0}{2},\frac14\right\}\le Jb_0.
\]
Moreover,
\[
\frac{\eta}{24}\min\{1,u_{\mathrm{act}}\}
\le\frac{\eta u_{\mathrm{act}}}{24}
=\frac{db_0}{24},
\qquad
\frac{\eta}{24}\min\{1,u_{\mathrm{act}}\}
\le\frac{\eta}{24}\le\frac1{48}.
\]
It follows that
\[
\frac{\eta}{24}\min\{1,u_{\mathrm{act}}\}
\le
\frac1{12}\min\left\{\frac{db_0}{2},\frac14\right\}
\le\frac{Jb_0}{12},
\]
and hence
\[
\frac{\eta}{2}\min\{1,u_{\mathrm{act}}\}\le Jb_0.
\]
For \(u_{\mathrm{act}}\ge0\), considering \(u_{\mathrm{act}}\le1\) and
\(u_{\mathrm{act}}\ge1\) gives
\[
\bigl(\min\{1,u_{\mathrm{act}}\}\bigr)^2
=\min\{1,u_{\mathrm{act}}^2\}.
\]
Squaring the mass bound therefore yields
\[
\frac{\eta^2}{4}\min\{1,u_{\mathrm{act}}^2\}
\le(Jb_0)^2.
\]
The choice of \(\underline c_{\mathrm{act}}\) and the calibrated fuzzy bound
now imply
\[
\underline c_{\mathrm{act}}\min\{1,b_{\mathrm{dim}}\}
\le
\frac{\eta^2}{16384}\min\{1,u_{\mathrm{act}}^2\}
\le\frac{(Jb_0)^2}{4096}
\le\mathfrak R_{n,d,q}.
\]
% lean: CausalSmith.Stat.MarRareqLogfrontier.interval_prior_segment_width_lower

\item \emph{Large effective size with fewer rare cells.}
Suppose instead that \(N\ge N_{\mathrm{act}}\) and \(J<2048\).
The capacity bound established above gives
\[
\left\lfloor\frac1{2b_0}\right\rfloor\ge2048.
\]
The definition of \(J\) therefore forces
\[
d-1<2048,\qquad d\le2048.
\]
This includes \(J=0\), for which the same argument gives \(d=1\).

Since \(N\ge1\) and \(\ell\ge1\),
\[
\sqrt N\le N\le N\ell.
\]
Consequently,
\[
\min\{1,u_{\mathrm{act}}\}
\le\frac{d}{N\ell}
\le\frac{2048}{\sqrt N},
\qquad
\min\{1,u_{\mathrm{act}}^2\}
\le\frac{2048^2}{N}.
\]
Also,
\[
\min\{1,N^{-1}\}=N^{-1},
\qquad
\underline c_{\mathrm{act}}\,2048^2\le\frac1{256}.
\]
The numerical one-cell estimate thus gives
\[
\underline c_{\mathrm{act}}\min\{1,b_{\mathrm{dim}}\}
\le
\frac{\underline c_{\mathrm{act}}\,2048^2}{N}
\le\frac1{256N}
\le\mathfrak R_{n,d,q}.
\]
% lean: CausalSmith.Stat.MarRareqLogfrontier.activated_fuzzy_lower

\item \emph{Small effective size.}
In the remaining effective-size branch, \(N<N_{\mathrm{act}}\).
Because \(N>0\) and \(N_{\mathrm{act}}\ge1\),
\[
N_{\mathrm{act}}^{-1}\le1,\qquad
N_{\mathrm{act}}^{-1}\le N^{-1},
\qquad
N_{\mathrm{act}}^{-1}\le\min\{1,N^{-1}\}.
\]
The definition of \(\underline c_{\mathrm{act}}\) and the numerical
one-cell estimate give
\[
\underline c_{\mathrm{act}}\min\{1,b_{\mathrm{dim}}\}
\le\underline c_{\mathrm{act}}
\le\frac1{256N_{\mathrm{act}}}
\le\frac1{256}\min\{1,N^{-1}\}
\le\mathfrak R_{n,d,q}.
\]
The three effective-size and rare-count branches establish the complementary
bound whenever \(\sqrt N\le d\).
% lean: CausalSmith.Stat.MarRareqLogfrontier.activated_fuzzy_lower

\item \emph{Combining the bounds when \(d\ge\sqrt N\).}
First suppose \(\sqrt N\le d\). Both component bounds apply, and their
nonnegative capped factors give
\[
\underline c_*\min\{1,a_{\mathrm{eff}}\}
\le
\underline c_{\mathrm{one}}\min\{1,a_{\mathrm{eff}}\}
\le\mathfrak R_{n,d,q},
\]
\[
\underline c_*\min\{1,b_{\mathrm{dim}}\}
\le
\underline c_{\mathrm{act}}\min\{1,b_{\mathrm{dim}}\}
\le\mathfrak R_{n,d,q}.
\]

For the nonnegative quantities \(a_{\mathrm{eff}},b_{\mathrm{dim}}\),
\[
\min\{1,a_{\mathrm{eff}}+b_{\mathrm{dim}}\}
\le
\min\{1,a_{\mathrm{eff}}\}+\min\{1,b_{\mathrm{dim}}\}.
\]
Indeed, if both quantities are at most one, the right-hand side is their sum;
if either is at least one, the right-hand side is at least one.
Each capped quantity is at most their maximum, so
\[
\min\{1,a_{\mathrm{eff}}+b_{\mathrm{dim}}\}
\le
2\max\{\min\{1,a_{\mathrm{eff}}\},\min\{1,b_{\mathrm{dim}}\}\},
\]
while the component risk bounds imply
\[
\underline c_*
\max\{\min\{1,a_{\mathrm{eff}}\},\min\{1,b_{\mathrm{dim}}\}\}
\le\mathfrak R_{n,d,q}.
\]
Since \(\mathfrak R_{n,d,q}\ge0\), we conclude that
\[
\underline c\,r_{n,d,q}
\le
\frac{\underline c_*}{2}
\max\{\min\{1,a_{\mathrm{eff}}\},\min\{1,b_{\mathrm{dim}}\}\}
\le\frac12\mathfrak R_{n,d,q}
\le\mathfrak R_{n,d,q}.
\]
% lean: CausalSmith.Stat.MarRareqLogfrontier.full_experiment_lower

\item \emph{Combining the bounds when \(d<\sqrt N\).}
Suppose \(\sqrt N>d\). Since \(d\ge0\),
\[
d^2\le(\sqrt N)^2=N.
\]
Also \(\ell\ge1\), so
\[
\ell^2\ge1,\qquad
(N\ell)^2>0,\qquad
N^{-1}(N\ell)^2=N\ell^2.
\]
Therefore
\[
b_{\mathrm{dim}}
=\frac{d^2}{(N\ell)^2}
\le\frac{N\ell^2}{(N\ell)^2}
=N^{-1}
=a_{\mathrm{eff}},
\]
and monotonicity of the minimum gives
\[
\min\{1,b_{\mathrm{dim}}\}
\le\min\{1,a_{\mathrm{eff}}\}.
\]
Applying the capped-sum inequality established above,
\[
r_{n,d,q}
\le
\min\{1,a_{\mathrm{eff}}\}+\min\{1,b_{\mathrm{dim}}\}
\le2\min\{1,a_{\mathrm{eff}}\}.
\]
Finally, the one-cell contribution and
\(\underline c_*\le\underline c_{\mathrm{one}}\) yield
\[
\underline c\,r_{n,d,q}
\le
\frac{\underline c_*}{2}\min\{1,a_{\mathrm{eff}}\}
\le
\underline c_{\mathrm{one}}\min\{1,a_{\mathrm{eff}}\}
\le\mathfrak R_{n,d,q}.
\]
The two dimension branches cover all admissible parameters.
% lean: CausalSmith.Stat.MarRareqLogfrontier.full_experiment_lower
\end{enumerate}
\end{proof}
